%============================================================
%  Bd. 26 - Bäume
%============================================================

\documentclass[main.tex]{subfiles}

\ifSubfilesClassLoaded{
  \BandSetupStandalone{B26}
}{
}

\title{Bd. 26 - Bäume}
\author{Martin Kunze}
\date{}
\setcounter{file}{26}

\hypersetup{
  pdftitle={Bd. 26 - Bäume},
  pdfauthor={Martin Kunze}
}

% Die allgemeinen Graphsymbole stehen in tex/impl/commands/graphs.tex.
% Die folgenden Bezeichnungen sind auf diesen Band beschränkt und werden mit
% \providecommand definiert, damit ein späteres Herausziehen in die zentrale
% Befehlsdatei ohne Kollision möglich bleibt.
\providecommand{\GraphPathsBetween}[4]{%
  \ensuremath{\mathcal P_{#1,#2}\!\left(#3,#4\right)}%
}
\providecommand{\ForestStruct}[1]{%
  \ensuremath{\mathsf{Wald}\!\left(#1\right)}%
}
\providecommand{\TreeStruct}[1]{%
  \ensuremath{\mathsf{Baum}\!\left(#1\right)}%
}
\providecommand{\PathAcyclicGraphStruct}[1]{%
  \ensuremath{\mathsf{WegAzyk}\!\left(#1\right)}%
}
\providecommand{\RootedTreeStruct}[1]{%
  \ensuremath{\mathsf{WBaum}\!\left(#1\right)}%
}
\providecommand{\TreeParentRel}[5]{%
  \ensuremath{\operatorname{Elter}_{#1,#2,#3}\!\left(#4,#5\right)}%
}
\providecommand{\TreeChildRel}[5]{%
  \ensuremath{\operatorname{Kind}_{#1,#2,#3}\!\left(#4,#5\right)}%
}
\providecommand{\TreeAncestorRel}[5]{%
  \ensuremath{\operatorname{Vorfahr}_{#1,#2,#3}\!\left(#4,#5\right)}%
}
\providecommand{\TreeChildren}[4]{%
  \ensuremath{\operatorname{Kind}_{#1,#2,#3}\!\left[#4\right]}%
}
\providecommand{\TreeLeafPred}[4]{%
  \ensuremath{\operatorname{Blatt}_{#1,#2,#3}\!\left(#4\right)}%
}
\providecommand{\TreeInnerPred}[4]{%
  \ensuremath{\operatorname{Innen}_{#1,#2,#3}\!\left(#4\right)}%
}
\providecommand{\TreeInnerVertices}[3]{%
  \ensuremath{\operatorname{Innen}_{#1,#2,#3}}%
}
\providecommand{\BinaryTreeStruct}[1]{%
  \ensuremath{\mathsf{BinBaum}\!\left(#1\right)}%
}
\providecommand{\FullBinaryTreeStruct}[1]{%
  \ensuremath{\mathsf{VollBinBaum}\!\left(#1\right)}%
}
\providecommand{\OrderedFullBinaryTreeStruct}[1]{%
  \ensuremath{\mathsf{GVollBinBaum}\!\left(#1\right)}%
}
\providecommand{\FinitePathTransportProperty}[3]{%
  \ensuremath{\mathsf{Transp}_{#1,#2}\!\left(#3\right)}%
}
\providecommand{\FinitePathReversalProperty}[3]{%
  \ensuremath{\mathsf{Umkehr}_{#1,#2}\!\left(#3\right)}%
}
\providecommand{\FinitePathSuffixProperty}[5]{%
  \ensuremath{\mathsf{Schwanz}_{#1,#2,#3,#4}\!\left(#5\right)}%
}
\providecommand{\FinitePathSpliceProperty}[3]{%
  \ensuremath{\mathsf{Füge}_{#1,#2}\!\left(#3\right)}%
}
\providecommand{\FinitePathCycleCore}[3]{%
  \ensuremath{%
    #3\subseteq #1\land #3\neq\varnothing\land
    \forall x\in #3\,\exists y\in #3\,
    \exists m\in\mathbb N\,\exists n\in\mathbb N\,
    \exists p\in\powerset(\mathbb N\times #1)\,
    \exists q\in\powerset(\mathbb N\times #1)\,\Bigl(
      x\neq y\land
      \DirectedPath{#1}{#2}{p}{m}\land p(0)=x\land p(m)=y\land
      \DirectedPath{#1}{#2}{q}{n}\land q(0)=x\land q(n)=y\land
      (m,p)\neq(n,q)\land
      p[\NatSeg{m+1}]\cap q[\NatSeg{n+1}]=\{x,y\}\land
      p[\NatSeg{m+1}]\cup q[\NatSeg{n+1}]\subseteq #3
    \Bigr)%
  }%
}

\begin{document}

\title{Bd. 26 - Bäume}
\author{Martin Kunze}
\date{}
\hypersetup{pageanchor=false}
\maketitle
\hypersetup{pageanchor=true}
\setcounter{tocdepth}{3}
\tableofcontents
\setcounter{file}{26}
\renewcommand{\FormulaBandID}{26}

\chapter{Einleitung}

Ein Baum ist in diesem Band kein Bild, sondern ein schlichter ungerichteter
Graph mit einer ausgezeichneten Wegeigenschaft. Ein endlicher einfacher Weg
wird dabei durch eine natürliche Kantenlänge und eine injektive endliche Folge
seiner Knoten dargestellt. Der Graph selbst darf unendlich sein; nur jeder
einzelne Weg ist endlich.

Wälder werden dadurch charakterisiert, dass zwischen zwei Knoten höchstens
ein endlicher einfacher Weg liegt. Bei einem Baum ist der Träger nicht leer
und zwischen je zwei Knoten liegt genau ein solcher Weg. Diese Definition
liefert Azyklizität und Zusammenhang getrennt und ohne ein zusätzliches
Endlichkeitsaxiom. Nach der Wahl einer Wurzel entstehen Eltern, Kinder und
Vorfahren aus den eindeutigen Wurzelwegen. Den Abschluss bilden binäre, volle
binäre und geordnete volle Binärbäume. Im folgenden Band werden die dort
konstruierten Klammerungsbäume als besondere geordnete volle Binärbäume
erkannt.

\chapter{Endliche einfache Wege}

\section{Wege zwischen festgelegten Endpunkten}

Mit \(\DirectedPath{V}{E}{p}{n}\) bezeichnet Band~22 einen endlichen
einfachen Weg. Dabei ist \(p\colon\NatSeg{n+1}\to V\) injektiv, und
aufeinanderfolgende Knoten sind durch Kanten aus \(E\) verbunden. Für
ungerichtete Graphen darf dieselbe Aussage verwendet werden, weil ihre
Kantenrelation symmetrisch ist.

\FormulaDefDeltaK[Menge der endlichen einfachen Wege]{%
  \begin{aligned}
  \GraphPathsBetween{V}{E}{u}{v}\coloneqq
  \bigl\{(n,p)\in
    \mathbb N\times\powerset(\mathbb N\times V)\bigm|{}&
      \DirectedPath{V}{E}{p}{n}\land{}\\[-2pt]
      &p(0)=u\land p(n)=v\bigr\}
  \end{aligned}%
}{FiniteSimplePathFamilyDef}{%
  \DeltaRow{Mengen}{V\dsep E\dsep u\dsep v\dsep n\dsep p}
  \DeltaPrem{Graph}{\GraphStruct{V,E}}
  \DeltaRow{Neue Schreibweise}{\GraphPathsBetween{V}{E}{u}{v}}
}

Ein Element von \(\GraphPathsBetween{V}{E}{u}{v}\) enthält die Länge und
die Knotenfolge. Dadurch werden Wege verschiedener Länge nicht versehentlich
miteinander identifiziert.

\FormulaThmDeltaK[Zugriff auf einen endlichen einfachen Weg]{%
  \begin{aligned}[t]
  z\in\GraphPathsBetween{V}{E}{u}{v}\eqvdash
  \exists n\in\mathbb N\,\exists p\in\powerset(\mathbb N\times V)\,\bigl(&
    z=(n,p)\land\DirectedPath{V}{E}{p}{n}\land{}\\[-2pt]
    &p(0)=u\land p(n)=v\bigr)
  \end{aligned}%
}{FiniteSimplePathFamilyMembership}{%
  \DeltaRow{Mengen}{V\dsep E\dsep u\dsep v\dsep z\dsep n\dsep p}
  \DeltaPrem{Graph}{\GraphStruct{V,E}}
}
\begin{tabproofwide}
  \proofstepwidestarbreak[1]{z\in\GraphPathsBetween{V}{E}{u}{v}}{\rA}
  \proofstepwidestarbreak[1]{%
    \exists n\in\mathbb N\,\exists p\in\powerset(\mathbb N\times V)\,
      \bigl(z=(n,p)\land\DirectedPath{V}{E}{p}{n}\land
      p(0)=u\land p(n)=v\bigr)}{%
    \FormulaRefAuto{FiniteSimplePathFamilyDef}[def]{1}}
  \proofstepwidestarbreak[3]{%
    \exists n\in\mathbb N\,\exists p\in\powerset(\mathbb N\times V)\,
      \bigl(z=(n,p)\land\DirectedPath{V}{E}{p}{n}\land
      p(0)=u\land p(n)=v\bigr)}{\rA}
  \proofstepwidestarbreak[3]{z\in\GraphPathsBetween{V}{E}{u}{v}}{%
    \FormulaRefAuto{FiniteSimplePathFamilyDef}[def]{3}}
  \proofstepwidestarbreak[]{%
    z\in\GraphPathsBetween{V}{E}{u}{v}\leftrightarrow
    \exists n\in\mathbb N\,\exists p\in\powerset(\mathbb N\times V)\,
      \bigl(z=(n,p)\land\DirectedPath{V}{E}{p}{n}\land
      p(0)=u\land p(n)=v\bigr)}{\rLRI{\rRI{1,2},\rRI{3,4}}}
\end{tabproofwide}

\FormulaThmDeltaK[Knotenfolge eines gespeicherten Weges]{%
  \GraphStruct{V,E}\dsep(n,p)\in\GraphPathsBetween{V}{E}{u}{v}
  \vdash p\colon\NatSeg{n+1}\to V%
}{FiniteSimplePathSequence}{%
  \DeltaRow{Mengen}{V\dsep E\dsep u\dsep v\dsep n\dsep p}
}
\begin{tabproofwide}
  \proofstepwidestar[1]{\GraphStruct{V,E}}{\rA}
  \proofstepwidestar[2]{(n,p)\in\GraphPathsBetween{V}{E}{u}{v}}{\rA}
  \proofstepwidestar[1,2]{\DirectedPath{V}{E}{p}{n}}{%
    \FormulaRefAuto{FiniteSimplePathFamilyDef}[def]{1,2}}
  \proofstepwidestar[1,2]{p\colon\NatSeg{n+1}\to V}{%
    \FormulaRefAuto{DirectedPathFunction}[thm]{1,3}}
\end{tabproofwide}

\FormulaThmDeltaK[Endpunkte eines Weges gehören zum Träger]{%
  \GraphStruct{V,E}\dsep
  (n,p)\in\GraphPathsBetween{V}{E}{u}{v}
  \vdash u\in V\land v\in V%
}{FiniteSimplePathEndpointTyping}{%
  \DeltaRow{Mengen}{V\dsep E\dsep u\dsep v\dsep n\dsep p}
}
\begin{tabproofwide}
  \proofstepwidestarbreak[1]{\GraphStruct{V,E}}{\rA}
  \proofstepwidestarbreak[2]{(n,p)\in\GraphPathsBetween{V}{E}{u}{v}}{\rA}
  \proofstepwidestarbreak[2]{\DirectedPath{V}{E}{p}{n}\land
    p(0)=u\land p(n)=v}{%
    \FormulaRefAuto{FiniteSimplePathFamilyMembership}[thm]{2}}
  \proofstepwidestar[1,2]{\DirectedWalk V E p n}{%
      \FormulaRefAuto{DirectedPathWalkAxiom}[ax]{1,\rAEa{3}}}
  \proofstepwidestarbreak[1,2]{p(0)\in V\land p(n)\in V}{%
    \FormulaRefAuto{DirectedWalkEndpointTyping}[thm]{%
      1,4}}
  \proofstepwidestarbreak[1,2]{u\in V\land v\in V}{%
    \rIE{\rAEb{3},\rIE{\rAEn{3},5}}}
\end{tabproofwide}

\section{Elementare Operationen mit endlichen Wegen}

Die folgenden Hilfssätze formulieren die drei endlichen Operationen, die
später benötigt werden: Abschneiden eines Anfangsstücks, Fortsetzen um eine
Kante unter sofortiger Entfernung einer entstehenden Wiederholung und
Umkehrung in einem ungerichteten Graphen. Die Beweise benutzen durchgehend
die Plus-eins-Indizierung aus Band~22.

\FormulaThmDeltaK[Indexangaben für einen Weganfang]{%
  n\in\mathbb N\dsep k\in\NatSeg{n+1}\vdash
  k\in\mathbb N\land k\leq n\land
  \NatSeg{k+1}\subseteq\NatSeg{n+1}%
}{FinitePathPrefixIndexData}{%
  \DeltaRow{Mengen}{n\dsep k}
}
\begin{tabproofsplitwide}
  \proofpartwide{Typisierung und Schranke des Schnittindex}
  \proofstepwidestarbreak[1]{n\in\mathbb N}{\rA}
  \proofstepwidestarbreak[2]{k\in\NatSeg{n+1}}{\rA}
  \proofstepwidestarbreak[]{1\in\mathbb N}{\FormulaRefAuto{PeanoOneInN}[thm]}
  \proofstepwidestarbreak[1]{n+1\in\mathbb N}{%
    \FormulaRefAuto{PeanoAddClosure}[thm]{1,3}}
  \proofstepwidestarbreak[1]{\NatSeg{n+1}\subseteq\mathbb N}{%
    \FormulaRefAuto{PeanoNatSegSubsetN}[thm]{4}}
  \proofstepwidestarbreak[1,2]{k\in\mathbb N}{%
    \FormulaRefAuto{A\subseteq B\dsep x\in A\vdash x\in B}[thm]{5,2}}
  \proofstepwidestarbreak[1,2]{k<n+1}{%
    \FormulaRefAuto{PeanoNatSegStrictMembership}[thm]{6,4,2}}
  \proofstepwidestarbreak[1,2]{k+1\leq n+1}{%
    \FormulaRefAuto{PeanoLtAddOneLeq}[thm]{6,4,7}}
  \proofstepwidestar[1,2]{k\leq n\leftrightarrow k\in\NatSeg{n+1}}{%
      \FormulaRefAuto{PeanoNatSegLeqAddOneMembership}[thm]{6,1}}
  \proofstepwidestarbreak[1,2]{k\leq n}{%
    \FormulaRefAuto{P\leftrightarrow Q\dsep Q\vdash P}[thm]{%
      9,2}}
  \closeproofpartwide

  \proofpartwide{Inklusion der Anfangsabschnitte}
  \setcounter{proofstepnr}{10}
  \proofstepwidestarbreak[1,2]{k+1\in\mathbb N}{%
    \FormulaRefAuto{PeanoAddClosure}[thm]{6,3}}
  \proofstepwidestarbreak[1,2]{\NatSeg{k+1}\subseteq\NatSeg{n+1}}{%
    \FormulaRefAuto{PeanoNatSegLeqSubsetForward}[thm]{11,4,8}}
  \proofstepwidestarbreak[1,2]{%
    k\in\mathbb N\land k\leq n\land
    \NatSeg{k+1}\subseteq\NatSeg{n+1}}{\rAI{6,10,12}}
  \closeproofpartwide
\end{tabproofsplitwide}

\FormulaThmDeltaK[Verschiebung eines Wegindex um eins]{%
  n\in\mathbb N\dsep i\in\NatSeg n\vdash
  i+1\in\NatSeg{n+1}%
}{FinitePathIndexShiftOne}{%
  \DeltaRow{Mengen}{n\dsep i}
}
\begin{tabproofwide}
  \proofstepwidestarbreak[1]{n\in\mathbb N}{\rA}
  \proofstepwidestarbreak[2]{i\in\NatSeg n}{\rA}
  \proofstepwidestarbreak[1]{\NatSeg n\subseteq\mathbb N}{%
    \FormulaRefAuto{PeanoNatSegSubsetN}[thm]{1}}
  \proofstepwidestarbreak[1,2]{i\in\mathbb N}{%
    \FormulaRefAuto{A\subseteq B\dsep x\in A\vdash x\in B}[thm]{3,2}}
  \proofstepwidestarbreak[1,2]{i<n}{%
    \FormulaRefAuto{PeanoNatSegStrictMembership}[thm]{4,1,2}}
  \proofstepwidestarbreak[]{1\in\mathbb N}{\FormulaRefAuto{PeanoOneInN}[thm]}
  \proofstepwidestarbreak[1,2]{1+i<1+n}{%
    \FormulaRefAuto{PeanoAddLeftStrictMonotone}[thm]{6,4,1,5}}
  \proofstepwidestar[1,2]{i + 1 = 1 + i}{%
      \FormulaRefAuto{PeanoAddCommutative}[thm]{4,6}}
  \proofstepwidestar[1,2]{n + 1 = 1 + n}{%
      \FormulaRefAuto{PeanoAddCommutative}[thm]{1,6}}
  \proofstepwidestarbreak[1,2]{i+1<n+1}{%
    \rIE{8,
      \rIE{9,7}}}
  \proofstepwidestarbreak[1,2]{i+1\in\mathbb N}{%
    \FormulaRefAuto{PeanoAddClosure}[thm]{4,6}}
  \proofstepwidestarbreak[1]{n+1\in\mathbb N}{%
    \FormulaRefAuto{PeanoAddClosure}[thm]{1,6}}
  \proofstepwidestarbreak[1,2]{i+1\in\NatSeg{n+1}}{%
    \FormulaRefAuto{PeanoNatSegStrictMembership}[thm]{11,12,10}}
\end{tabproofwide}

\FormulaThmDeltaK[Zugriff auf die Typdaten eines endlichen Pfads]{%
  \GraphStruct{V,E}\dsep\DirectedPath{V}{E}{p}{n}\vdash
  n\in\mathbb N\land p\colon\NatSeg{n+1}\to V\land
  p\colon\NatSeg{n+1}\inj V%
}{FinitePathTyping}{%
  \DeltaRow{Mengen}{V\dsep E\dsep p\dsep n}
}
\begin{tabproofwide}
  \proofstepwidestarbreak[1]{\GraphStruct{V,E}}{\rA}
  \proofstepwidestarbreak[2]{\DirectedPath{V}{E}{p}{n}}{\rA}
  \proofstepwidestarbreak[1,2]{\DirectedWalk{V}{E}{p}{n}}{%
    \FormulaRefAuto{DirectedPathWalkAxiom}[ax]{1,2}}
  \proofstepwidestarbreak[1,2]{n\in\mathbb N\land
    p\colon\NatSeg{n+1}\to V\land
    \forall i\in\NatSeg n\,((p(i),p(i+1))\in E)}{%
    \FormulaRefAuto{DirectedWalkCharacterization}[thm]{1,3}}
  \proofstepwidestar[1,2]{\DirectedWalk{V}{E}{p}{n}\land p\colon\NatSeg{n+1}\inj V}{%
      \FormulaRefAuto{DirectedPathCharacterization}[thm]{1,2}}
  \proofstepwidestarbreak[1,2]{p\colon\NatSeg{n+1}\inj V}{%
    \rAEb{5}}
  \proofstepwidestar[1,2]{{ p \colon \NatSeg { n + 1 } \to V }}{%
      \FormulaRefAuto{P\land Q\land R\vdash Q}[thm]{4}}
  \proofstepwidestarbreak[1,2]{n\in\mathbb N\land
    p\colon\NatSeg{n+1}\to V\land p\colon\NatSeg{n+1}\inj V}{%
    \rAI{\rAEa{4},
      7,6}}
\end{tabproofwide}

\FormulaThmDeltaK[Ein endlicher Walk ist als Relation typisiert]{%
  \GraphStruct{V,E}\dsep\DirectedWalk{V}{E}{p}{n}\vdash
  p\in\powerset(\mathbb N\times V)%
}{FiniteWalkPowerSetTyping}{%
  \DeltaRow{Mengen}{V\dsep E\dsep p\dsep n}
}
\begin{tabproofwide}
  \proofstepwidestarbreak[1]{\GraphStruct{V,E}}{\rA}
  \proofstepwidestarbreak[2]{\DirectedWalk{V}{E}{p}{n}}{\rA}
  \proofstepwidestarbreak[1,2]{n\in\mathbb N\land
    p\colon\NatSeg{n+1}\to V\land
    \forall i\in\NatSeg n\,((p(i),p(i+1))\in E)}{%
    \FormulaRefAuto{DirectedWalkCharacterization}[thm]{1,2}}
  \proofstepwidestar[1,2]{{ p \colon \NatSeg { n + 1 } \to V }}{%
      \FormulaRefAuto{P\land Q\land R\vdash Q}[thm]{3}}
  \proofstepwidestarbreak[1,2]{p\subseteq\NatSeg{n+1}\times V}{%
    \FormulaRefAuto{F\colon A\to B\vdash F\subseteq A\times B}[thm]{%
      4}}
  \proofstepwidestar[1,2]{n+1\in\mathbb N}{%
      \FormulaRefAuto{PeanoAddOneClosure}[thm]{\rAEa{3}}}
  \proofstepwidestar[1,2]{\NatSeg{n+1}\subseteq\mathbb N}{\FormulaRefAuto{PeanoNatSegSubsetN}[thm]{6}}
  \proofstepwidestarbreak[1,2]{\NatSeg{n+1}\times V
      \subseteq\mathbb N\times V}{%
    \FormulaRefAuto{CartesianProductMonotoneLeft}[thm]{%
      7}}
  \proofstepwidestarbreak[1,2]{p\subseteq\mathbb N\times V}{%
    \FormulaRefAuto{A\subseteq B\dsep B\subseteq C\vdash A\subseteq C}[thm]{5,8}}
  \proofstepwidestar[1,2]{p \in \powerset ( { \mathbb N \times V } ) \leftrightarrow p \subseteq { \mathbb N \times V }}{%
      \FormulaRefAuto{x \in \powerset(A)\;\eqvdash\; x \subseteq A}[thm]}
  \proofstepwidestarbreak[1,2]{p\in\powerset(\mathbb N\times V)}{%
    \FormulaRefAuto{P\leftrightarrow Q\dsep Q\vdash P}[thm]{%
      10,9}}
\end{tabproofwide}

\FormulaThmDeltaK[Ein endlicher Pfad ist als Relation typisiert]{%
  \GraphStruct{V,E}\dsep\DirectedPath{V}{E}{p}{n}\vdash
  p\in\powerset(\mathbb N\times V)%
}{FinitePathPowerSetTyping}{%
  \DeltaRow{Mengen}{V\dsep E\dsep p\dsep n}
}
\begin{tabproofwide}
  \proofstepwidestarbreak[1]{\GraphStruct{V,E}}{\rA}
  \proofstepwidestarbreak[2]{\DirectedPath{V}{E}{p}{n}}{\rA}
  \proofstepwidestar[1,2]{n \in \mathbb N \land p \colon \NatSeg { n + 1 } \to V \land p \colon \NatSeg { n + 1 } \inj V}{%
      \FormulaRefAuto{FinitePathTyping}[thm]{1,2}}
  \proofstepwidestarbreak[1,2]{p\colon\NatSeg{n+1}\to V}{%
    \FormulaRefAuto{P\land Q\land R\vdash Q}[thm]{%
      3}}
  \proofstepwidestarbreak[1,2]{p\subseteq\NatSeg{n+1}\times V}{%
    \FormulaRefAuto{F\colon A\to B\vdash F\subseteq A\times B}[thm]{4}}
  \proofstepwidestar[1,2]{n \in \mathbb N \land p \colon \NatSeg { n + 1 } \to V \land p \colon \NatSeg { n + 1 } \inj V}{%
      \FormulaRefAuto{FinitePathTyping}[thm]{1,2}}
  \proofstepwidestar[1,2]{n+1\in\mathbb N}{%
      \FormulaRefAuto{PeanoAddOneClosure}[thm]{%
        \rAEa{6}}}
  \proofstepwidestar[1,2]{\NatSeg{n+1}\subseteq\mathbb N}{\FormulaRefAuto{PeanoNatSegSubsetN}[thm]{7}}
  \proofstepwidestarbreak[1,2]{\NatSeg{n+1}\times V
      \subseteq\mathbb N\times V}{%
    \FormulaRefAuto{CartesianProductMonotoneLeft}[thm]{%
      8}}
  \proofstepwidestarbreak[1,2]{p\subseteq\mathbb N\times V}{%
    \FormulaRefAuto{A\subseteq B\dsep B\subseteq C\vdash A\subseteq C}[thm]{5,9}}
  \proofstepwidestar[1,2]{p \in \powerset ( { \mathbb N \times V } ) \leftrightarrow p \subseteq { \mathbb N \times V }}{%
      \FormulaRefAuto{x \in \powerset(A)\;\eqvdash\; x \subseteq A}[thm]}
  \proofstepwidestarbreak[1,2]{p\in\powerset(\mathbb N\times V)}{%
    \FormulaRefAuto{P\leftrightarrow Q\dsep Q\vdash P}[thm]{%
      11,10}}
\end{tabproofwide}

\FormulaThmDeltaK[Anfangsabschnitt eines endlichen Walks]{%
  \begin{aligned}[t]
  &\GraphStruct{V,E}\dsep\DirectedWalk{V}{E}{p}{n}\dsep
    k\in\NatSeg{n+1}\vdash{}\\[-2pt]
  &\quad\DirectedWalk{V}{E}{p\restriction_{\NatSeg{k+1}}}{k}
    \land p\restriction_{\NatSeg{k+1}}(0)=p(0)
    \land p\restriction_{\NatSeg{k+1}}(k)=p(k)
  \end{aligned}%
}{FiniteWalkInitialSegment}{%
  \DeltaRow{Mengen}{V\dsep E\dsep p\dsep q\dsep n\dsep k\dsep i}
}
\begin{tabproofsplitwide}
  \proofpartwide{Typisierung der eingeschränkten Folge}
  \proofstepwidestarbreak[1]{\GraphStruct{V,E}}{\rA}
  \proofstepwidestarbreak[2]{\DirectedWalk{V}{E}{p}{n}}{\rA}
  \proofstepwidestarbreak[3]{k\in\NatSeg{n+1}}{\rA}
  \proofstepwidestarbreak[1,2]{n\in\mathbb N\land
    p\colon\NatSeg{n+1}\to V\land
    \forall i\in\NatSeg n\,((p(i),p(i+1))\in E)}{%
    \FormulaRefAuto{DirectedWalkCharacterization}[thm]{1,2}}
  \proofstepwidestarbreak[1,2,3]{k\in\mathbb N\land k\leq n\land
    \NatSeg{k+1}\subseteq\NatSeg{n+1}}{%
    \FormulaRefAuto{FinitePathPrefixIndexData}[thm]{\rAEa{4},3}}
  \proofstepwidestar[1,2,3]{{ p \colon \NatSeg { n + 1 } \to V }}{%
      \FormulaRefAuto{P\land Q\land R\vdash Q}[thm]{4}}
  \proofstepwidestarbreak[1,2,3]{%
    p\restriction_{\NatSeg{k+1}}\colon\NatSeg{k+1}\to V}{%
    \FormulaRefAuto{C\subseteq A \vdash F\restriction_C\colon C\to B}[thm]{%
      \rAEn{5},6}}
  \closeproofpartwide

  \proofpartwide{Kantenbedingung des Anfangsabschnitts}
  \setcounter{proofstepnr}{7}
  \proofstepwidestarbreak[8]{i\in\NatSeg k}{\rA}
  \proofstepwidestarbreak[1,2,3]{\NatSeg k\subseteq\NatSeg n}{%
    \FormulaRefAuto{PeanoNatSegLeqSubsetForward}[thm]{%
      \rAEa{5},\rAEa{4},\rAEb{5}}}
  \proofstepwidestarbreak[1,2,3,8]{i\in\NatSeg n}{%
    \FormulaRefAuto{A\subseteq B\dsep x\in A\vdash x\in B}[thm]{9,8}}
  \proofstepwidestarbreak[1,2,3,8]{i+1\in\NatSeg{k+1}}{%
    \FormulaRefAuto{FinitePathIndexShiftOne}[thm]{\rAEa{5},8}}
  \proofstepwidestar[1,2,3,8]{k\leq k+1}{%
      \FormulaRefAuto{PeanoLeqAddOne}[thm]{\rAEa{5}}}
  \proofstepwidestar[1,2,3,8]{\NatSeg k \subseteq \NatSeg { { k + 1 } }}{%
      \FormulaRefAuto{PeanoNatSegLeqSubsetForward}[thm]{%
        \rAEa{5},\rAEa{5},12}}
  \proofstepwidestarbreak[1,2,3,8]{i\in\NatSeg{k+1}}{%
    \FormulaRefAuto{A\subseteq B\dsep x\in A\vdash x\in B}[thm]{%
      13,8}}
  \proofstepwidestarbreak[1,2,3,8]{%
    p\restriction_{\NatSeg{k+1}}(i)=p(i)}{%
    \FormulaRefAuto{C\subseteq A \dsep x\in C \vdash
      F\restriction_C(x)=F(x)}[thm]{\rAEn{5},14}}
  \proofstepwidestarbreak[1,2,3,8]{%
    p\restriction_{\NatSeg{k+1}}(i+1)=p(i+1)}{%
    \FormulaRefAuto{C\subseteq A \dsep x\in C \vdash
      F\restriction_C(x)=F(x)}[thm]{\rAEn{5},11}}
  \proofstepwidestarbreak[1,2,3,8]{(p(i),p(i+1))\in E}{%
    \FormulaRefAuto{DirectedWalkStepEdge}[thm]{1,2,10}}
  \proofstepwidestarbreak[1,2,3,8]{%
    (p\restriction_{\NatSeg{k+1}}(i),
      p\restriction_{\NatSeg{k+1}}(i+1))\in E}{\rIE{15,\rIE{16,17}}}
  \proofstepwidestarbreak[1,2,3]{%
    \forall i\in\NatSeg k\,
      \bigl((p\restriction_{\NatSeg{k+1}}(i),
        p\restriction_{\NatSeg{k+1}}(i+1))\in E\bigr)}{%
    \rUI{\rRI{8,18}}}
  \proofstepwidestarbreak[1,2,3]{%
    \DirectedWalk{V}{E}{p\restriction_{\NatSeg{k+1}}}{k}}{%
    \FormulaRefAuto{DirectedWalkCharacterization}[thm]{1,\rAEa{5},7,19}}
  \closeproofpartwide

  \proofpartwide{Erhaltung der Endpunkte}
  \setcounter{proofstepnr}{20}
  \proofstepwidestar[1,2,3]{0 \in \mathbb N}{%
      \FormulaRefAuto{PeanoZero}[ax]}
  \proofstepwidestar[1,2,3]{k\leq k}{%
      \FormulaRefAuto{PeanoLeqReflexive}[thm]{\rAEa{5}}}
  \proofstepwidestar[1,2,3]{k\leq k\leftrightarrow k\in\NatSeg{k+1}}{%
      \FormulaRefAuto{PeanoNatSegLeqAddOneMembership}[thm]{%
              \rAEa{5},\rAEa{5}}}
  \proofstepwidestar[1,2,3]{{ k \in { \NatSeg { k + 1 } } }}{%
      \FormulaRefAuto{P\leftrightarrow Q\dsep P\vdash Q}[thm]{%
            23,
            22}}
  \proofstepwidestar[1,2,3]{0\leq k}{%
      \FormulaRefAuto{PeanoZeroLeq}[thm]{\rAEa{5}}}
  \proofstepwidestar[1,2,3]{0\leq k\leftrightarrow 0\in\NatSeg{k+1}}{%
      \FormulaRefAuto{PeanoNatSegLeqAddOneMembership}[thm]{%
              21,\rAEa{5}}}
  \proofstepwidestar[1,2,3]{{ 0 \in { \NatSeg { k + 1 } } }}{%
      \FormulaRefAuto{P\leftrightarrow Q\dsep P\vdash Q}[thm]{%
            26,
            25}}
  \proofstepwidestar[1,2,3]{p \restriction _ { { \NatSeg { k + 1 } } } ( k ) = p ( k )}{%
      \FormulaRefAuto{C\subseteq A \dsep x\in C \vdash
        F\restriction_C(x)=F(x)}[thm]{\rAEn{5},
          24}}
  \proofstepwidestar[1,2,3]{p \restriction _ { { \NatSeg { k + 1 } } } ( 0 ) = p ( 0 )}{%
      \FormulaRefAuto{C\subseteq A \dsep x\in C \vdash
        F\restriction_C(x)=F(x)}[thm]{\rAEn{5},
          27}}
  \proofstepwidestarbreak[1,2,3]{%
    p\restriction_{\NatSeg{k+1}}(0)=p(0)\land
    p\restriction_{\NatSeg{k+1}}(k)=p(k)}{%
    \rAI{%
      29,
      28}}
  \proofstepwidestarbreak[1,2,3]{%
    \DirectedWalk{V}{E}{p\restriction_{\NatSeg{k+1}}}{k}
    \land p\restriction_{\NatSeg{k+1}}(0)=p(0)
    \land p\restriction_{\NatSeg{k+1}}(k)=p(k)}{%
    \rAI{20,\rAEa{30},\rAEb{30}}}
  \closeproofpartwide
\end{tabproofsplitwide}

\FormulaThmDeltaK[Anfangsabschnitt eines endlichen Pfads]{%
  \begin{aligned}[t]
  &\GraphStruct{V,E}\dsep\DirectedPath{V}{E}{p}{n}\dsep
    k\in\NatSeg{n+1}\vdash{}\\[-2pt]
  &\quad\DirectedPath{V}{E}{p\restriction_{\NatSeg{k+1}}}{k}
    \land p\restriction_{\NatSeg{k+1}}(0)=p(0)
    \land p\restriction_{\NatSeg{k+1}}(k)=p(k)
  \end{aligned}%
}{FinitePathInitialSegment}{%
  \DeltaRow{Mengen}{V\dsep E\dsep p\dsep q\dsep n\dsep k}
}
\begin{tabproofwide}
  \proofstepwidestarbreak[1]{\GraphStruct{V,E}}{\rA}
  \proofstepwidestarbreak[2]{\DirectedPath{V}{E}{p}{n}}{\rA}
  \proofstepwidestarbreak[3]{k\in\NatSeg{n+1}}{\rA}
  \proofstepwidestarbreak[1,2]{\DirectedWalk{V}{E}{p}{n}}{%
    \FormulaRefAuto{DirectedPathWalkAxiom}[ax]{1,2}}
  \proofstepwidestarbreak[1,2,3]{%
    \DirectedWalk{V}{E}{p\restriction_{\NatSeg{k+1}}}{k}
    \land p\restriction_{\NatSeg{k+1}}(0)=p(0)
    \land p\restriction_{\NatSeg{k+1}}(k)=p(k)}{%
    \FormulaRefAuto{FiniteWalkInitialSegment}[thm]{1,4,3}}
  \proofstepwidestar[1,2]{\DirectedWalk{V}{E}{p}{n}\land p\colon\NatSeg{n+1}\inj V}{%
      \FormulaRefAuto{DirectedPathCharacterization}[thm]{1,2}}
  \proofstepwidestarbreak[1,2]{p\colon\NatSeg{n+1}\inj V}{%
    \rAEb{6}}
  \proofstepwidestar[1,2,3]{\begin{aligned}[t]&n\in\mathbb N\land p\colon\NatSeg{n+1}\to V\land{}\\[-2pt]&\forall i\in\NatSeg{n}\,((p(i),p(i+1))\in E)\end{aligned}}{%
      \FormulaRefAuto{DirectedWalkCharacterization}[thm]{1,4}}
  \proofstepwidestar[1,2,3]{k\in\mathbb N\land k\leq n\land\NatSeg{k+1}\subseteq\NatSeg{n+1}}{%
      \FormulaRefAuto{FinitePathPrefixIndexData}[thm]{%
        \rAEa{8},3}}
  \proofstepwidestarbreak[1,2,3]{%
    p\restriction_{\NatSeg{k+1}}\colon\NatSeg{k+1}\inj V}{%
    \FormulaRefAuto{F\colon A\inj B\dsep C\subseteq A\vdash
      F\restriction_C\colon C\inj B}[thm]{7,
      \rAEn{9}}}
  \proofstepwidestarbreak[1,2,3]{%
    \DirectedPath{V}{E}{p\restriction_{\NatSeg{k+1}}}{k}}{%
    \FormulaRefAuto{DirectedPathCharacterization}[thm]{1,\rAEa{5},10}}
  \proofstepwidestarbreak[1,2,3]{%
    \DirectedPath{V}{E}{p\restriction_{\NatSeg{k+1}}}{k}
    \land p\restriction_{\NatSeg{k+1}}(0)=p(0)
    \land p\restriction_{\NatSeg{k+1}}(k)=p(k)}{%
    \rAI{11,\rAEa{5},\rAEb{5}}}
\end{tabproofwide}

\FormulaThmDeltaK[Schwanzsegment eines endlichen Pfads]{\begin{aligned}[t]
  &\GraphStruct{V,E}\dsep n\in\mathbb N\dsep
    \DirectedPath{V}{E}{p}{n+1}\vdash{}\\[-2pt]
  &\quad\exists q\in\powerset(\mathbb N\times V)\,\Bigl(
    \DirectedPath{V}{E}{q}{n}\land{}\\[-2pt]
  &\qquad q(0)=p(1)\land q(n)=p(n+1)\land{}\\[-2pt]
  &\qquad\forall i\in\NatSeg{n+1}\,(q(i)=p(i+1))\Bigr)
  \end{aligned}}{FinitePathTailSegment}{%
  \DeltaRow{Mengen}{V\dsep E\dsep p\dsep q\dsep n\dsep i\dsep j}
}
\begin{tabproofsplitwide}
  \proofpartwide{Konstruktion der verschobenen Folge}
  \proofstepwidestarbreak[1]{\GraphStruct{V,E}}{\rA}
  \proofstepwidestarbreak[2]{n\in\mathbb N}{\rA}
  \proofstepwidestarbreak[3]{\DirectedPath{V}{E}{p}{n+1}}{\rA}
  \proofstepwidestarbreak[1,3]{\DirectedWalk{V}{E}{p}{n+1}}{%
    \FormulaRefAuto{DirectedPathWalkAxiom}[ax]{1,3}}
  \proofstepwidestar[1,2,3]{\DirectedWalk{V}{E}{p}{(n+1)}\land p\colon\NatSeg{(n+1)+1}\inj V}{%
      \FormulaRefAuto{DirectedPathCharacterization}[thm]{1,3}}
  \proofstepwidestarbreak[1,2,3]{p\colon\NatSeg{(n+1)+1}\inj V}{%
    \rAEb{5}}
  \proofstepwidestarbreak[7]{i\in\NatSeg{n+1}}{\rA}
  \proofstepwidestar[2,7]{1 \in \mathbb N}{%
      \FormulaRefAuto{PeanoOneInN}[thm]}
  \proofstepwidestar[2,7]{n+1\in\mathbb N}{%
      \FormulaRefAuto{PeanoAddClosure}[thm]{2,
        8}}
  \proofstepwidestarbreak[2,7]{i+1\in\NatSeg{(n+1)+1}}{%
    \FormulaRefAuto{FinitePathIndexShiftOne}[thm]{%
      9,7}}
  \proofstepwidestar[1,2,3,7]{p \colon { \NatSeg { ( n + 1 ) + 1 } } \to V}{%
      \FormulaRefAuto{Injektive Funktion}[def]{6}}
  \proofstepwidestarbreak[1,2,3,7]{p(i+1)\in V}{%
    \FormulaRefAuto{F\colon A\to B\dsep x\in A\vdash F(x)\in B}[thm]{%
      11,10}}
  \proofstepwidestarbreak[1,2,3]{%
    \forall i\in\NatSeg{n+1}\,p(i+1)\in V}{\rUI{\rRI{7,12}}}
  \proofstepwidestarbreak[1,2,3]{%
    \exists!q\,\Bigl(q\colon\NatSeg{n+1}\to V\land
      \forall i\in\NatSeg{n+1}\,(q(i)=p(i+1))\Bigr)}{%
    \FormulaRefAuto{TypedFunctionDefinitionPrinciple}[thm]{13}}
  \proofstepwidestarbreak[15]{q\colon\NatSeg{n+1}\to V\land
    \forall i\in\NatSeg{n+1}\,(q(i)=p(i+1))}{\rA}
  \closeproofpartwide

  \proofpartwide{Injektivität}
  \setcounter{proofstepnr}{15}
  \proofstepwidestarbreak[16]{i\in\NatSeg{n+1}}{\rA}
  \proofstepwidestarbreak[17]{j\in\NatSeg{n+1}}{\rA}
  \proofstepwidestarbreak[18]{q(i)=q(j)}{\rA}
  \proofstepwidestar[2,16]{1 \in \mathbb N}{%
      \FormulaRefAuto{PeanoOneInN}[thm]}
  \proofstepwidestar[2,16]{n+1\in\mathbb N}{%
      \FormulaRefAuto{PeanoAddClosure}[thm]{2,
        19}}
  \proofstepwidestarbreak[2,16]{i+1\in\NatSeg{(n+1)+1}}{%
    \FormulaRefAuto{FinitePathIndexShiftOne}[thm]{%
      20,16}}
  \proofstepwidestar[2,17]{1 \in \mathbb N}{%
      \FormulaRefAuto{PeanoOneInN}[thm]}
  \proofstepwidestar[2,17]{n+1\in\mathbb N}{%
      \FormulaRefAuto{PeanoAddClosure}[thm]{2,
        22}}
  \proofstepwidestarbreak[2,17]{j+1\in\NatSeg{(n+1)+1}}{%
    \FormulaRefAuto{FinitePathIndexShiftOne}[thm]{%
      23,17}}
  \proofstepwidestarbreak[15,16]{q(i)=p(i+1)}{\rRE{16,\rUE{\rAEb{15}}}}
  \proofstepwidestarbreak[15,17]{q(j)=p(j+1)}{\rRE{17,\rUE{\rAEb{15}}}}
  \proofstepwidestarbreak[15,16,17,18]{p(i+1)=p(j+1)}{%
    \rIE{25,\rIE{26,18}}}
  \proofstepwidestarbreak[1,2,3,15,16,17,18]{i+1=j+1}{%
    \FormulaRefAuto{F\colon A\inj B\dsep x\in A\dsep y\in A\dsep
      F(x)=F(y)\vdash x=y}[ax]{6,21,24,27}}
  \proofstepwidestar[2,16]{1 \in \mathbb N}{%
      \FormulaRefAuto{PeanoOneInN}[thm]}
  \proofstepwidestar[2,16]{n + 1 \in \mathbb N}{%
      \FormulaRefAuto{PeanoAddClosure}[thm]{2,
          29}}
  \proofstepwidestar[2,16]{\NatSeg { ( n + 1 ) } \subseteq \mathbb N}{%
      \FormulaRefAuto{PeanoNatSegSubsetN}[thm]{%
        30}}
  \proofstepwidestarbreak[2,16]{i\in\mathbb N}{%
    \FormulaRefAuto{A\subseteq B\dsep x\in A\vdash x\in B}[thm]{%
      31,16}}
  \proofstepwidestar[2,17]{1 \in \mathbb N}{%
      \FormulaRefAuto{PeanoOneInN}[thm]}
  \proofstepwidestar[2,17]{n + 1 \in \mathbb N}{%
      \FormulaRefAuto{PeanoAddClosure}[thm]{2,
          33}}
  \proofstepwidestar[2,17]{\NatSeg { ( n + 1 ) } \subseteq \mathbb N}{%
      \FormulaRefAuto{PeanoNatSegSubsetN}[thm]{%
        34}}
  \proofstepwidestarbreak[2,17]{j\in\mathbb N}{%
    \FormulaRefAuto{A\subseteq B\dsep x\in A\vdash x\in B}[thm]{%
      35,17}}
  \proofstepwidestar[1,2,3,15,16,17,18]{1 \in \mathbb N}{%
      \FormulaRefAuto{PeanoOneInN}[thm]}
  \proofstepwidestarbreak[1,2,3,15,16,17,18]{i=j}{%
    \FormulaRefAuto{PeanoAddRightCancellation}[thm]{%
      32,36,37,28}}
  \proofstepwidestarbreak[1,2,3,15]{q\colon\NatSeg{n+1}\inj V}{%
    \FormulaRefAuto{Injektive Funktion}[def]{\rAEa{15},16,17,18,38}}
  \closeproofpartwide

  \proofpartwide{Kantenbedingung und Pfadstruktur}
  \setcounter{proofstepnr}{39}
  \proofstepwidestarbreak[40]{i\in\NatSeg n}{\rA}
  \proofstepwidestarbreak[2,40]{i+1\in\NatSeg{n+1}}{%
    \FormulaRefAuto{FinitePathIndexShiftOne}[thm]{2,40}}
  \proofstepwidestarbreak[1,2,3,40]{(p(i+1),p((i+1)+1))\in E}{%
    \FormulaRefAuto{DirectedWalkStepEdge}[thm]{1,4,41}}
  \proofstepwidestar[15,40]{1 \in \mathbb N}{%
      \FormulaRefAuto{PeanoOneInN}[thm]}
  \proofstepwidestar[15,40]{n + 1 \in \mathbb N}{%
      \FormulaRefAuto{PeanoAddClosure}[thm]{2,
            43}}
  \proofstepwidestar[15,40]{n\leq n+1}{%
      \FormulaRefAuto{PeanoLeqAddOne}[thm]{2}}
  \proofstepwidestar[15,40]{\NatSeg n \subseteq \NatSeg { ( n + 1 ) }}{%
      \FormulaRefAuto{PeanoNatSegLeqSubsetForward}[thm]{%
          2,44,
          45}}
  \proofstepwidestar[15,40]{i\in\NatSeg{n+1}}{%
      \FormulaRefAuto{A\subseteq B\dsep x\in A\vdash x\in B}[thm]{%
        46,40}}
  \proofstepwidestarbreak[15,40]{q(i)=p(i+1)}{%
    \rRE{%
      47,
      \rUE{\rAEb{15}}}}
  \proofstepwidestarbreak[15,2,40]{q(i+1)=p((i+1)+1)}{%
    \rRE{41,\rUE{\rAEb{15}}}}
  \proofstepwidestarbreak[1,2,3,15,40]{(q(i),q(i+1))\in E}{%
    \rIE{48,\rIE{49,42}}}
  \proofstepwidestarbreak[1,2,3,15]{%
    \forall i\in\NatSeg n\,((q(i),q(i+1))\in E)}{\rUI{\rRI{40,50}}}
  \proofstepwidestarbreak[1,2,3,15]{\DirectedWalk{V}{E}{q}{n}}{%
    \FormulaRefAuto{DirectedWalkCharacterization}[thm]{1,2,\rAEa{15},51}}
  \proofstepwidestarbreak[1,2,3,15]{\DirectedPath{V}{E}{q}{n}}{%
    \FormulaRefAuto{DirectedPathCharacterization}[thm]{1,52,39}}
  \closeproofpartwide

  \proofpartwide{Endpunkte und Existenzschluss}
  \setcounter{proofstepnr}{53}
  \proofstepwidestar[15]{1 \in \mathbb N}{%
      \FormulaRefAuto{PeanoOneInN}[thm]}
  \proofstepwidestar[15]{0 + 1 = 1}{%
      \FormulaRefAuto{PeanoAddLeftZero}[thm]{%
        54}}
  \proofstepwidestar[2]{0\in\mathbb N}{%
      \FormulaRefAuto{PeanoZero}[ax]}
  \proofstepwidestar[2]{0\leq n}{%
      \FormulaRefAuto{PeanoZeroLeq}[thm]{2}}
  \proofstepwidestar[2]{0\leq n\leftrightarrow0\in\NatSeg{n+1}}{%
      \FormulaRefAuto{PeanoNatSegLeqAddOneMembership}[thm]{56,2}}
  \proofstepwidestar[2]{0\in\NatSeg{n+1}}{\FormulaRefAuto{P\leftrightarrow Q\dsep P\vdash Q}[thm]{58,57}}
  \proofstepwidestarbreak[2,15]{q(0)=p(1)}{%
    \rIE{55,
      \rRE{%
        59,
        \rUE{\rAEb{15}}}}}
  \proofstepwidestar[2,15]{n\leq n}{%
      \FormulaRefAuto{PeanoLeqReflexive}[thm]{2}}
  \proofstepwidestar[2,15]{n\leq n\leftrightarrow n\in\NatSeg{n+1}}{%
      \FormulaRefAuto{PeanoNatSegLeqAddOneMembership}[thm]{2,2}}
  \proofstepwidestar[2,15]{n\in\NatSeg{n+1}}{%
      \FormulaRefAuto{P\leftrightarrow Q\dsep P\vdash Q}[thm]{%
        62,
        61}}
  \proofstepwidestarbreak[2,15]{q(n)=p(n+1)}{%
    \rRE{%
      63,
      \rUE{\rAEb{15}}}}
  \proofstepwidestarbreak[1,2,3,15]{q\in\powerset(\mathbb N\times V)}{%
    \FormulaRefAuto{FinitePathPowerSetTyping}[thm]{1,52}}
  \proofstepwidestarbreak[1,2,3,15]{%
    \exists q\in\powerset(\mathbb N\times V)\,\bigl(
      \DirectedPath{V}{E}{q}{n}\land q(0)=p(1)\land q(n)=p(n+1)\land
      \forall i\in\NatSeg{n+1}\,(q(i)=p(i+1))\bigr)}{%
    \rEI{\rAI{65,53,60,64,\rAEb{15}}}}
  \proofstepwidestar[1,2,3]{\exists q\,(q\colon\NatSeg{n+1}\to V\land\forall i\in\NatSeg{n+1}\,(q(i)=p(i+1)))}{%
      \FormulaRefAuto{\exists! x\,P(x)\vdash\exists x\,P(x)}[thm]{14}}
  \proofstepwidestarbreak[1,2,3]{%
    \exists q\in\powerset(\mathbb N\times V)\,\bigl(
      \DirectedPath{V}{E}{q}{n}\land q(0)=p(1)\land q(n)=p(n+1)\land
      \forall i\in\NatSeg{n+1}\,(q(i)=p(i+1))\bigr)}{%
    \rEE{67,15,66}}
  \closeproofpartwide
\end{tabproofsplitwide}

\FormulaThmDeltaK[Das Schwanzsegment meidet den ersten Knoten]{%
  \begin{aligned}[t]
  &\GraphStruct{V,E}\dsep n\in\mathbb N\dsep
    \DirectedPath{V}{E}{p}{n+1}\dsep
    q\colon\NatSeg{n+1}\to V\dsep{}\\[-2pt]
  &\forall i\in\NatSeg{n+1}\,(q(i)=p(i+1))
    \vdash p(0)\notin q[\NatSeg{n+1}]
  \end{aligned}%
}{FinitePathTailAvoidsFirstVertex}{%
  \DeltaRow{Mengen}{V\dsep E\dsep p\dsep q\dsep n\dsep i}
}
\begin{tabproofwide}
  \proofstepwidestarbreak[1]{\GraphStruct{V,E}}{\rA}
  \proofstepwidestarbreak[2]{n\in\mathbb N}{\rA}
  \proofstepwidestarbreak[3]{\DirectedPath{V}{E}{p}{n+1}}{\rA}
  \proofstepwidestarbreak[4]{q\colon\NatSeg{n+1}\to V}{\rA}
  \proofstepwidestarbreak[5]{\forall i\in\NatSeg{n+1}\,(q(i)=p(i+1))}{\rA}
  \proofstepwidestarbreak[6]{p(0)\in q[\NatSeg{n+1}]}{\rA}
  \proofstepwidestarbreak[4,6]{\exists i\in\NatSeg{n+1}\,(p(0)=q(i))}{%
    \FormulaRefAuto{F\colon A\to B\dsep y\in F[A]\vdash
      \exists x\in A\,y=F(x)}[thm]{4,6}}
  \proofstepwidestarbreak[8]{i\in\NatSeg{n+1}\land p(0)=q(i)}{\rA}
  \proofstepwidestarbreak[5,8]{q(i)=p(i+1)}{\rRE{\rAEa{8},\rUE{5}}}
  \proofstepwidestarbreak[5,8]{p(0)=p(i+1)}{\rIE{9,\rAEb{8}}}
  \proofstepwidestar[1,2,3]{\DirectedWalk{V}{E}{p}{(n+1)}\land p\colon\NatSeg{(n+1)+1}\inj V}{%
      \FormulaRefAuto{DirectedPathCharacterization}[thm]{1,3}}
  \proofstepwidestarbreak[1,2,3]{p\colon\NatSeg{(n+1)+1}\inj V}{%
    \rAEb{11}}
  \proofstepwidestar[2]{0 \in \mathbb N}{%
      \FormulaRefAuto{PeanoZero}[ax]}
  \proofstepwidestar[2]{1 \in \mathbb N}{%
      \FormulaRefAuto{PeanoOneInN}[thm]}
  \proofstepwidestar[2]{1 \in \mathbb N}{%
      \FormulaRefAuto{PeanoOneInN}[thm]}
  \proofstepwidestar[2]{n + 1 \in \mathbb N}{%
      \FormulaRefAuto{PeanoAddClosure}[thm]{2,
        14}}
  \proofstepwidestar[2]{n + 1 \in \mathbb N}{%
      \FormulaRefAuto{PeanoAddClosure}[thm]{2,
          15}}
  \proofstepwidestar[2]{0\leq n+1}{%
      \FormulaRefAuto{PeanoZeroLeq}[thm]{%
        17}}
  \proofstepwidestarbreak[2]{0\in\NatSeg{(n+1)+1}}{%
    \FormulaRefAuto{PeanoNatSegLeqAddOneMembership}[thm]{%
      13,
      16,
      18}}
  \proofstepwidestar[2,8]{1 \in \mathbb N}{%
      \FormulaRefAuto{PeanoOneInN}[thm]}
  \proofstepwidestar[2,8]{n+1\in\mathbb N}{%
      \FormulaRefAuto{PeanoAddClosure}[thm]{2,
        20}}
  \proofstepwidestarbreak[2,8]{i+1\in\NatSeg{(n+1)+1}}{%
    \FormulaRefAuto{FinitePathIndexShiftOne}[thm]{%
      21,\rAEa{8}}}
  \proofstepwidestarbreak[1,2,3,5,8]{0=i+1}{%
    \FormulaRefAuto{F\colon A\inj B\dsep x\in A\dsep y\in A\dsep
      F(x)=F(y)\vdash x=y}[ax]{12,19,22,10}}
  \proofstepwidestar[2,8]{1 \in \mathbb N}{%
      \FormulaRefAuto{PeanoOneInN}[thm]}
  \proofstepwidestar[2,8]{1 \in \mathbb N}{%
      \FormulaRefAuto{PeanoOneInN}[thm]}
  \proofstepwidestar[2,8]{1 \in \mathbb N}{%
      \FormulaRefAuto{PeanoOneInN}[thm]}
  \proofstepwidestar[2,8]{1 \in \mathbb N}{%
      \FormulaRefAuto{PeanoOneInN}[thm]}
  \proofstepwidestar[2,8]{n + 1 \in \mathbb N}{%
      \FormulaRefAuto{PeanoAddClosure}[thm]{2,
            24}}
  \proofstepwidestar[2,8]{n + 1 \in \mathbb N}{%
      \FormulaRefAuto{PeanoAddClosure}[thm]{2,
              25}}
  \proofstepwidestar[2,8]{n + 1 \in \mathbb N}{%
      \FormulaRefAuto{PeanoAddClosure}[thm]{2,
              27}}
  \proofstepwidestar[2,8]{\NatSeg { ( n + 1 ) } \subseteq \mathbb N}{%
      \FormulaRefAuto{PeanoNatSegSubsetN}[thm]{%
          28}}
  \proofstepwidestar[2,8]{\NatSeg { ( n + 1 ) } \subseteq \mathbb N}{%
      \FormulaRefAuto{PeanoNatSegSubsetN}[thm]{%
            29}}
  \proofstepwidestar[2,8]{\NatSeg { ( n + 1 ) } \subseteq \mathbb N}{%
      \FormulaRefAuto{PeanoNatSegSubsetN}[thm]{%
            30}}
  \proofstepwidestar[2,8]{i \in \mathbb N}{%
      \FormulaRefAuto{A\subseteq B\dsep x\in A\vdash x\in B}[thm]{%
          32,\rAEa{8}}}
  \proofstepwidestar[2,8]{i + 1 \in \mathbb N}{%
      \FormulaRefAuto{PeanoAddClosure}[thm]{%
        34,
        26}}
  \proofstepwidestar[2,8]{i\in\mathbb N}{%
      \FormulaRefAuto{A\subseteq B\dsep x\in A\vdash x\in B}[thm]{%
        31,\rAEa{8}}}
  \proofstepwidestar[2,8]{i\in\mathbb N}{%
      \FormulaRefAuto{A\subseteq B\dsep x\in A\vdash x\in B}[thm]{%
          33,\rAEa{8}}}
  \proofstepwidestar[2,8]{i<i+1}{%
      \FormulaRefAuto{PeanoLtAddOne}[thm]{%
        37}}
  \proofstepwidestarbreak[2,8]{i+1\neq0}{%
    \FormulaRefAuto{PeanoStrictUpperNonzero}[thm]{%
      36,
      35,
      38}}
  \proofstepwidestar[1,2,3,4,5,8]{{ ( i + 1 ) } = 0}{%
      \FormulaRefAuto{a=b\vdash b=a}[thm]{23}}
  \proofstepwidestarbreak[1,2,3,4,5,8]{\bot}{%
    \rBI{40,39}}
  \proofstepwidestarbreak[1,2,3,4,5,6]{\bot}{\rEE{7,8,41}}
  \proofstepwidestarbreak[1,2,3,4,5]{p(0)\notin q[\NatSeg{n+1}]}{%
    \rCI{6,42}}
\end{tabproofwide}

\FormulaThmDeltaK[Das Bild des Schwanzsegments bleibt im Pfadbild]{%
  \begin{aligned}[t]
  &n\in\mathbb N\dsep p\colon\NatSeg{(n+1)+1}\to V\dsep
    q\colon\NatSeg{n+1}\to V\dsep{}\\[-2pt]
  &\forall i\in\NatSeg{n+1}\,(q(i)=p(i+1))\vdash
    q[\NatSeg{n+1}]\subseteq p[\NatSeg{(n+1)+1}]
  \end{aligned}%
}{FinitePathTailImageSubset}{%
  \DeltaRow{Mengen}{V\dsep p\dsep q\dsep n\dsep i\dsep y}
}
\begin{tabproofwide}
  \proofstepwidestarbreak[1]{n\in\mathbb N}{\rA}
  \proofstepwidestarbreak[2]{p\colon\NatSeg{(n+1)+1}\to V}{\rA}
  \proofstepwidestarbreak[3]{q\colon\NatSeg{n+1}\to V}{\rA}
  \proofstepwidestarbreak[4]{\forall i\in\NatSeg{n+1}\,(q(i)=p(i+1))}{\rA}
  \proofstepwidestarbreak[5]{y\in q[\NatSeg{n+1}]}{\rA}
  \proofstepwidestarbreak[3,5]{\exists i\in\NatSeg{n+1}\,(y=q(i))}{%
    \FormulaRefAuto{F\colon A\to B\dsep y\in F[A]\vdash
      \exists x\in A\,y=F(x)}[thm]{3,5}}
  \proofstepwidestarbreak[7]{i\in\NatSeg{n+1}\land y=q(i)}{\rA}
  \proofstepwidestar[1,7]{1 \in \mathbb N}{%
      \FormulaRefAuto{PeanoOneInN}[thm]}
  \proofstepwidestar[1,7]{n+1\in\mathbb N}{%
      \FormulaRefAuto{PeanoAddClosure}[thm]{1,
        8}}
  \proofstepwidestarbreak[1,7]{i+1\in\NatSeg{(n+1)+1}}{%
    \FormulaRefAuto{FinitePathIndexShiftOne}[thm]{%
      9,\rAEa{7}}}
  \proofstepwidestarbreak[4,7]{q(i)=p(i+1)}{\rRE{\rAEa{7},\rUE{4}}}
  \proofstepwidestarbreak[2,4,7]{y\in p[\NatSeg{(n+1)+1}]}{%
    \FormulaRefAuto{F\colon A\to B\dsep\exists x\in A\,y=F(x)\vdash
      y\in F[A]}[thm]{2,\rEI{\rAI{10,\rIE{11,\rAEb{7}}}}}}
  \proofstepwidestarbreak[1,2,3,4,5]{y\in p[\NatSeg{(n+1)+1}]}{%
    \rEE{6,7,12}}
  \proofstepwidestarbreak[1,2,3,4]{%
    q[\NatSeg{n+1}]\subseteq p[\NatSeg{(n+1)+1}]}{%
    \FormulaRefAuto{SubsetIntroductionRule}[thm]{[5]\vdots 13}}
\end{tabproofwide}

\FormulaThmDeltaK[Frische Fortsetzung eines endlichen Pfads]{\begin{aligned}[t]
  &\GraphStruct{V,E}\dsep\DirectedPath{V}{E}{p}{n}\dsep y\in V\dsep{}\\[-2pt]
  &y\notin p[\NatSeg{n+1}]\dsep(p(n),y)\in E\vdash{}\\[-2pt]
  &\quad\exists q\in\powerset(\mathbb N\times V)\,\Bigl(
    \DirectedPath{V}{E}{q}{n+1}\land{}\\[-2pt]
  &\qquad q(0)=p(0)\land q(n+1)=y\land{}\\[-2pt]
  &\qquad q[\NatSeg{(n+1)+1}]=p[\NatSeg{n+1}]\cup\{y\}\Bigr)
  \end{aligned}}{FinitePathFreshEndpointExtension}{%
  \DeltaRow{Mengen}{V\dsep E\dsep p\dsep q\dsep n\dsep i\dsep y}
}
\begin{tabproofsplitwide}
  \proofpartwide{Konstruktion, Injektivität und Bild}
  \proofstepwidestarbreak[1]{\GraphStruct{V,E}}{\rA}
  \proofstepwidestarbreak[2]{\DirectedPath{V}{E}{p}{n}}{\rA}
  \proofstepwidestarbreak[3]{y\in V}{\rA}
  \proofstepwidestarbreak[4]{y\notin p[\NatSeg{n+1}]}{\rA}
  \proofstepwidestarbreak[5]{(p(n),y)\in E}{\rA}
  \proofstepwidestarbreak[1,2]{\DirectedWalk{V}{E}{p}{n}}{%
    \FormulaRefAuto{DirectedPathWalkAxiom}[ax]{1,2}}
  \proofstepwidestar[1,2]{\begin{aligned}[t]&n\in\mathbb N\land p\colon\NatSeg{n+1}\to V\land{}\\[-2pt]&\forall i\in\NatSeg{n}\,((p(i),p(i+1))\in E)\end{aligned}}{%
      \FormulaRefAuto{DirectedWalkCharacterization}[thm]{1,6}}
  \proofstepwidestarbreak[1,2]{n\in\mathbb N}{%
    \rAEa{7}}
  \proofstepwidestar[1,2]{\DirectedWalk{V}{E}{p}{n}\land p\colon\NatSeg{n+1}\inj V}{%
      \FormulaRefAuto{DirectedPathCharacterization}[thm]{1,2}}
  \proofstepwidestarbreak[1,2]{p\colon\NatSeg{n+1}\inj V}{%
    \rAEb{9}}
  \proofstepwidestarbreak[1,2]{p\colon\NatSeg{n+1}\bij p[\NatSeg{n+1}]}{%
    \FormulaRefAuto{F\colon A\inj B\vdash F\colon A\bij F[A]}[thm]{10}}
  \proofstepwidestar[1,2]{p \colon { \NatSeg { n + 1 } } \to V}{%
      \FormulaRefAuto{Injektive Funktion}[def]{10}}
  \proofstepwidestarbreak[1,2]{p[\NatSeg{n+1}]\subseteq V}{%
    \FormulaRefAuto{F\colon A\to B\vdash F[A]\subseteq B}[thm]{%
      12}}
  \proofstepwidestar[1,2,3,4]{1 \in \mathbb N}{%
      \FormulaRefAuto{PeanoOneInN}[thm]}
  \proofstepwidestar[1,2,3,4]{n + 1 \in \mathbb N}{%
      \FormulaRefAuto{PeanoAddClosure}[thm]{8,
          14}}
  \proofstepwidestar[1,2,3,4]{{ ( n + 1 ) } \notin \NatSeg { ( n + 1 ) }}{%
      \FormulaRefAuto{PeanoNatSegEndpointNotIn}[thm]{%
        15}}
  \proofstepwidestarbreak[1,2,3,4]{%
    \ExtMap{\iota_{p[\NatSeg{n+1}],V}\circ p}{n+1}{y}
      \colon\NatSeg{n+1}\cup\{n+1\}
        \bij p[\NatSeg{n+1}]\cup\{y\}}{%
    \FormulaRefAuto{FreshExtMapBijective}[thm]{%
      11,
      16,4}}
  \proofstepwidestar[1,2]{1 \in \mathbb N}{%
      \FormulaRefAuto{PeanoOneInN}[thm]}
  \proofstepwidestar[1,2]{n + 1 \in \mathbb N}{%
      \FormulaRefAuto{PeanoAddClosure}[thm]{8,
        18}}
  \proofstepwidestarbreak[1,2]{%
    \NatSeg{(n+1)+1}=\NatSeg{n+1}\cup\{n+1\}}{%
    \FormulaRefAuto{PeanoNatSegAddOne}[thm]{%
      19}}
  \proofstepwidestarbreak[1,2,3,4]{%
    \ExtMap{\iota_{p[\NatSeg{n+1}],V}\circ p}{n+1}{y}
      \colon\NatSeg{(n+1)+1}\bij p[\NatSeg{n+1}]\cup\{y\}}{%
    \rIE{17,20}}
  \proofstepwidestar[1,2,3]{\{ y \} \subseteq V}{%
      \FormulaRefAuto{a\in A\vdash\{a\}\subseteq A}[thm]{3}}
  \proofstepwidestarbreak[1,2,3]{p[\NatSeg{n+1}]\cup\{y\}\subseteq V}{%
    \FormulaRefAuto{A\subseteq C\dsep B\subseteq C\vdash A\cup B\subseteq C}[thm]{%
      13,22}}
  \proofstepwidestar[1,2,3,4]{\ExtMap{\iota_{p[\NatSeg{n+1}],V}\circ p}{n+1}{y}\colon\NatSeg{(n+1)+1}\inj (p[\NatSeg{n+1}]\cup\{y\})}{%
      \FormulaRefAuto{Bijektive Funktion}[def]{21}}
  \proofstepwidestarbreak[1,2,3,4]{%
    \ExtMap{\iota_{p[\NatSeg{n+1}],V}\circ p}{n+1}{y}
      \colon\NatSeg{(n+1)+1}\inj V}{%
    \FormulaRefAuto{F\colon A\inj B\dsep B\subseteq C\vdash
      F\colon A\inj C}[thm]{%
        24,23}}
  \proofstepwidestarbreak[1,2,3,4]{%
    \ExtMap{\iota_{p[\NatSeg{n+1}],V}\circ p}{n+1}{y}
      [\NatSeg{(n+1)+1}]=p[\NatSeg{n+1}]\cup\{y\}}{%
    \FormulaRefAuto{F\colon A\bij B\vdash F[A]=B}[thm]{21}}
  \closeproofpartwide

  \proofpartwide{Werte auf dem alten und dem neuen Index}
  \setcounter{proofstepnr}{26}
  \proofstepwidestarbreak[27]{i\in\NatSeg{n+1}}{\rA}
  \proofstepwidestar[1,2,3,4,27]{1 \in \mathbb N}{%
      \FormulaRefAuto{PeanoOneInN}[thm]}
  \proofstepwidestar[1,2,3,4,27]{p \colon { \NatSeg { n + 1 } } \to { p [ \NatSeg { n + 1 } ] }}{%
      \FormulaRefAuto{Bijektive Funktion}[def]{11}}
  \proofstepwidestar[1,2,3,4,27]{n+1\in\mathbb N}{%
      \FormulaRefAuto{PeanoAddClosure}[thm]{8,
          28}}
  \proofstepwidestar[1,2,3,4,27]{{ { n + 1 } } \notin \NatSeg { { n + 1 } }}{%
      \FormulaRefAuto{PeanoNatSegEndpointNotIn}[thm]{%
        30}}
  \proofstepwidestarbreak[1,2,3,4,27]{%
    \ExtMap{\iota_{p[\NatSeg{n+1}],V}\circ p}{n+1}{y}(i)=p(i)}{%
    \FormulaRefAuto{FreshExtMapOnBase}[thm]{13,
      29,
      31,3,27}}
  \proofstepwidestar[1,2,3,4]{1 \in \mathbb N}{%
      \FormulaRefAuto{PeanoOneInN}[thm]}
  \proofstepwidestar[1,2,3,4]{p \colon { \NatSeg { n + 1 } } \to { p [ \NatSeg { n + 1 } ] }}{%
      \FormulaRefAuto{Bijektive Funktion}[def]{11}}
  \proofstepwidestar[1,2,3,4]{n + 1 \in \mathbb N}{%
      \FormulaRefAuto{PeanoAddClosure}[thm]{8,
          33}}
  \proofstepwidestar[1,2,3,4]{{ ( n + 1 ) } \notin \NatSeg { ( n + 1 ) }}{%
      \FormulaRefAuto{PeanoNatSegEndpointNotIn}[thm]{%
        35}}
  \proofstepwidestarbreak[1,2,3,4]{%
    \ExtMap{\iota_{p[\NatSeg{n+1}],V}\circ p}{n+1}{y}(n+1)=y}{%
    \FormulaRefAuto{FreshExtMapAtNew}[thm]{13,
      34,
      36,3}}
  \closeproofpartwide

  \proofpartwide{Kantenbedingung und Pfadstruktur}
  \setcounter{proofstepnr}{37}
  \proofstepwidestarbreak[38]{i\in\NatSeg{n+1}}{\rA}
  \proofstepwidestar[8,38]{\NatSeg { ( n + 1 ) } = \NatSeg n \cup \{ n \}}{%
      \FormulaRefAuto{PeanoNatSegAddOne}[thm]{8}}
  \proofstepwidestarbreak[8,38]{i\in\NatSeg n\lor i=n}{%
    \rIE{39,25}}
  \proofstepwidestarbreak[41]{i\in\NatSeg n}{\rA}
  \proofstepwidestarbreak[1,2,41]{(p(i),p(i+1))\in E}{%
    \FormulaRefAuto{DirectedWalkStepEdge}[thm]{1,6,41}}
  \proofstepwidestar[1,2,3,4,41]{1 \in \mathbb N}{%
      \FormulaRefAuto{PeanoOneInN}[thm]}
  \proofstepwidestar[1,2,3,4,41]{p \colon { \NatSeg { n + 1 } } \to { p [ \NatSeg { n + 1 } ] }}{%
      \FormulaRefAuto{Bijektive Funktion}[def]{11}}
  \proofstepwidestar[1,2,3,4,41]{i + 1 \in \NatSeg { n + 1 }}{%
      \FormulaRefAuto{FinitePathIndexShiftOne}[thm]{8,41}}
  \proofstepwidestar[1,2,3,4,41]{n + 1 \in \mathbb N}{%
      \FormulaRefAuto{PeanoAddClosure}[thm]{8,
            43}}
  \proofstepwidestar[1,2,3,4,41]{{ ( n + 1 ) } \notin \NatSeg { ( n + 1 ) }}{%
      \FormulaRefAuto{PeanoNatSegEndpointNotIn}[thm]{%
          46}}
  \proofstepwidestar[1,2,3,4,41]{\ExtMap{\iota_{p[\NatSeg{n+1}],V}\circ p}{n+1}{y}(i+1)=(\iota_{p[\NatSeg{n+1}],V}\circ p)(i+1)}{%
      \FormulaRefAuto{FreshExtMapOnBase}[thm]{13,
        44,
        47,3,
        45}}
  \proofstepwidestarbreak[1,2,3,4,41]{\begin{aligned}[t]
      &\bigl(\ExtMap{\iota_{p[\NatSeg{n+1}],V}\circ p}{n+1}{y}(i),\\[-2pt]
      &\quad \ExtMap{\iota_{p[\NatSeg{n+1}],V}\circ p}{n+1}{y}(i+1)\bigr)\in E
    \end{aligned}}{%
    \rIE{32,\rIE{%
      48,42}}}
  \proofstepwidestarbreak[50]{i=n}{\rA}
  \proofstepwidestarbreak[1,2,3,4,5,50]{\begin{aligned}[t]
      &\bigl(\ExtMap{\iota_{p[\NatSeg{n+1}],V}\circ p}{n+1}{y}(i),\\[-2pt]
      &\quad \ExtMap{\iota_{p[\NatSeg{n+1}],V}\circ p}{n+1}{y}(i+1)\bigr)\in E
    \end{aligned}}{%
    \rIE{50,\rIE{32,\rIE{37,5}}}}
  \proofstepwidestarbreak[1,2,3,4,5,38]{\begin{aligned}[t]
      &\bigl(\ExtMap{\iota_{p[\NatSeg{n+1}],V}\circ p}{n+1}{y}(i),\\[-2pt]
      &\quad \ExtMap{\iota_{p[\NatSeg{n+1}],V}\circ p}{n+1}{y}(i+1)\bigr)\in E
    \end{aligned}}{%
    \rOE{40,41,49,50,51}}
  \proofstepwidestarbreak[1,2,3,4,5]{\begin{aligned}[t]
      &%
    \forall i\in\NatSeg{n+1}\,{}\\[-2pt]
      &\bigl(\ExtMap{\iota_{p[\NatSeg{n+1}],V}\circ p}{n+1}{y}(i),\\[-2pt]
      &\quad \ExtMap{\iota_{p[\NatSeg{n+1}],V}\circ p}{n+1}{y}(i+1)\bigr)\in E
    \end{aligned}}{%
    \rUI{\rRI{38,52}}}
  \proofstepwidestar[1,2,3,4,5]{1 \in \mathbb N}{%
      \FormulaRefAuto{PeanoOneInN}[thm]}
  \proofstepwidestar[1,2,3,4,5]{{ \ExtMap { \iota _ { p [ \NatSeg { n + 1 } ] , V } \circ p } { n + 1 } y } \colon { \NatSeg { ( n + 1 ) + 1 } } \to V}{%
      \FormulaRefAuto{Injektive Funktion}[def]{25}}
  \proofstepwidestar[1,2,3,4,5]{n + 1 \in \mathbb N}{%
      \FormulaRefAuto{PeanoAddClosure}[thm]{8,
        54}}
  \proofstepwidestarbreak[1,2,3,4,5]{%
    \DirectedWalk{V}{E}{%
      \ExtMap{\iota_{p[\NatSeg{n+1}],V}\circ p}{n+1}{y}}{n+1}}{%
    \FormulaRefAuto{DirectedWalkCharacterization}[thm]{1,
      56,
      55,53}}
  \proofstepwidestarbreak[1,2,3,4,5]{%
    \DirectedPath{V}{E}{%
      \ExtMap{\iota_{p[\NatSeg{n+1}],V}\circ p}{n+1}{y}}{n+1}}{%
    \FormulaRefAuto{DirectedPathCharacterization}[thm]{1,57,25}}
  \closeproofpartwide

  \proofpartwide{Endpunkte und Existenzschluss}
  \setcounter{proofstepnr}{58}
  \proofstepwidestarbreak[1,2,3,4]{%
    \ExtMap{\iota_{p[\NatSeg{n+1}],V}\circ p}{n+1}{y}
      \in\powerset(\mathbb N\times V)}{%
    \FormulaRefAuto{FinitePathPowerSetTyping}[thm]{1,58}}
  \proofstepwidestar[1,2,3,4]{0 \in \mathbb N}{%
      \FormulaRefAuto{PeanoZero}[ax]}
  \proofstepwidestar[1,2,3,4]{1 \in \mathbb N}{%
      \FormulaRefAuto{PeanoOneInN}[thm]}
  \proofstepwidestar[1,2,3,4]{p \colon { \NatSeg { n + 1 } } \to { p [ \NatSeg { n + 1 } ] }}{%
      \FormulaRefAuto{Bijektive Funktion}[def]{11}}
  \proofstepwidestar[1,2,3,4]{n+1\in\mathbb N}{%
      \FormulaRefAuto{PeanoAddClosure}[thm]{8,
          61}}
  \proofstepwidestar[1,2,3,4]{{ { n + 1 } } \notin \NatSeg { { n + 1 } }}{%
      \FormulaRefAuto{PeanoNatSegEndpointNotIn}[thm]{%
        63}}
  \proofstepwidestar[1,2,3,4]{0\leq n}{%
      \FormulaRefAuto{PeanoZeroLeq}[thm]{8}}
  \proofstepwidestar[1,2,3,4]{0\in\NatSeg{n+1}}{%
      \FormulaRefAuto{PeanoNatSegLeqAddOneMembership}[thm]{%
        60,8,
        65}}
  \proofstepwidestarbreak[1,2,3,4]{%
    \ExtMap{\iota_{p[\NatSeg{n+1}],V}\circ p}{n+1}{y}(0)=p(0)}{%
    \FormulaRefAuto{FreshExtMapOnBase}[thm]{13,
      62,
      64,3,
      66}}
  \proofstepwidestarbreak[1,2,3,4,5]{%
    \exists q\in\powerset(\mathbb N\times V)\,\bigl(
      \DirectedPath{V}{E}{q}{n+1}\land q(0)=p(0)\land q(n+1)=y\land
      q[\NatSeg{(n+1)+1}]=p[\NatSeg{n+1}]\cup\{y\}\bigr)}{%
    \rEI{\rAI{59,58,67,37,26}}}
  \closeproofpartwide
\end{tabproofsplitwide}

\FormulaThmDeltaK[Fortsetzung eines Pfads um eine Kante]{%
  \begin{aligned}[t]
  &\GraphStruct{V,E}\dsep\DirectedPath{V}{E}{p}{n}\dsep
    (p(n),y)\in E\vdash{}\\[-2pt]
  &\quad\exists m\in\mathbb N\,\exists q\in\powerset(\mathbb N\times V)\,
    \bigl(\DirectedPath{V}{E}{q}{m}\land q(0)=p(0)\land q(m)=y\bigr)
  \end{aligned}%
}{FinitePathEndpointStep}{%
  \DeltaRow{Mengen}{V\dsep E\dsep p\dsep q\dsep n\dsep m\dsep i\dsep y}
}
\begin{tabproofwide}
  \proofstepwidestarbreak[1]{\GraphStruct{V,E}}{\rA}
  \proofstepwidestarbreak[2]{\DirectedPath{V}{E}{p}{n}}{\rA}
  \proofstepwidestarbreak[3]{(p(n),y)\in E}{\rA}
  \proofstepwidestarbreak[1,2]{\DirectedWalk{V}{E}{p}{n}}{%
    \FormulaRefAuto{DirectedPathWalkAxiom}[ax]{1,2}}
  \proofstepwidestar[1,2]{\begin{aligned}[t]&n\in\mathbb N\land p\colon\NatSeg{n+1}\to V\land{}\\[-2pt]&\forall i\in\NatSeg{n}\,((p(i),p(i+1))\in E)\end{aligned}}{%
      \FormulaRefAuto{DirectedWalkCharacterization}[thm]{1,4}}
  \proofstepwidestarbreak[1,2]{n\in\mathbb N}{%
    \rAEa{5}}
  \proofstepwidestar[1,3]{{ p ( n ) } \in V \land y \in V}{%
      \FormulaRefAuto{GraphEdgeEndpointTyping}[thm]{1,3}}
  \proofstepwidestarbreak[1,3]{y\in V}{%
    \rAEb{7}}
  \proofstepwidestar[1,2]{\begin{aligned}[t]&n\in\mathbb N\land p\colon\NatSeg{n+1}\to V\land{}\\[-2pt]&\forall i\in\NatSeg{n}\,((p(i),p(i+1))\in E)\end{aligned}}{%
      \FormulaRefAuto{DirectedWalkCharacterization}[thm]{1,4}}
  \proofstepwidestarbreak[1,2]{p\colon\NatSeg{n+1}\to V}{%
    \FormulaRefAuto{P\land Q\land R\vdash Q}[thm]{%
      9}}
  \proofstepwidestarbreak[]{y\in p[\NatSeg{n+1}]\lor
    y\notin p[\NatSeg{n+1}]}{\FormulaRefAuto{P\lor\neg P}[thm]}
  \proofstepwidestarbreak[12]{y\in p[\NatSeg{n+1}]}{\rA}
  \proofstepwidestarbreak[1,2,12]{%
    \exists i\in\NatSeg{n+1}\,(y=p(i))}{%
    \FormulaRefAuto{F\colon A\to B\dsep y\in F[A]\vdash
      \exists x\in A\,y=F(x)}[thm]{10,12}}
  \proofstepwidestarbreak[14]{i\in\NatSeg{n+1}\land y=p(i)}{\rA}
  \proofstepwidestarbreak[1,2,14]{%
    \DirectedPath{V}{E}{p\restriction_{\NatSeg{i+1}}}{i}\land
    p\restriction_{\NatSeg{i+1}}(0)=p(0)\land
    p\restriction_{\NatSeg{i+1}}(i)=p(i)}{%
    \FormulaRefAuto{FinitePathInitialSegment}[thm]{1,2,\rAEa{14}}}
  \proofstepwidestar[1,2,14]{i\in\mathbb N\land i\leq n\land\NatSeg{i+1}\subseteq\NatSeg{n+1}}{%
      \FormulaRefAuto{FinitePathPrefixIndexData}[thm]{6,\rAEa{14}}}
  \proofstepwidestarbreak[1,2,14]{i\in\mathbb N}{%
    \rAEa{16}}
  \proofstepwidestarbreak[1,2,14]{%
    p\restriction_{\NatSeg{i+1}}\in\powerset(\mathbb N\times V)}{%
    \FormulaRefAuto{FinitePathPowerSetTyping}[thm]{1,\rAEa{15}}}
  \proofstepwidestarbreak[1,2,14]{%
    \exists m\in\mathbb N\,\exists q\in\powerset(\mathbb N\times V)\,
    \bigl(\DirectedPath{V}{E}{q}{m}\land q(0)=p(0)\land q(m)=y\bigr)}{%
    \rEI{\rAI{17,\rEI{\rAI{18,\rAEa{15},\rAEb{15},
      \rIE{\rAEb{14},\rAEn{15}}}}}}}
  \proofstepwidestarbreak[1,2,12]{%
    \exists m\in\mathbb N\,\exists q\in\powerset(\mathbb N\times V)\,
    \bigl(\DirectedPath{V}{E}{q}{m}\land q(0)=p(0)\land q(m)=y\bigr)}{%
    \rEE{13,14,19}}
  \proofstepwidestarbreak[21]{y\notin p[\NatSeg{n+1}]}{\rA}
  \proofstepwidestarbreak[1,2,3,8,21]{%
    \exists q\in\powerset(\mathbb N\times V)\,\bigl(
      \DirectedPath{V}{E}{q}{n+1}\land q(0)=p(0)\land q(n+1)=y\bigr)}{%
    \FormulaRefAuto{FinitePathFreshEndpointExtension}[thm]{1,2,8,21,3}}
  \proofstepwidestarbreak[23]{q\in\powerset(\mathbb N\times V)\land
    \DirectedPath{V}{E}{q}{n+1}\land q(0)=p(0)\land q(n+1)=y\land
    q[\NatSeg{(n+1)+1}]=p[\NatSeg{n+1}]\cup\{y\}}{\rA}
  \proofstepwidestar[1,2,3,21,23]{1 \in \mathbb N}{%
      \FormulaRefAuto{PeanoOneInN}[thm]}
  \proofstepwidestar[1,2,3,21,23]{n + 1 \in \mathbb N}{%
      \FormulaRefAuto{PeanoAddClosure}[thm]{6,24}}
  \proofstepwidestarbreak[1,2,3,21,23]{%
    \exists m\in\mathbb N\,\exists q\in\powerset(\mathbb N\times V)\,
    \bigl(\DirectedPath{V}{E}{q}{m}\land q(0)=p(0)\land q(m)=y\bigr)}{%
    \rEI{\rAI{%
      25,
      \rEI{23}}}}
  \proofstepwidestarbreak[1,2,3,8,21]{%
    \exists m\in\mathbb N\,\exists q\in\powerset(\mathbb N\times V)\,
    \bigl(\DirectedPath{V}{E}{q}{m}\land q(0)=p(0)\land q(m)=y\bigr)}{%
    \rEE{22,23,26}}
  \proofstepwidestarbreak[1,2,3]{%
    \exists m\in\mathbb N\,\exists q\in\powerset(\mathbb N\times V)\,
    \bigl(\DirectedPath{V}{E}{q}{m}\land q(0)=p(0)\land q(m)=y\bigr)}{%
    \rOE{11,12,20,21,27}}
\end{tabproofwide}

\FormulaDefDeltaK[Transporteigenschaft für Walklängen]{\begin{aligned}[t]
  &\FinitePathTransportProperty{V}{E}{n}\coloneqq{}\\[-2pt]
  &\quad\forall p\in\powerset(\mathbb N\times V)\,\forall m\in\mathbb N\,
    \forall q\in\powerset(\mathbb N\times V)\,\Bigl(\\[-2pt]
  &\qquad\DirectedWalk{V}{E}{p}{n}\land\DirectedPath{V}{E}{q}{m}\land
    q(m)=p(0)\rightarrow{}\\[-2pt]
  &\qquad\exists l\in\mathbb N\,\exists s\in\powerset(\mathbb N\times V)\,
    \bigl(\DirectedPath{V}{E}{s}{l}\land{}\\[-2pt]
  &\qquad\quad s(0)=q(0)\land s(l)=p(n)\bigr)\Bigr)
  \end{aligned}}{FinitePathTransportPropertyDef}{%
  \DeltaRow{Mengen}{V\dsep E\dsep n\dsep m\dsep l\dsep p\dsep q\dsep s}
  \DeltaPrem{Graph}{\GraphStruct{V,E}}
  \DeltaRow{Neue Eigenschaft}{\FinitePathTransportProperty{V}{E}{n}}
}

\FormulaThmDeltaK[Transportbasis]{%
  \GraphStruct{V,E}\vdash\FinitePathTransportProperty{V}{E}{0}%
}{FinitePathTransportBase}{%
  \DeltaRow{Mengen}{V\dsep E\dsep m\dsep p\dsep q}
}
\begin{tabproofwide}
  \proofstepwidestarbreak[1]{\GraphStruct{V,E}}{\rA}
  \proofstepwidestarbreak[2]{p\in\powerset(\mathbb N\times V)}{\rA}
  \proofstepwidestarbreak[3]{m\in\mathbb N}{\rA}
  \proofstepwidestarbreak[4]{q\in\powerset(\mathbb N\times V)}{\rA}
  \proofstepwidestarbreak[5]{\DirectedWalk{V}{E}{p}{0}\land
    \DirectedPath{V}{E}{q}{m}\land q(m)=p(0)}{\rA}
  \proofstepwidestarbreak[3,4,5]{%
    \exists l\in\mathbb N\,\exists s\in\powerset(\mathbb N\times V)\,
    (\DirectedPath{V}{E}{s}{l}\land s(0)=q(0)\land s(l)=p(0))}{%
    \rEI{\rAI{3,\rEI{\rAI{4,\rAEb{5},\rII,\rAEn{5}}}}}}
  \proofstepwidestarbreak[1]{\FinitePathTransportProperty{V}{E}{0}}{%
    \FormulaRefAuto{FinitePathTransportPropertyDef}[def]{%
      \rUI{\rRI{2,\rUI{\rRI{3,\rUI{\rRI{4,\rRI{5,6}}}}}}}}}
\end{tabproofwide}

\FormulaThmDeltaK[Transportschritt]{%
  \GraphStruct{V,E}\dsep n\in\mathbb N\dsep
  \FinitePathTransportProperty{V}{E}{n}\vdash
  \FinitePathTransportProperty{V}{E}{n+1}%
}{FinitePathTransportStep}{%
  \DeltaRow{Mengen}{V\dsep E\dsep n\dsep m\dsep l\dsep h\dsep
    p\dsep q\dsep s\dsep t}
}
\begin{tabproofwide}
  \proofstepwidestarbreak[1]{\GraphStruct{V,E}}{\rA}
  \proofstepwidestarbreak[2]{n\in\mathbb N}{\rA}
  \proofstepwidestarbreak[3]{\FinitePathTransportProperty{V}{E}{n}}{\rA}
  \proofstepwidestarbreak[4]{p\in\powerset(\mathbb N\times V)}{\rA}
  \proofstepwidestarbreak[5]{m\in\mathbb N}{\rA}
  \proofstepwidestarbreak[6]{q\in\powerset(\mathbb N\times V)}{\rA}
  \proofstepwidestarbreak[7]{\DirectedWalk{V}{E}{p}{n+1}\land
    \DirectedPath{V}{E}{q}{m}\land q(m)=p(0)}{\rA}
  \proofstepwidestarbreak[2]{n\leq n+1}{\FormulaRefAuto{PeanoLeqAddOne}[thm]{2}}
  \proofstepwidestar[2]{1 \in \mathbb N}{%
      \FormulaRefAuto{PeanoOneInN}[thm]}
  \proofstepwidestarbreak[2]{n+1\in\mathbb N}{%
    \FormulaRefAuto{PeanoAddClosure}[thm]{2,9}}
  \proofstepwidestarbreak[2]{n\in\NatSeg{(n+1)+1}}{%
    \FormulaRefAuto{PeanoNatSegLeqAddOneMembership}[thm]{2,10,8}}
  \proofstepwidestarbreak[1,2,7]{%
    \DirectedWalk{V}{E}{p\restriction_{\NatSeg{n+1}}}{n}\land
    p\restriction_{\NatSeg{n+1}}(0)=p(0)\land
    p\restriction_{\NatSeg{n+1}}(n)=p(n)}{%
    \FormulaRefAuto{FiniteWalkInitialSegment}[thm]{1,\rAEa{7},11}}
  \proofstepwidestarbreak[1,2,3,5,6,7]{%
    \exists l\in\mathbb N\,\exists s\in\powerset(\mathbb N\times V)\,
    (\DirectedPath{V}{E}{s}{l}\land s(0)=q(0)\land s(l)=p(n))}{%
    \FormulaRefAuto{FinitePathTransportPropertyDef}[def]{%
      3,\rAEa{12},\rAEb{7},\rIE{\rAEb{12},\rAEn{7}}}}
  \proofstepwidestarbreak[14]{l\in\mathbb N\land
    s\in\powerset(\mathbb N\times V)\land
    \DirectedPath{V}{E}{s}{l}\land s(0)=q(0)\land s(l)=p(n)}{\rA}
  \proofstepwidestar[2]{n\leq n}{%
      \FormulaRefAuto{PeanoLeqReflexive}[thm]{2}}
  \proofstepwidestarbreak[2]{n\in\NatSeg{n+1}}{%
    \FormulaRefAuto{PeanoNatSegLeqAddOneMembership}[thm]{%
      2,2,15}}
  \proofstepwidestarbreak[1,2,7]{(p(n),p(n+1))\in E}{%
    \FormulaRefAuto{DirectedWalkStepEdge}[thm]{1,\rAEa{7},16}}
  \proofstepwidestarbreak[1,2,7,14]{%
    \exists h\in\mathbb N\,\exists t\in\powerset(\mathbb N\times V)\,
    (\DirectedPath{V}{E}{t}{h}\land t(0)=s(0)\land t(h)=p(n+1))}{%
    \FormulaRefAuto{FinitePathEndpointStep}[thm]{%
      1,\rAEn{14},\rIE{\rAEn{14},17}}}
  \proofstepwidestarbreak[19]{h\in\mathbb N\land
    t\in\powerset(\mathbb N\times V)\land
    \DirectedPath{V}{E}{t}{h}\land t(0)=s(0)\land t(h)=p(n+1)}{\rA}
  \proofstepwidestarbreak[1,2,3,5,6,7,14,19]{%
    \exists h\in\mathbb N\,\exists t\in\powerset(\mathbb N\times V)\,
    (\DirectedPath{V}{E}{t}{h}\land t(0)=q(0)\land t(h)=p(n+1))}{%
    \rEI{\rAI{\rAEa{19},\rEI{\rAI{\rAEb{19},\rAEn{19},
      \rIE{\rAEn{14},\rAEb{\rAEn{19}}},\rAEn{19}}}}}}
  \proofstepwidestarbreak[1,2,3,5,6,7]{%
    \exists h\in\mathbb N\,\exists t\in\powerset(\mathbb N\times V)\,
    (\DirectedPath{V}{E}{t}{h}\land t(0)=q(0)\land t(h)=p(n+1))}{%
    \rEE{13,14,\rEE{18,19,20}}}
  \proofstepwidestarbreak[1,2,3]{\FinitePathTransportProperty{V}{E}{n+1}}{%
    \FormulaRefAuto{FinitePathTransportPropertyDef}[def]{%
      \rUI{\rRI{4,\rUI{\rRI{5,\rUI{\rRI{6,\rRI{7,21}}}}}}}}}
\end{tabproofwide}

\FormulaThmDeltaK[Transport entlang eines endlichen Walks]{%
  \begin{aligned}[t]
  &\GraphStruct{V,E}\dsep n,m\in\mathbb N\dsep
    \DirectedWalk{V}{E}{p}{n}\dsep\DirectedPath{V}{E}{q}{m}\dsep
    q(m)=p(0)\vdash{}\\[-2pt]
  &\quad\exists l\in\mathbb N\,\exists s\in\powerset(\mathbb N\times V)\,
    \bigl(\DirectedPath{V}{E}{s}{l}\land s(0)=q(0)\land s(l)=p(n)\bigr)
  \end{aligned}%
}{FinitePathTransportAlongWalk}{%
  \DeltaRow{Mengen}{V\dsep E\dsep n\dsep m\dsep l\dsep p\dsep q\dsep s}
}
\begin{tabproofwide}
  \proofstepwidestarbreak[1]{\GraphStruct{V,E}}{\rA}
  \proofstepwidestarbreak[2]{n,m\in\mathbb N}{\rA}
  \proofstepwidestarbreak[3]{\DirectedWalk{V}{E}{p}{n}}{\rA}
  \proofstepwidestarbreak[4]{\DirectedPath{V}{E}{q}{m}}{\rA}
  \proofstepwidestarbreak[5]{q(m)=p(0)}{\rA}
  \proofstepwidestarbreak[1]{\FinitePathTransportProperty{V}{E}{0}}{%
    \FormulaRefAuto{FinitePathTransportBase}[thm]{1}}
  \proofstepwidestar[1,2]{n\in\mathbb N\dsep\FinitePathTransportProperty{V}{E}{n}\vdash\FinitePathTransportProperty{V}{E}{n+1}}{%
      \FormulaRefAuto{FinitePathTransportStep}[thm]{1}}
  \proofstepwidestarbreak[1,2]{\FinitePathTransportProperty{V}{E}{n}}{%
    \FormulaRefAuto{PeanoInductionRuleAddOne}[thm]{%
      6,7,\rAEa{2}}}
  \proofstepwidestarbreak[1,2,3,4,5]{%
    \exists l\in\mathbb N\,\exists s\in\powerset(\mathbb N\times V)\,
    \bigl(\DirectedPath{V}{E}{s}{l}\land s(0)=q(0)\land s(l)=p(n)\bigr)}{%
    \FormulaRefAuto{FinitePathTransportPropertyDef}[def]{%
      8,3,\rAEb{2},4,5}}
\end{tabproofwide}

\FormulaDefDeltaK[Umkehreigenschaft für Pfadlängen]{\begin{aligned}[t]
  &\FinitePathReversalProperty{V}{E}{n}\coloneqq{}\\[-2pt]
  &\quad\forall p\in\powerset(\mathbb N\times V)\,\Bigl(
    \DirectedPath{V}{E}{p}{n}\rightarrow{}\\[-2pt]
  &\qquad\exists q\in\powerset(\mathbb N\times V)\,\bigl(
    \DirectedPath{V}{E}{q}{n}\land{}\\[-2pt]
  &\qquad\quad q(0)=p(n)\land q(n)=p(0)\land{}\\[-2pt]
  &\qquad\quad q[\NatSeg{n+1}]\subseteq p[\NatSeg{n+1}]\bigr)\Bigr)
  \end{aligned}}{FinitePathReversalPropertyDef}{%
  \DeltaRow{Mengen}{V\dsep E\dsep n\dsep p\dsep q}
  \DeltaPrem{Ungerichteter Graph}{\UndirectedGraphStruct{V,E}}
  \DeltaRow{Neue Eigenschaft}{\FinitePathReversalProperty{V}{E}{n}}
}

\FormulaThmDeltaK[Umkehrbasis]{%
  \UndirectedGraphStruct{V,E}\vdash
  \FinitePathReversalProperty{V}{E}{0}%
}{FinitePathReversalBase}{%
  \DeltaRow{Mengen}{V\dsep E\dsep p}
}
\begin{tabproofwide}
  \proofstepwidestarbreak[1]{\UndirectedGraphStruct{V,E}}{\rA}
  \proofstepwidestarbreak[2]{p\in\powerset(\mathbb N\times V)}{\rA}
  \proofstepwidestarbreak[3]{\DirectedPath{V}{E}{p}{0}}{\rA}
  \proofstepwidestar[2,3]{p[\NatSeg{0+1}]\subseteq p[\NatSeg{0+1}]}{%
      \FormulaRefAuto{A\subseteq A}[thm]}
  \proofstepwidestarbreak[2,3]{%
    \exists q\in\powerset(\mathbb N\times V)\,\bigl(
      \DirectedPath{V}{E}{q}{0}\land q(0)=p(0)\land q(0)=p(0)\land
      q[\NatSeg{0+1}]\subseteq p[\NatSeg{0+1}]\bigr)}{%
    \rEI{\rAI{2,3,\rII,\rII,4}}}
  \proofstepwidestarbreak[1]{\FinitePathReversalProperty{V}{E}{0}}{%
    \FormulaRefAuto{FinitePathReversalPropertyDef}[def]{%
      \rUI{\rRI{2,\rRI{3,5}}}}}
\end{tabproofwide}

\FormulaThmDeltaK[Umkehrschritt]{%
  \UndirectedGraphStruct{V,E}\dsep n\in\mathbb N\dsep
  \FinitePathReversalProperty{V}{E}{n}\vdash
  \FinitePathReversalProperty{V}{E}{n+1}%
}{FinitePathReversalStep}{%
  \DeltaRow{Mengen}{V\dsep E\dsep n\dsep p\dsep q\dsep s\dsep t\dsep i}
}
\begin{tabproofwide}
  \proofstepwidestarbreak[1]{\UndirectedGraphStruct{V,E}}{\rA}
  \proofstepwidestarbreak[2]{n\in\mathbb N}{\rA}
  \proofstepwidestarbreak[3]{\FinitePathReversalProperty{V}{E}{n}}{\rA}
  \proofstepwidestarbreak[1]{\GraphStruct{V,E}}{%
    \FormulaRefAuto{UndirectedGraphAccess}[thm]{1}}
  \proofstepwidestarbreak[5]{p\in\powerset(\mathbb N\times V)}{\rA}
  \proofstepwidestarbreak[6]{\DirectedPath{V}{E}{p}{n+1}}{\rA}
  \proofstepwidestarbreak[1,2,5,6]{%
    \exists q\in\powerset(\mathbb N\times V)\,\bigl(
      \DirectedPath{V}{E}{q}{n}\land q(0)=p(1)\land q(n)=p(n+1)\land
      \forall i\in\NatSeg{n+1}\,(q(i)=p(i+1))\bigr)}{%
    \FormulaRefAuto{FinitePathTailSegment}[thm]{4,2,6}}
  \proofstepwidestarbreak[8]{q\in\powerset(\mathbb N\times V)\land
    \DirectedPath{V}{E}{q}{n}\land q(0)=p(1)\land q(n)=p(n+1)\land
    \forall i\in\NatSeg{n+1}\,(q(i)=p(i+1))}{\rA}
  \proofstepwidestarbreak[1,2,3,8]{%
    \exists s\in\powerset(\mathbb N\times V)\,\bigl(
      \DirectedPath{V}{E}{s}{n}\land s(0)=q(n)\land s(n)=q(0)\land
      s[\NatSeg{n+1}]\subseteq q[\NatSeg{n+1}]\bigr)}{%
    \FormulaRefAuto{FinitePathReversalPropertyDef}[def]{3,\rAEa{8},\rAEb{8}}}
  \proofstepwidestarbreak[10]{s\in\powerset(\mathbb N\times V)\land
    \DirectedPath{V}{E}{s}{n}\land s(0)=q(n)\land s(n)=q(0)\land
    s[\NatSeg{n+1}]\subseteq q[\NatSeg{n+1}]}{\rA}
  \proofstepwidestarbreak[4,8]{q\colon\NatSeg{n+1}\to V}{%
    \FormulaRefAuto{DirectedPathFunction}[thm]{4,\rAEb{8}}}
  \proofstepwidestarbreak[4,8]{p\colon\NatSeg{(n+1)+1}\to V}{%
    \FormulaRefAuto{DirectedPathFunction}[thm]{4,6}}
  \proofstepwidestarbreak[2,8]{%
    q[\NatSeg{n+1}]\subseteq p[\NatSeg{(n+1)+1}]}{%
    \FormulaRefAuto{FinitePathTailImageSubset}[thm]{%
      2,12,11,\rAEn{8}}}
  \proofstepwidestarbreak[1,2,8]{p(0)\notin q[\NatSeg{n+1}]}{%
    \FormulaRefAuto{FinitePathTailAvoidsFirstVertex}[thm]{%
      4,2,6,11,\rAEn{8}}}
  \proofstepwidestar[2,8,10]{\forall x\,(x\in s[\NatSeg{n+1}]\rightarrow x\in q[\NatSeg{n+1}])}{%
      \FormulaRefAuto{A\subseteq B\coloneq\forall x\,(x\in A\rightarrow x\in B)}[def]{%
        \rAEn{10}}}
  \proofstepwidestarbreak[2,8,10]{%
    p(0)\notin s[\NatSeg{n+1}]}{%
    \FormulaRefAuto{P\rightarrow Q, \neg Q\vdash\neg P}[thm]{%
      \rUE{15},14}}
  \proofstepwidestar[4,6]{\DirectedWalk V E p { { n + 1 } }}{%
      \FormulaRefAuto{DirectedPathWalkAxiom}[ax]{4,6}}
  \proofstepwidestarbreak[4,6]{p(0),p(1)\in V}{%
    \FormulaRefAuto{DirectedWalkEndpointTyping}[thm]{%
      4,17}}
  \proofstepwidestar[2]{0 \in \mathbb N}{%
      \FormulaRefAuto{PeanoZero}[ax]}
  \proofstepwidestar[2]{0\leq n}{%
      \FormulaRefAuto{PeanoZeroLeq}[thm]{2}}
  \proofstepwidestarbreak[2]{0\in\NatSeg{n+1}}{%
    \FormulaRefAuto{PeanoNatSegLeqAddOneMembership}[thm]{%
      19,2,
      20}}
  \proofstepwidestar[1,2,6]{\DirectedWalk V E p { { n + 1 } }}{%
      \FormulaRefAuto{DirectedPathWalkAxiom}[ax]{4,6}}
  \proofstepwidestarbreak[1,2,6]{(p(0),p(1))\in E}{%
    \FormulaRefAuto{DirectedWalkStepEdge}[thm]{%
      4,22,21}}
  \proofstepwidestarbreak[1,2,6]{(p(1),p(0))\in E}{%
    \FormulaRefAuto{UndirectedAdjacencySymmetric}[thm]{%
      1,\rAEa{18},\rAEb{18},23}}
  \proofstepwidestarbreak[1,2,8,10]{s(n)=p(1)}{%
    \rIE{\rAEa{\rAEn{8}},\rAEn{10}}}
  \proofstepwidestarbreak[1,2,8,10]{s(0)=p(n+1)}{%
    \rIE{\rAEb{\rAEn{8}},\rAEb{10}}}
  \proofstepwidestarbreak[1,2,6,8,10]{%
    \exists t\in\powerset(\mathbb N\times V)\,\bigl(
      \DirectedPath{V}{E}{t}{n+1}\land t(0)=s(0)\land t(n+1)=p(0)\land
      t[\NatSeg{(n+1)+1}]=s[\NatSeg{n+1}]\cup\{p(0)\}\bigr)}{%
    \FormulaRefAuto{FinitePathFreshEndpointExtension}[thm]{%
      4,\rAEb{10},\rAEa{18},16,\rIE{25,24}}}
  \proofstepwidestarbreak[28]{t\in\powerset(\mathbb N\times V)\land
    \DirectedPath{V}{E}{t}{n+1}\land t(0)=s(0)\land t(n+1)=p(0)\land
    t[\NatSeg{(n+1)+1}]=s[\NatSeg{n+1}]\cup\{p(0)\}}{\rA}
  \proofstepwidestar[2,5,6,8,10,28]{\NatSeg{n+1}\subseteq\NatSeg{n+1}}{%
      \FormulaRefAuto{A\subseteq A}[thm]}
  \proofstepwidestar[2,5,6,8,10,28]{p ( 0 ) \in p [ { \NatSeg { n + 1 } } ]}{%
      \FormulaRefAuto{C\subseteq A\dsep F\colon A\to B\dsep x\in C\vdash F(x)\in F[C]}[thm]{%
              29,12,21}}
  \proofstepwidestar[2,5,6,8,10,28]{\{ { p ( 0 ) } \} \subseteq { p [ \NatSeg { n + 1 } ] }}{%
      \FormulaRefAuto{a\in A\vdash\{a\}\subseteq A}[thm]{%
            30}}
  \proofstepwidestar[2,5,6,8,10,28]{s[\NatSeg{n+1}]\subseteq p[\NatSeg{(n+1)+1}]}{%
      \FormulaRefAuto{A\subseteq B\dsep B\subseteq C\vdash A\subseteq C}[thm]{%
            \rAEn{10},13}}
  \proofstepwidestar[2,5,6,8,10,28]{s[\NatSeg{n+1}]\cup\{p(0)\}\subseteq p[\NatSeg{(n+1)+1}]}{%
      \FormulaRefAuto{A\subseteq C\dsep B\subseteq C\vdash
        A\cup B\subseteq C}[thm]{%
          32,
          31}}
  \proofstepwidestarbreak[2,5,6,8,10,28]{%
    t[\NatSeg{(n+1)+1}]\subseteq p[\NatSeg{(n+1)+1}]}{%
    \rIE{\rAEn{28},
      33}}
  \proofstepwidestarbreak[1,2,5,6,8,10,28]{%
    \exists t\in\powerset(\mathbb N\times V)\,\bigl(
      \DirectedPath{V}{E}{t}{n+1}\land t(0)=p(n+1)\land t(n+1)=p(0)\land
      t[\NatSeg{(n+1)+1}]\subseteq p[\NatSeg{(n+1)+1}]\bigr)}{%
    \rEI{\rAI{\rAEa{28},\rAEb{28},\rIE{26,\rAEn{28}},
      \rAEn{28},34}}}
  \proofstepwidestarbreak[1,2,5,6]{%
    \exists t\in\powerset(\mathbb N\times V)\,\bigl(
      \DirectedPath{V}{E}{t}{n+1}\land t(0)=p(n+1)\land t(n+1)=p(0)\land
      t[\NatSeg{(n+1)+1}]\subseteq p[\NatSeg{(n+1)+1}]\bigr)}{%
    \rEE{7,8,\rEE{9,10,\rEE{27,28,35}}}}
  \proofstepwidestarbreak[1,2,3]{\FinitePathReversalProperty{V}{E}{n+1}}{%
    \FormulaRefAuto{FinitePathReversalPropertyDef}[def]{%
      \rUI{\rRI{5,\rRI{6,36}}}}}
\end{tabproofwide}

\FormulaThmDeltaK[Umkehrung eines Pfads im ungerichteten Graphen]{%
  \begin{aligned}[t]
  &\UndirectedGraphStruct{V,E}\dsep\DirectedPath{V}{E}{p}{n}\vdash{}\\[-2pt]
  &\quad\exists q\in\powerset(\mathbb N\times V)\,\bigl(
    \DirectedPath{V}{E}{q}{n}\land q(0)=p(n)\land q(n)=p(0)\bigr)
  \end{aligned}%
}{UndirectedFinitePathReversal}{%
  \DeltaRow{Mengen}{V\dsep E\dsep n\dsep p\dsep q}
}
\begin{tabproofwide}
  \proofstepwidestarinline[1]{\UndirectedGraphStruct{V,E}}{\rA}
  \proofstepwidestarinline[2]{\DirectedPath{V}{E}{p}{n}}{\rA}
  \proofstepwidestarinline[1]{\GraphStruct{V,E}}{%
    \FormulaRefAuto{UndirectedGraphAccess}[thm]{1}}
  \proofstepwidestar[1,2]{\DirectedWalk V E p n}{%
      \FormulaRefAuto{DirectedPathWalkAxiom}[ax]{3,2}}
  \proofstepwidestar[1,2]{\begin{aligned}[t]&n\in\mathbb N\land p\colon\NatSeg{n+1}\to V\land{}\\[-2pt]&\forall i\in\NatSeg{n}\,((p(i),p(i+1))\in E)\end{aligned}}{%
      \FormulaRefAuto{DirectedWalkCharacterization}[thm]{%
      3,4}}
  \proofstepwidestarbreakkeep[1,2]{n\in\mathbb N}{%
    \rAEa{5}}
  \proofstepwidestarinline[1]{\FinitePathReversalProperty{V}{E}{0}}{%
    \FormulaRefAuto{FinitePathReversalBase}[thm]{1}}
  \proofstepwidestar[1,2]{n\in\mathbb N\dsep\FinitePathReversalProperty{V}{E}{n}\vdash\FinitePathReversalProperty{V}{E}{n+1}}{%
      \FormulaRefAuto{FinitePathReversalStep}[thm]{1}}
  \proofstepwidestarbreakkeep[1,2]{\FinitePathReversalProperty{V}{E}{n}}{%
    \FormulaRefAuto{PeanoInductionRuleAddOne}[thm]{%
      7,8,6}}
  \proofstepwidestarinline[1,2]{p\in\powerset(\mathbb N\times V)}{%
    \FormulaRefAuto{FinitePathPowerSetTyping}[thm]{3,2}}
  \proofstepwidestarbreakkeep[1,2]{%
    \exists q\in\powerset(\mathbb N\times V)\,\bigl(
      \DirectedPath{V}{E}{q}{n}\land q(0)=p(n)\land q(n)=p(0)\land
      q[\NatSeg{n+1}]\subseteq p[\NatSeg{n+1}]\bigr)}{%
    \FormulaRefAuto{FinitePathReversalPropertyDef}[def]{9,10,2}}
  \proofstepwidestarbreakkeep[12]{%
    q\in\powerset(\mathbb N\times V)\land\Bigl(
      \bigl(\DirectedPath{V}{E}{q}{n}\land q(0)=p(n)\land q(n)=p(0)\bigr)\land
      q[\NatSeg{n+1}]\subseteq p[\NatSeg{n+1}]\Bigr)}{\rA}
  \proofstepwidestarbreakkeep[12]{%
    \exists q\in\powerset(\mathbb N\times V)\,\bigl(
      \DirectedPath{V}{E}{q}{n}\land q(0)=p(n)\land q(n)=p(0)\bigr)}{%
    \rEI{\rAI{\rAEa{12},\rAEa{\rAEb{12}}}}}
  \proofstepwidestarbreakkeep[1,2]{%
    \exists q\in\powerset(\mathbb N\times V)\,\bigl(
      \DirectedPath{V}{E}{q}{n}\land q(0)=p(n)\land q(n)=p(0)\bigr)}{%
    \rEE{11,12,13}}
\end{tabproofwide}

\FormulaDefDeltaK[Schwanzeigenschaft an einer Wegstelle]{%
  \begin{aligned}[t]
  &\FinitePathSuffixProperty{V}{E}{p}{n}{k}\coloneqq{}\\[-2pt]
  &\quad\forall l\in\mathbb N\,\Bigl(k+l=n\rightarrow{}\\[-2pt]
  &\qquad\exists q\in\powerset(\mathbb N\times V)\,\bigl(
    \DirectedPath{V}{E}{q}{l}\land q(0)=p(k)\land q(l)=p(n)\land{}\\[-2pt]
  &\qquad\quad\forall i\in\NatSeg{l+1}\,(q(i)=p(k+i))\bigr)\Bigr)
  \end{aligned}%
}{FinitePathSuffixPropertyDef}{%
  \DeltaRow{Mengen}{V\dsep E\dsep p\dsep q\dsep n\dsep k\dsep l\dsep i}
  \DeltaRow{Neue Eigenschaft}{\FinitePathSuffixProperty{V}{E}{p}{n}{k}}
}

\FormulaThmDeltaK[Schwanzbasis]{%
  \GraphStruct{V,E}\dsep n\in\mathbb N\dsep
  p\in\powerset(\mathbb N\times V)\dsep
  \DirectedPath{V}{E}{p}{n}\vdash
  \FinitePathSuffixProperty{V}{E}{p}{n}{0}%
}{FinitePathSuffixBase}{%
  \DeltaRow{Mengen}{V\dsep E\dsep p\dsep q\dsep n\dsep l\dsep i}
}
\begin{tabproofwide}
  \proofstepwidestarbreak[1]{\GraphStruct{V,E}}{\rA}
  \proofstepwidestarbreak[2]{n\in\mathbb N}{\rA}
  \proofstepwidestarbreak[3]{p\in\powerset(\mathbb N\times V)}{\rA}
  \proofstepwidestarbreak[4]{\DirectedPath{V}{E}{p}{n}}{\rA}
  \proofstepwidestarbreak[5]{l\in\mathbb N}{\rA}
  \proofstepwidestarbreak[6]{0+l=n}{\rA}
  \proofstepwidestarbreak[5]{0+l=l}{%
    \FormulaRefAuto{PeanoAddLeftZero}[thm]{5}}
  \proofstepwidestarbreak[5,6]{l=n}{\rIE{7,6}}
  \proofstepwidestarbreak[4,5,6]{\DirectedPath{V}{E}{p}{l}}{\rIE{8,4}}
  \proofstepwidestarbreak[5,6]{p(l)=p(n)}{\rIE{8,\rII}}
  \proofstepwidestarbreak[11]{i\in\NatSeg{l+1}}{\rA}
  \proofstepwidestar[5,11]{1 \in \mathbb N}{%
      \FormulaRefAuto{PeanoOneInN}[thm]}
  \proofstepwidestar[5,11]{l + 1 \in \mathbb N}{%
      \FormulaRefAuto{PeanoAddClosure}[thm]{5,
          12}}
  \proofstepwidestar[5,11]{\NatSeg { ( l + 1 ) } \subseteq \mathbb N}{%
      \FormulaRefAuto{PeanoNatSegSubsetN}[thm]{%
        13}}
  \proofstepwidestarbreak[5,11]{i\in\mathbb N}{%
    \FormulaRefAuto{A\subseteq B\dsep x\in A\vdash x\in B}[thm]{%
      14,11}}
  \proofstepwidestarbreak[5,11]{0+i=i}{%
    \FormulaRefAuto{PeanoAddLeftZero}[thm]{15}}
  \proofstepwidestar[5,11]{i = { ( 0 + i ) }}{%
      \FormulaRefAuto{a=b\vdash b=a}[thm]{16}}
  \proofstepwidestarbreak[5,11]{p(i)=p(0+i)}{%
    \rIE{17,\rII}}
  \proofstepwidestarbreak[5]{%
    \forall i\in\NatSeg{l+1}\,(p(i)=p(0+i))}{%
    \rUI{\rRI{11,18}}}
  \proofstepwidestarbreak[2,3,4,5,6]{%
    \exists q\in\powerset(\mathbb N\times V)\,\bigl(
      \DirectedPath{V}{E}{q}{l}\land q(0)=p(0)\land q(l)=p(n)\land
      \forall i\in\NatSeg{l+1}\,(q(i)=p(0+i))\bigr)}{%
    \rEI{\rAI{3,9,\rII,10,19}}}
  \proofstepwidestarbreak[1,2,3,4]{%
    \FinitePathSuffixProperty{V}{E}{p}{n}{0}}{%
    \FormulaRefAuto{FinitePathSuffixPropertyDef}[def]{%
      \rUI{\rRI{5,\rRI{6,20}}}}}
\end{tabproofwide}

\FormulaThmDeltaK[Schwanzschritt]{%
  \begin{aligned}[t]
  &\GraphStruct{V,E}\dsep n,k\in\mathbb N\dsep
    p\in\powerset(\mathbb N\times V)\dsep
    \DirectedPath{V}{E}{p}{n}\dsep
    \FinitePathSuffixProperty{V}{E}{p}{n}{k}\vdash{}\\[-2pt]
  &\qquad\FinitePathSuffixProperty{V}{E}{p}{n}{k+1}
  \end{aligned}%
}{FinitePathSuffixStep}{%
  \DeltaRow{Mengen}{V\dsep E\dsep p\dsep q\dsep s\dsep n\dsep k\dsep l\dsep i}
}
\begin{tabproofwide}
  \proofstepwidestarbreak[1]{\GraphStruct{V,E}}{\rA}
  \proofstepwidestarbreak[2]{n\in\mathbb N}{\rA}
  \proofstepwidestarbreak[3]{k\in\mathbb N}{\rA}
  \proofstepwidestarbreak[4]{p\in\powerset(\mathbb N\times V)}{\rA}
  \proofstepwidestarbreak[5]{\DirectedPath{V}{E}{p}{n}}{\rA}
  \proofstepwidestarbreak[6]{\FinitePathSuffixProperty{V}{E}{p}{n}{k}}{\rA}
  \proofstepwidestarbreak[7]{l\in\mathbb N}{\rA}
  \proofstepwidestarbreak[8]{(k+1)+l=n}{\rA}
  \proofstepwidestar[3,7]{1 \in \mathbb N}{%
      \FormulaRefAuto{PeanoOneInN}[thm]}
  \proofstepwidestar[3,7]{1 \in \mathbb N}{%
      \FormulaRefAuto{PeanoOneInN}[thm]}
  \proofstepwidestar[3,7]{l + 1 = 1 + l}{%
      \FormulaRefAuto{PeanoAddCommutative}[thm]{%
          7,9}}
  \proofstepwidestar[3,7]{( k + 1 ) + l = k + ( 1 + l )}{%
      \FormulaRefAuto{PeanoAddAssociative}[thm]{%
          3,10,7}}
  \proofstepwidestar[3,7]{{ ( k + ( 1 + l ) ) } = { ( ( k + 1 ) + l ) }}{%
      \FormulaRefAuto{a=b\vdash b=a}[thm]{%
        12}}
  \proofstepwidestarbreak[3,7]{k+(l+1)=(k+1)+l}{%
    \rChain{%
      \rIE{%
        11,\rII},
      13}}
  \proofstepwidestarbreak[3,7,8]{k+(l+1)=n}{\rIE{14,8}}
  \proofstepwidestar[7]{1 \in \mathbb N}{%
      \FormulaRefAuto{PeanoOneInN}[thm]}
  \proofstepwidestarbreak[7]{l+1\in\mathbb N}{%
    \FormulaRefAuto{PeanoAddClosure}[thm]{7,
      16}}
  \proofstepwidestarbreak[3,6,7,8]{%
    \exists q\in\powerset(\mathbb N\times V)\,\bigl(
      \DirectedPath{V}{E}{q}{l+1}\land q(0)=p(k)\land q(l+1)=p(n)\land
      \forall i\in\NatSeg{(l+1)+1}\,(q(i)=p(k+i))\bigr)}{%
    \FormulaRefAuto{FinitePathSuffixPropertyDef}[def]{6,17,15}}
  \proofstepwidestarbreak[19]{q\in\powerset(\mathbb N\times V)\land
      \DirectedPath{V}{E}{q}{l+1}\land q(0)=p(k)\land q(l+1)=p(n)\land
      \forall i\in\NatSeg{(l+1)+1}\,(q(i)=p(k+i))}{\rA}
  \proofstepwidestarbreak[1,7,19]{%
    \exists s\in\powerset(\mathbb N\times V)\,\bigl(
      \DirectedPath{V}{E}{s}{l}\land s(0)=q(1)\land s(l)=q(l+1)\land
      \forall i\in\NatSeg{l+1}\,(s(i)=q(i+1))\bigr)}{%
    \FormulaRefAuto{FinitePathTailSegment}[thm]{1,7,\rAEb{19}}}
  \proofstepwidestarbreak[21]{s\in\powerset(\mathbb N\times V)\land
      \DirectedPath{V}{E}{s}{l}\land s(0)=q(1)\land s(l)=q(l+1)\land
      \forall i\in\NatSeg{l+1}\,(s(i)=q(i+1))}{\rA}
  \proofstepwidestar[7]{1 \in \mathbb N}{%
      \FormulaRefAuto{PeanoOneInN}[thm]}
  \proofstepwidestar[7]{1 \in \mathbb N}{%
      \FormulaRefAuto{PeanoOneInN}[thm]}
  \proofstepwidestar[7]{1 \in \mathbb N}{%
      \FormulaRefAuto{PeanoOneInN}[thm]}
  \proofstepwidestar[7]{1 \in \mathbb N}{%
      \FormulaRefAuto{PeanoOneInN}[thm]}
  \proofstepwidestar[7]{1 \in \mathbb N}{%
      \FormulaRefAuto{PeanoOneInN}[thm]}
  \proofstepwidestar[7]{l + 1 \in \mathbb N}{%
      \FormulaRefAuto{PeanoAddClosure}[thm]{7,
          23}}
  \proofstepwidestar[7]{l + 1 \in \mathbb N}{%
      \FormulaRefAuto{PeanoAddClosure}[thm]{7,
            25}}
  \proofstepwidestar[7]{1 + l = l + 1}{%
      \FormulaRefAuto{PeanoAddCommutative}[thm]{%
            26,7}}
  \proofstepwidestar[7]{1\leq (l+1)\leftrightarrow 1\in\NatSeg{(l+1)+1}}{%
      \FormulaRefAuto{PeanoNatSegLeqAddOneMembership}[thm]{%
        22,
        27}}
  \proofstepwidestar[7]{1\leq l+1\leftrightarrow\exists k\in\mathbb N\,(1+k=l+1)}{%
      \FormulaRefAuto{m\in\mathbb{N}\dsep n\in\mathbb{N}\vdash
          m\leq n\leftrightarrow\exists k\in\mathbb{N}\,(m+k=n)}[thm]{%
          24,
          28}}
  \proofstepwidestar[7]{1\leq l+1}{%
      \FormulaRefAuto{P\leftrightarrow Q\dsep Q\vdash P}[thm]{%
        31,
        \rEI{\rAI{7,
          29}}}}
  \proofstepwidestarbreak[7]{1\in\NatSeg{(l+1)+1}}{%
    \FormulaRefAuto{P\leftrightarrow Q\dsep P\vdash Q}[thm]{%
      30,
      32}}
  \proofstepwidestarbreak[7,19]{q(1)=p(k+1)}{\rUE{\rAEn{19},33}}
  \proofstepwidestarbreak[7,19,21]{s(0)=p(k+1)}{\rIE{34,\rAEa{\rAEn{21}}}}
  \proofstepwidestarbreak[8,19,21]{s(l)=p(n)}{%
    \rIE{\rAEb{\rAEn{21}},\rAEn{19}}}
  \proofstepwidestarbreak[37]{i\in\NatSeg{l+1}}{\rA}
  \proofstepwidestar[7,21]{1 \in \mathbb N}{%
      \FormulaRefAuto{PeanoOneInN}[thm]}
  \proofstepwidestar[7,36]{l+1\in\mathbb N}{%
      \FormulaRefAuto{PeanoAddClosure}[thm]{7,
        38}}
  \proofstepwidestarbreak[7,37]{i+1\in\NatSeg{(l+1)+1}}{\FormulaRefAuto{FinitePathIndexShiftOne}[thm]{39,37}}
  \proofstepwidestarbreak[19,37]{q(i+1)=p(k+(i+1))}{\rUE{\rAEn{19},40}}
  \proofstepwidestar[3,7,37]{1 \in \mathbb N}{%
      \FormulaRefAuto{PeanoOneInN}[thm]}
  \proofstepwidestar[3,7,37]{1 \in \mathbb N}{%
      \FormulaRefAuto{PeanoOneInN}[thm]}
  \proofstepwidestar[3,7,37]{1 \in \mathbb N}{%
      \FormulaRefAuto{PeanoOneInN}[thm]}
  \proofstepwidestar[3,7,37]{1 \in \mathbb N}{%
      \FormulaRefAuto{PeanoOneInN}[thm]}
  \proofstepwidestar[3,7,37]{1 \in \mathbb N}{%
      \FormulaRefAuto{PeanoOneInN}[thm]}
  \proofstepwidestar[3,7,37]{1 \in \mathbb N}{%
      \FormulaRefAuto{PeanoOneInN}[thm]}
  \proofstepwidestar[3,7,37]{l + 1 \in \mathbb N}{%
      \FormulaRefAuto{PeanoAddClosure}[thm]{7,
              42}}
  \proofstepwidestar[3,7,37]{l + 1 \in \mathbb N}{%
      \FormulaRefAuto{PeanoAddClosure}[thm]{7,
              44}}
  \proofstepwidestar[3,7,37]{l + 1 \in \mathbb N}{%
      \FormulaRefAuto{PeanoAddClosure}[thm]{7,
                47}}
  \proofstepwidestar[3,7,37]{\NatSeg { ( l + 1 ) } \subseteq \mathbb N}{%
      \FormulaRefAuto{PeanoNatSegSubsetN}[thm]{%
            48}}
  \proofstepwidestar[3,7,37]{\NatSeg { ( l + 1 ) } \subseteq \mathbb N}{%
      \FormulaRefAuto{PeanoNatSegSubsetN}[thm]{%
            49}}
  \proofstepwidestar[3,7,37]{\NatSeg { ( l + 1 ) } \subseteq \mathbb N}{%
      \FormulaRefAuto{PeanoNatSegSubsetN}[thm]{%
              50}}
  \proofstepwidestar[3,7,37]{i\in\mathbb N}{\FormulaRefAuto{A\subseteq B\dsep x\in A\vdash x\in B}[thm]{51,37}}
  \proofstepwidestar[3,7,37]{i\in\mathbb N}{\FormulaRefAuto{A\subseteq B\dsep x\in A\vdash x\in B}[thm]{52,37}}
  \proofstepwidestar[3,7,37]{i\in\mathbb N}{\FormulaRefAuto{A\subseteq B\dsep x\in A\vdash x\in B}[thm]{53,37}}
  \proofstepwidestar[3,7,37]{i+1=1+i}{%
      \FormulaRefAuto{PeanoAddCommutative}[thm]{%
        55,
        45}}
  \proofstepwidestar[3,7,37]{(k+1)+i=k+(1+i)}{%
      \FormulaRefAuto{PeanoAddAssociative}[thm]{%
          3,46,
          56}}
  \proofstepwidestar[3,7,37]{(k+i)+1=k+(i+1)}{%
      \FormulaRefAuto{PeanoAddAssociative}[thm]{%
        3,
        54,
        43}}
  \proofstepwidestar[3,7,37]{k+(1+i)=(k+1)+i}{%
      \FormulaRefAuto{a=b\vdash b=a}[thm]{%
        58}}
  \proofstepwidestarbreak[3,7,37]{k+(i+1)=(k+1)+i}{\rChain{\rIE{57,\rII},60}}
  \proofstepwidestarbreak[19,21,37]{s(i)=q(i+1)}{\rUE{\rAEn{21},37}}
  \proofstepwidestarbreak[3,7,19,21,37]{s(i)=p((k+1)+i)}{\rIE{61,\rIE{41,62}}}
  \proofstepwidestarbreak[3,7,19,21]{%
    \forall i\in\NatSeg{l+1}\,(s(i)=p((k+1)+i))}{\rUI{\rRI{37,63}}}
  \proofstepwidestarbreak[3,7,8,19,21]{%
    \exists s\in\powerset(\mathbb N\times V)\,\bigl(
      \DirectedPath{V}{E}{s}{l}\land s(0)=p(k+1)\land s(l)=p(n)\land
      \forall i\in\NatSeg{l+1}\,(s(i)=p((k+1)+i))\bigr)}{\rEI{\rAI{\rAEa{21},\rAEb{21},35,36,64}}}
  \proofstepwidestarbreak[1,3,6,7,8]{%
    \exists s\in\powerset(\mathbb N\times V)\,\bigl(
      \DirectedPath{V}{E}{s}{l}\land s(0)=p(k+1)\land s(l)=p(n)\land
      \forall i\in\NatSeg{l+1}\,(s(i)=p((k+1)+i))\bigr)}{\rEE{18,19,\rEE{20,21,65}}}
  \proofstepwidestarbreak[1,2,3,4,5,6]{%
    \FinitePathSuffixProperty{V}{E}{p}{n}{k+1}}{\FormulaRefAuto{FinitePathSuffixPropertyDef}[def]{\rUI{\rRI{7,\rRI{8,66}}}}}
\end{tabproofwide}

\FormulaThmDeltaK[Endabschnitt eines endlichen Pfads]{\begin{aligned}[t]
  &\GraphStruct{V,E}\dsep n\in\mathbb N\dsep
    p\in\powerset(\mathbb N\times V)\dsep{}\\[-2pt]
  &\DirectedPath{V}{E}{p}{n}\dsep k\in\NatSeg{n+1}\vdash{}\\[-2pt]
  &\quad\exists l\in\mathbb N\,\exists q\in\powerset(\mathbb N\times V)\,
    \bigl(k+l=n\land{}\\[-2pt]
  &\qquad\DirectedPath{V}{E}{q}{l}\land q(0)=p(k)\land q(l)=p(n)\land{}\\[-2pt]
  &\qquad\forall i\in\NatSeg{l+1}\,(q(i)=p(k+i))\bigr)
  \end{aligned}}{FinitePathFinalSegment}{%
  \DeltaRow{Mengen}{V\dsep E\dsep p\dsep q\dsep n\dsep k\dsep l\dsep i}
}
\begin{tabproofwide}
  \proofstepwidestarbreak[1]{\GraphStruct{V,E}}{\rA}
  \proofstepwidestarbreak[2]{n\in\mathbb N}{\rA}
  \proofstepwidestarbreak[3]{p\in\powerset(\mathbb N\times V)}{\rA}
  \proofstepwidestarbreak[4]{\DirectedPath{V}{E}{p}{n}}{\rA}
  \proofstepwidestarbreak[5]{k\in\NatSeg{n+1}}{\rA}
  \proofstepwidestar[2,5]{k\in\mathbb N\land k\leq n\land\NatSeg{k+1}\subseteq\NatSeg{n+1}}{%
      \FormulaRefAuto{FinitePathPrefixIndexData}[thm]{2,5}}
  \proofstepwidestarbreak[2,5]{k\in\mathbb N}{%
    \rAEa{6}}
  \proofstepwidestar[2,5]{k\in\mathbb N\land k\leq n\land\NatSeg{k+1}\subseteq\NatSeg{n+1}}{%
      \FormulaRefAuto{FinitePathPrefixIndexData}[thm]{2,5}}
  \proofstepwidestarbreak[2,5]{k\leq n}{%
    \rAEb{8}}
  \proofstepwidestar[2,5]{k \leq n \leftrightarrow \exists l \in \mathbb N ( k + l = n )}{%
      \FormulaRefAuto{m\in\mathbb{N}\dsep n\in\mathbb{N}\vdash
        m\leq n\leftrightarrow\exists k\in\mathbb{N}\,(m+k=n)}[thm]{%
        7,2}}
  \proofstepwidestarbreak[2,5]{\exists l\in\mathbb N\,(k+l=n)}{%
    \FormulaRefAuto{P\leftrightarrow Q\dsep P\vdash Q}[thm]{%
      10,9}}
  \proofstepwidestar[1,2,3,4,5]{\FinitePathSuffixProperty V E p n 0}{%
      \FormulaRefAuto{FinitePathSuffixBase}[thm]{1,2,3,4}}
  \proofstepwidestar[1,2,3,4,5]{\begin{aligned}[t]&k\in\mathbb N\dsep\FinitePathSuffixProperty{V}{E}{p}{n}{k}\\[-2pt]&\quad\vdash\FinitePathSuffixProperty{V}{E}{p}{n}{k+1}\end{aligned}}{%
      \FormulaRefAuto{FinitePathSuffixStep}[thm]{1,2,3,4}}
  \proofstepwidestarbreak[1,2,3,4,5]{%
    \FinitePathSuffixProperty{V}{E}{p}{n}{k}}{%
    \FormulaRefAuto{PeanoInductionRuleAddOne}[thm]{%
      12,
      13,7}}
  \proofstepwidestarbreak[15]{l\in\mathbb N\land k+l=n}{\rA}
  \proofstepwidestarbreak[1,2,3,4,5,14,15]{%
    \exists q\in\powerset(\mathbb N\times V)\,\bigl(
      \DirectedPath{V}{E}{q}{l}\land q(0)=p(k)\land q(l)=p(n)\land
      \forall i\in\NatSeg{l+1}\,(q(i)=p(k+i))\bigr)}{%
    \FormulaRefAuto{FinitePathSuffixPropertyDef}[def]{14,\rAEa{15},\rAEb{15}}}
  \proofstepwidestarbreak[17]{q\in\powerset(\mathbb N\times V)\land
      \DirectedPath{V}{E}{q}{l}\land q(0)=p(k)\land q(l)=p(n)\land
      \forall i\in\NatSeg{l+1}\,(q(i)=p(k+i))}{\rA}
  \proofstepwidestarbreak[15,17]{%
    \exists l\in\mathbb N\,\exists q\in\powerset(\mathbb N\times V)\,
      \bigl(k+l=n\land\DirectedPath{V}{E}{q}{l}\land
        q(0)=p(k)\land q(l)=p(n)\land
        \forall i\in\NatSeg{l+1}\,(q(i)=p(k+i))\bigr)}{%
    \rEI{\rAI{\rAEa{15},\rEI{\rAI{\rAEa{17},\rAEb{15},
      \rAEb{17},\rAEn{17},\rAEn{17},\rAEn{17}}}}}}
  \proofstepwidestarbreak[1,2,3,4,5,14,15]{%
    \exists l\in\mathbb N\,\exists q\in\powerset(\mathbb N\times V)\,
      \bigl(k+l=n\land\DirectedPath{V}{E}{q}{l}\land
        q(0)=p(k)\land q(l)=p(n)\land
        \forall i\in\NatSeg{l+1}\,(q(i)=p(k+i))\bigr)}{%
    \rEE{16,17,18}}
  \proofstepwidestarbreak[1,2,3,4,5,14]{%
    \exists l\in\mathbb N\,\exists q\in\powerset(\mathbb N\times V)\,
      \bigl(k+l=n\land\DirectedPath{V}{E}{q}{l}\land
        q(0)=p(k)\land q(l)=p(n)\land
        \forall i\in\NatSeg{l+1}\,(q(i)=p(k+i))\bigr)}{%
    \rEE{11,15,19}}
\end{tabproofwide}

\FormulaThmDeltaK[Linkskürzung einer nichtstrikten Summenungleichung]{%
  a,b,c\in\mathbb N\dsep a+b\leq a+c\vdash b\leq c%
}{PeanoAddLeftLeqCancellation}{%
  \DeltaRow{Mengen}{a\dsep b\dsep c\dsep d}
}
\begin{tabproofwide}
  \proofstepwidestarbreak[1]{a\in\mathbb N}{\rA}
  \proofstepwidestarbreak[2]{b\in\mathbb N}{\rA}
  \proofstepwidestarbreak[3]{c\in\mathbb N}{\rA}
  \proofstepwidestarbreak[4]{a+b\leq a+c}{\rA}
  \proofstepwidestar[1,2,3,4]{a + b \in \mathbb N}{%
      \FormulaRefAuto{PeanoAddClosure}[thm]{1,2}}
  \proofstepwidestar[1,2,3,4]{a + c \in \mathbb N}{%
      \FormulaRefAuto{PeanoAddClosure}[thm]{1,3}}
  \proofstepwidestar[1,2,3,4]{a+b\leq a+c\leftrightarrow\exists d\in\mathbb N\,((a+b)+d=a+c)}{%
      \FormulaRefAuto{m\in\mathbb{N}\dsep n\in\mathbb{N}\vdash m\leq n\leftrightarrow\exists k\in\mathbb{N}\,(m+k=n)}[thm]{%
        5,
        6}}
  \proofstepwidestarbreak[1,2,3,4]{%
    \exists d\in\mathbb N\,((a+b)+d=a+c)}{%
    \FormulaRefAuto{P\leftrightarrow Q\dsep P\vdash Q}[thm]{%
      7,4}}
  \proofstepwidestarbreak[9]{d\in\mathbb N\land(a+b)+d=a+c}{\rA}
  \proofstepwidestar[1,2,9]{( a + b ) + d = a + ( b + d )}{%
      \FormulaRefAuto{PeanoAddAssociative}[thm]{%
      1,2,\rAEa{9}}}
  \proofstepwidestarbreak[1,2,9]{a+(b+d)=a+c}{%
    \rIE{10,\rAEb{9}}}
  \proofstepwidestar[1,2,3,9]{b + d \in \mathbb N}{%
      \FormulaRefAuto{PeanoAddClosure}[thm]{2,\rAEa{9}}}
  \proofstepwidestarbreak[1,2,3,9]{b+d=c}{%
    \FormulaRefAuto{PeanoAddLeftCancellation}[thm]{%
      1,12,3,11}}
  \proofstepwidestar[1,2,3,9]{b\leq c\leftrightarrow\exists d\in\mathbb N\,(b+d=c)}{%
      \FormulaRefAuto{m\in\mathbb{N}\dsep n\in\mathbb{N}\vdash m\leq n\leftrightarrow\exists k\in\mathbb{N}\,(m+k=n)}[thm]{2,3}}
  \proofstepwidestarbreak[1,2,3,9]{b\leq c}{%
    \FormulaRefAuto{P\leftrightarrow Q\dsep Q\vdash P}[thm]{%
      14,
      \rEI{\rAI{\rAEa{9},13}}}}
  \proofstepwidestarbreak[1,2,3,4]{b\leq c}{\rEE{8,9,15}}
\end{tabproofwide}

\FormulaThmDeltaK[Nichtstrikte Ordnung ohne Gleichheit ist strikt]{%
  m,n\in\mathbb N\dsep m\leq n\dsep m\neq n\vdash m<n%
}{PeanoLeqNeqStrict}{%
  \DeltaRow{Mengen}{m\dsep n}
}
\begin{tabproofwide}
  \proofstepwidestarbreak[1]{m\in\mathbb N}{\rA}
  \proofstepwidestarbreak[2]{n\in\mathbb N}{\rA}
  \proofstepwidestarbreak[3]{m\leq n}{\rA}
  \proofstepwidestarbreak[4]{m\neq n}{\rA}
  \proofstepwidestarbreak[1,2,3]{m<n\lor m=n}{%
    \FormulaRefAuto{PeanoLeqSplit}[thm]{1,2,3}}
  \proofstepwidestarbreak[6]{m<n}{\rA}
  \proofstepwidestarbreak[7]{m=n}{\rA}
  \proofstepwidestarbreak[4,7]{m<n}{%
    \FormulaRefAuto{P,\neg P\vdash Q}[thm]{7,4}}
  \proofstepwidestarbreak[1,2,3,4]{m<n}{\rOE{5,6,6,7,8}}
\end{tabproofwide}

\FormulaThmDeltaK[Strikte Monotonie der Summe in beiden Summanden]{%
  a,b,c,d\in\mathbb N\dsep a<b\dsep c<d\vdash a+c<b+d%
}{PeanoAddStrictMonotoneBoth}{%
  \DeltaRow{Mengen}{a\dsep b\dsep c\dsep d}
}
\begin{tabproofwide}
  \proofstepwidestarbreak[1]{a\in\mathbb N}{\rA}
  \proofstepwidestarbreak[2]{b\in\mathbb N}{\rA}
  \proofstepwidestarbreak[3]{c\in\mathbb N}{\rA}
  \proofstepwidestarbreak[4]{d\in\mathbb N}{\rA}
  \proofstepwidestarbreak[5]{a<b}{\rA}
  \proofstepwidestarbreak[6]{c<d}{\rA}
  \proofstepwidestarbreak[1,2,3,5]{c+a<c+b}{%
    \FormulaRefAuto{PeanoAddLeftStrictMonotone}[thm]{3,1,2,5}}
  \proofstepwidestar[1,2,3,5]{c + a = a + c}{%
      \FormulaRefAuto{PeanoAddCommutative}[thm]{3,1}}
  \proofstepwidestar[1,2,3,5]{c + b = b + c}{%
      \FormulaRefAuto{PeanoAddCommutative}[thm]{3,2}}
  \proofstepwidestarbreak[1,2,3,5]{a+c<b+c}{%
    \rIE{8,
      \rIE{9,7}}}
  \proofstepwidestarbreak[2,3,4,6]{b+c<b+d}{%
    \FormulaRefAuto{PeanoAddLeftStrictMonotone}[thm]{2,3,4,6}}
  \proofstepwidestar[1,2,3,4,5,6]{a + c \in \mathbb N}{%
      \FormulaRefAuto{PeanoAddClosure}[thm]{1,3}}
  \proofstepwidestar[1,2,3,4,5,6]{b + c \in \mathbb N}{%
      \FormulaRefAuto{PeanoAddClosure}[thm]{2,3}}
  \proofstepwidestar[1,2,3,4,5,6]{b + d \in \mathbb N}{%
      \FormulaRefAuto{PeanoAddClosure}[thm]{2,4}}
  \proofstepwidestarbreak[1,2,3,4,5,6]{a+c<b+d}{%
    \FormulaRefAuto{PeanoLtTransitive}[thm]{%
      12,
      13,
      14,10,11}}
\end{tabproofwide}

\FormulaThmDeltaK[Ein endlicher Weg ist durch zwei anschließende Abschnitte bestimmt]{\begin{aligned}[t]
  &m,n,i,j,h,k\in\mathbb N\dsep i+h=m\dsep j+k=n\dsep{}\\[-2pt]
  &p\colon\NatSeg{m+1}\to V\dsep q\colon\NatSeg{n+1}\to V\dsep{}\\[-2pt]
  &(i,p\restriction_{\NatSeg{i+1}})
      =(j,q\restriction_{\NatSeg{j+1}})\dsep{}\\[-2pt]
  &s\colon\NatSeg{h+1}\to V\dsep t\colon\NatSeg{k+1}\to V\dsep{}\\[-2pt]
  &\forall a\in\NatSeg{h+1}\,(s(a)=p(i+a))\dsep{}\\[-2pt]
  &\forall a\in\NatSeg{k+1}\,(t(a)=q(j+a))\dsep{}\\[-2pt]
  &(h,s)=(k,t)\vdash(m,p)=(n,q)
  \end{aligned}}{FinitePathSplitRigidity}{%
  \DeltaRow{Mengen}{V\dsep p\dsep q\dsep s\dsep t\dsep
    m\dsep n\dsep i\dsep j\dsep h\dsep k\dsep a\dsep b\dsep c}
}
\begin{tabproofsplitwide}
  \proofpartwide{Gleichheit der Abschnittslängen}
  \proofstepwidestarbreak[1]{m\in\mathbb N}{\rA}
  \proofstepwidestarbreak[2]{n\in\mathbb N}{\rA}
  \proofstepwidestarbreak[3]{i\in\mathbb N}{\rA}
  \proofstepwidestarbreak[4]{j\in\mathbb N}{\rA}
  \proofstepwidestarbreak[5]{h\in\mathbb N}{\rA}
  \proofstepwidestarbreak[6]{k\in\mathbb N}{\rA}
  \proofstepwidestarbreak[7]{i+h=m}{\rA}
  \proofstepwidestarbreak[8]{j+k=n}{\rA}
  \proofstepwidestarbreak[9]{p\colon\NatSeg{m+1}\to V}{\rA}
  \proofstepwidestarbreak[10]{q\colon\NatSeg{n+1}\to V}{\rA}
  \proofstepwidestarbreak[11]{%
    (i,p\restriction_{\NatSeg{i+1}})
      =(j,q\restriction_{\NatSeg{j+1}})}{\rA}
  \proofstepwidestarbreak[12]{s\colon\NatSeg{h+1}\to V}{\rA}
  \proofstepwidestarbreak[13]{t\colon\NatSeg{k+1}\to V}{\rA}
  \proofstepwidestarbreak[14]{%
    \forall a\in\NatSeg{h+1}\,(s(a)=p(i+a))}{\rA}
  \proofstepwidestarbreak[15]{%
    \forall a\in\NatSeg{k+1}\,(t(a)=q(j+a))}{\rA}
  \proofstepwidestarbreak[16]{(h,s)=(k,t)}{\rA}
  \proofstepwidestar[11]{i=j\land p\restriction_{\NatSeg{i+1}}=q\restriction_{\NatSeg{j+1}}}{%
      \FormulaRefAuto{(a,b)=(c,d)\eqvdash a=c\land b=d}[thm]{11}}
  \proofstepwidestarbreak[11]{i=j}{%
    \rAEa{17}}
  \proofstepwidestar[11]{i=j\land p\restriction_{\NatSeg{i+1}}=q\restriction_{\NatSeg{j+1}}}{%
      \FormulaRefAuto{(a,b)=(c,d)\eqvdash a=c\land b=d}[thm]{11}}
  \proofstepwidestarbreak[11]{%
    p\restriction_{\NatSeg{i+1}}=q\restriction_{\NatSeg{j+1}}}{%
    \rAEb{19}}
  \proofstepwidestar[16]{h=k\land s=t}{%
      \FormulaRefAuto{(a,b)=(c,d)\eqvdash a=c\land b=d}[thm]{16}}
  \proofstepwidestarbreak[16]{h=k}{%
    \rAEa{21}}
  \proofstepwidestar[16]{h=k\land s=t}{%
      \FormulaRefAuto{(a,b)=(c,d)\eqvdash a=c\land b=d}[thm]{16}}
  \proofstepwidestarbreak[16]{s=t}{%
    \rAEb{23}}
  \proofstepwidestarbreak[1,2,3,4,5,6,7,8,11,16]{m=n}{%
    \rIE{18,\rIE{22,\rIE{7,\rIE{8,\rII}}}}}
  \closeproofpartwide

  \proofpartwide{Übereinstimmung auf dem Anfangsabschnitt}
  \setcounter{proofstepnr}{25}
  \proofstepwidestarbreak[26]{a\in\NatSeg{m+1}}{\rA}
  \proofstepwidestar[1,3,5,7]{i\leq m\leftrightarrow\exists h\in\mathbb N\,(i+h=m)}{%
      \FormulaRefAuto{m\in\mathbb{N}\dsep n\in\mathbb{N}\vdash m\leq n\leftrightarrow\exists k\in\mathbb{N}\,(m+k=n)}[thm]{%
        3,1}}
  \proofstepwidestarbreak[1,3,5,7]{i\leq m}{%
    \FormulaRefAuto{P\leftrightarrow Q\dsep Q\vdash P}[thm]{%
      27,\rEI{\rAI{5,7}}}}
  \proofstepwidestar[1,3,26]{1 \in \mathbb N}{%
      \FormulaRefAuto{PeanoOneInN}[thm]}
  \proofstepwidestar[1,3,26]{m + 1 \in \mathbb N}{%
      \FormulaRefAuto{PeanoAddClosure}[thm]{1,
            29}}
  \proofstepwidestar[1,3,26]{\NatSeg { ( m + 1 ) } \subseteq \mathbb N}{%
      \FormulaRefAuto{PeanoNatSegSubsetN}[thm]{%
          30}}
  \proofstepwidestar[1,3,26]{a\in\mathbb N}{%
      \FormulaRefAuto{A\subseteq B\dsep x\in A\vdash x\in B}[thm]{%
        31,26}}
  \proofstepwidestarbreak[1,3,26]{a\leq i\lor i<a}{%
    \FormulaRefAuto{PeanoOrderComparison}[thm]{%
      32,3}}
  \proofstepwidestarbreak[34]{a\leq i}{\rA}
  \proofstepwidestar[3,26,34]{1 \in \mathbb N}{%
      \FormulaRefAuto{PeanoOneInN}[thm]}
  \proofstepwidestar[3,26,34]{m + 1 \in \mathbb N}{%
      \FormulaRefAuto{PeanoAddClosure}[thm]{1,
              35}}
  \proofstepwidestar[3,26,34]{\NatSeg { ( m + 1 ) } \subseteq \mathbb N}{%
      \FormulaRefAuto{PeanoNatSegSubsetN}[thm]{%
            36}}
  \proofstepwidestar[3,26,34]{a\in\mathbb N}{%
      \FormulaRefAuto{A\subseteq B\dsep x\in A\vdash x\in B}[thm]{%
          37,26}}
  \proofstepwidestar[3,26,34]{a\leq i\leftrightarrow a\in\NatSeg{i+1}}{%
      \FormulaRefAuto{PeanoNatSegLeqAddOneMembership}[thm]{%
        38,3}}
  \proofstepwidestarbreak[3,26,34]{a\in\NatSeg{i+1}}{%
    \FormulaRefAuto{P\leftrightarrow Q\dsep P\vdash Q}[thm]{%
      39,34}}
  \proofstepwidestar[9,26,34]{1 \in \mathbb N}{%
      \FormulaRefAuto{PeanoOneInN}[thm]}
  \proofstepwidestar[9,26,34]{1 \in \mathbb N}{%
      \FormulaRefAuto{PeanoOneInN}[thm]}
  \proofstepwidestar[9,26,34]{i + 1 \in \mathbb N}{%
      \FormulaRefAuto{PeanoAddClosure}[thm]{3,
          41}}
  \proofstepwidestar[9,26,34]{m + 1 \in \mathbb N}{%
      \FormulaRefAuto{PeanoAddClosure}[thm]{1,
          42}}
  \proofstepwidestar[9,26,34]{\NatSeg { ( i + 1 ) } \subseteq \NatSeg { ( m + 1 ) }}{%
      \FormulaRefAuto{PeanoNatSegLeqSubsetForward}[thm]{%
        43,
        44,28}}
  \proofstepwidestarbreak[9,26,34]{%
    p\restriction_{\NatSeg{i+1}}(a)=p(a)}{%
    \FormulaRefAuto{C\subseteq A\dsep x\in C\vdash
      F\restriction_C(x)=F(x)}[thm]{%
      45,40}}
  \proofstepwidestar[10,11,25,26,34]{1 \in \mathbb N}{%
      \FormulaRefAuto{PeanoOneInN}[thm]}
  \proofstepwidestar[10,11,25,26,34]{1 \in \mathbb N}{%
      \FormulaRefAuto{PeanoOneInN}[thm]}
  \proofstepwidestar[10,11,25,26,34]{j + 1 \in \mathbb N}{%
      \FormulaRefAuto{PeanoAddClosure}[thm]{4,
            47}}
  \proofstepwidestar[10,11,25,26,34]{n + 1 \in \mathbb N}{%
      \FormulaRefAuto{PeanoAddClosure}[thm]{2,
            48}}
  \proofstepwidestar[10,11,25,26,34]{\NatSeg { ( j + 1 ) } \subseteq \NatSeg { ( n + 1 ) }}{%
      \FormulaRefAuto{PeanoNatSegLeqSubsetForward}[thm]{%
          49,
          50,
          \rIE{18,\rIE{25,28}}}}
  \proofstepwidestar[10,11,25,26,34]{q\restriction_{\NatSeg{j+1}}(a)=q(a)}{%
      \FormulaRefAuto{C\subseteq A\dsep x\in C\vdash
        F\restriction_C(x)=F(x)}[thm]{%
        51,\rIE{18,40}}}
  \proofstepwidestarbreak[10,11,25,26,34]{%
    q\restriction_{\NatSeg{j+1}}(a)=q(a)}{%
    \rIE{18,
      52}}
  \proofstepwidestar[9,10,11,26,34]{{ p ( a ) } = { p \restriction _ { \NatSeg { i + 1 } } ( a ) }}{%
      \FormulaRefAuto{a=b\vdash b=a}[thm]{46}}
  \proofstepwidestarbreak[9,10,11,26,34]{p(a)=q(a)}{%
    \rChain{54,
      \rIE{20,\rII},53}}
  \closeproofpartwide

  \proofpartwide{Übereinstimmung auf dem Endabschnitt}
  \setcounter{proofstepnr}{55}
  \proofstepwidestarbreak[56]{i<a}{\rA}
  \proofstepwidestar[1,3,26,56]{1 \in \mathbb N}{%
      \FormulaRefAuto{PeanoOneInN}[thm]}
  \proofstepwidestar[1,3,26,56]{m + 1 \in \mathbb N}{%
      \FormulaRefAuto{PeanoAddClosure}[thm]{1,
              57}}
  \proofstepwidestar[1,3,26,56]{\NatSeg { ( m + 1 ) } \subseteq \mathbb N}{%
      \FormulaRefAuto{PeanoNatSegSubsetN}[thm]{%
            58}}
  \proofstepwidestar[1,3,26,56]{a \in \mathbb N}{%
      \FormulaRefAuto{A\subseteq B\dsep x\in A\vdash x\in B}[thm]{%
          59,26}}
  \proofstepwidestar[1,3,26,56]{i < a \leftrightarrow \exists b \in \mathbb N ( i + ( b + 1 ) = a )}{%
      \FormulaRefAuto{PeanoStrictAddOneDistance}[thm]{%
        3,
        60}}
  \proofstepwidestarbreak[1,3,26,56]{%
    \exists b\in\mathbb N\,(i+(b+1)=a)}{%
    \FormulaRefAuto{P\leftrightarrow Q\dsep P\vdash Q}[thm]{%
      61,56}}
  \proofstepwidestarbreak[63]{b\in\mathbb N\land i+(b+1)=a}{\rA}
  \proofstepwidestar[1,26]{1 \in \mathbb N}{%
      \FormulaRefAuto{PeanoOneInN}[thm]}
  \proofstepwidestar[1,26]{m + 1 \in \mathbb N}{%
      \FormulaRefAuto{PeanoAddClosure}[thm]{1,
              64}}
  \proofstepwidestar[1,26]{\NatSeg { ( m + 1 ) } \subseteq \mathbb N}{%
      \FormulaRefAuto{PeanoNatSegSubsetN}[thm]{%
            65}}
  \proofstepwidestar[1,26]{a\in\mathbb N}{%
      \FormulaRefAuto{A\subseteq B\dsep x\in A\vdash x\in B}[thm]{%
          66,26}}
  \proofstepwidestar[1,26]{a\leq m\leftrightarrow a\in\NatSeg{m+1}}{%
      \FormulaRefAuto{PeanoNatSegLeqAddOneMembership}[thm]{%
        67,1}}
  \proofstepwidestarbreak[1,26]{a\leq m}{%
    \FormulaRefAuto{P\leftrightarrow Q\dsep Q\vdash P}[thm]{%
      68,26}}
  \proofstepwidestar[1,3,5,7,26,63]{1 \in \mathbb N}{%
      \FormulaRefAuto{PeanoOneInN}[thm]}
  \proofstepwidestar[1,3,5,7,26,63]{b+1\in\mathbb N}{%
      \FormulaRefAuto{PeanoAddClosure}[thm]{\rAEa{63},
        70}}
  \proofstepwidestarbreak[1,3,5,7,26,63]{b+1\leq h}{%
    \FormulaRefAuto{PeanoAddLeftLeqCancellation}[thm]{%
      3,
      71,5,
      \rIE{\rAEb{63},\rIE{7,69}}}}
  \proofstepwidestar[5,63,72]{1 \in \mathbb N}{%
      \FormulaRefAuto{PeanoOneInN}[thm]}
  \proofstepwidestar[5,63,72]{b + 1 \in \mathbb N}{%
      \FormulaRefAuto{PeanoAddClosure}[thm]{\rAEa{63},
          73}}
  \proofstepwidestar[5,63,72]{b+1\leq h\leftrightarrow b+1\in\NatSeg{h+1}}{%
      \FormulaRefAuto{PeanoNatSegLeqAddOneMembership}[thm]{%
        74,5}}
  \proofstepwidestarbreak[5,63,72]{b+1\in\NatSeg{h+1}}{%
    \FormulaRefAuto{P\leftrightarrow Q\dsep P\vdash Q}[thm]{%
      75,72}}
  \proofstepwidestarbreak[14,63,72]{s(b+1)=p(a)}{%
    \rIE{\rAEb{63},\rUE{14,76}}}
  \proofstepwidestarbreak[2,5,6,8,11,15,16,63,72]{t(b+1)=q(a)}{%
    \rIE{18,\rIE{22,\rIE{\rAEb{63},\rUE{15,\rIE{22,76}}}}}}
  \proofstepwidestar[14,15,16,63,72]{{ p ( a ) } = { s ( b + 1 ) }}{%
      \FormulaRefAuto{a=b\vdash b=a}[thm]{77}}
  \proofstepwidestarbreak[14,15,16,63,72]{p(a)=q(a)}{%
    \rChain{79,
      \rIE{24,\rII},78}}
  \closeproofpartwide

  \proofpartwide{Zusammenführung und Funktionsextensionalität}
  \setcounter{proofstepnr}{80}
  \proofstepwidestarbreak[1,2,3,4,5,6,7,8,9,10,11,12,13,14,15,16,26]{%
    p(a)=q(a)}{%
    \rOE{33,34,55,56,\rEE{62,63,80}}}
  \proofstepwidestarbreak[1,2,3,4,5,6,7,8,9,10,11,12,13,14,15,16]{%
    \forall a\in\NatSeg{m+1}\,(p(a)=q(a))}{%
    \rUI{\rRI{26,81}}}
  \proofstepwidestarbreak[1,2,9,10,25,82]{p=q}{%
    \FormulaRefAuto{FunctionExtensionality}[thm]{9,\rIE{25,10},82}}
  \proofstepwidestarbreak[1,2,9,10,25,82]{(m,p)=(n,q)}{%
    \FormulaRefAuto{a=c,\, b=d\vdash(a,b)=(c,d)}[thm]{25,83}}
  \closeproofpartwide
\end{tabproofsplitwide}

\FormulaThmDeltaK[Reduktion an einem gemeinsamen inneren Wegknoten]{\begin{aligned}[t]
  &\UndirectedGraphStruct{V,E}\dsep m,n\in\mathbb N\dsep
    p,q\in\powerset(\mathbb N\times V)\dsep{}\\[-2pt]
  &\DirectedPath{V}{E}{p}{m}\dsep\DirectedPath{V}{E}{q}{n}\dsep{}\\[-2pt]
  &p(0)=q(0)\dsep p(m)=q(n)\dsep(m,p)\neq(n,q)\dsep{}\\[-2pt]
  &z\in\bigl(p[\NatSeg{m+1}]\cap q[\NatSeg{n+1}]\bigr)
      \setminus\{p(0),p(m)\}\vdash{}\\[-2pt]
  &\quad\exists h\in\mathbb N\,\exists k\in\mathbb N\,
    \exists s\in\powerset(\mathbb N\times V)\,
    \exists t\in\powerset(\mathbb N\times V)\,\Bigl(\\[-2pt]
  &\qquad\DirectedPath{V}{E}{s}{h}\land\DirectedPath{V}{E}{t}{k}\land{}\\[-2pt]
  &\qquad s(0)=t(0)\land s(h)=t(k)\land
    (h,s)\neq(k,t)\land h+k<m+n\Bigr)
  \end{aligned}}{FinitePathInteriorPairReduction}{%
  \DeltaRow{Mengen}{V\dsep E\dsep p\dsep q\dsep s\dsep t\dsep
    m\dsep n\dsep h\dsep k\dsep i\dsep j\dsep z}
}
\begin{tabproofsplitwide}
  \proofpartwide{Wahl der Schnittstellen und Wegabschnitte}
  \proofstepwidestarbreak[1]{\UndirectedGraphStruct{V,E}}{\rA}
  \proofstepwidestarbreak[2]{m\in\mathbb N}{\rA}
  \proofstepwidestarbreak[3]{n\in\mathbb N}{\rA}
  \proofstepwidestarbreak[4]{p\in\powerset(\mathbb N\times V)}{\rA}
  \proofstepwidestarbreak[5]{q\in\powerset(\mathbb N\times V)}{\rA}
  \proofstepwidestarbreak[6]{\DirectedPath{V}{E}{p}{m}}{\rA}
  \proofstepwidestarbreak[7]{\DirectedPath{V}{E}{q}{n}}{\rA}
  \proofstepwidestarbreak[8]{p(0)=q(0)}{\rA}
  \proofstepwidestarbreak[9]{p(m)=q(n)}{\rA}
  \proofstepwidestarbreak[10]{(m,p)\neq(n,q)}{\rA}
  \proofstepwidestarbreak[11]{%
    z\in\bigl(p[\NatSeg{m+1}]\cap q[\NatSeg{n+1}]\bigr)
      \setminus\{p(0),p(m)\}}{\rA}
  \proofstepwidestar[4,11]{\GraphStruct { V , E }}{%
      \FormulaRefAuto{UndirectedGraphAccess}[thm]{1}}
  \proofstepwidestar[4,11]{\GraphStruct { V , E }}{%
      \FormulaRefAuto{UndirectedGraphAccess}[thm]{1}}
  \proofstepwidestar[4,11]{\DirectedWalk V E p m}{%
      \FormulaRefAuto{DirectedPathWalkAxiom}[ax]{%
            13,6}}
  \proofstepwidestar[4,11]{z\in p[\NatSeg{m+1}]\cap q[\NatSeg{n+1}]}{%
      \FormulaRefAuto{x\in A\setminus B\vdash x\in A}[thm]{11}}
  \proofstepwidestar[4,11]{\begin{aligned}[t]&m\in\mathbb N\land p\colon\NatSeg{m+1}\to V\land{}\\[-2pt]&\forall i\in\NatSeg{m}\,((p(i),p(i+1))\in E)\end{aligned}}{%
      \FormulaRefAuto{DirectedWalkCharacterization}[thm]{%
          12,
          14}}
  \proofstepwidestar[4,11]{z\in p[\NatSeg{m+1}]}{%
      \FormulaRefAuto{x\in A\cap B\vdash x\in A}[thm]{%
        15}}
  \proofstepwidestar[4,11]{p\colon\NatSeg{m+1}\to V}{%
      \FormulaRefAuto{P\land Q\land R\vdash Q}[thm]{%
        16}}
  \proofstepwidestarbreak[4,11]{%
    \exists i\in\NatSeg{m+1}\,(p(i)=z)}{%
    \FormulaRefAuto{F\colon A\to B\dsep y\in F[A]\vdash
      \exists x\in A\,y=F(x)}[thm]{%
      18,
      17}}
  \proofstepwidestar[5,11]{\GraphStruct { V , E }}{%
      \FormulaRefAuto{UndirectedGraphAccess}[thm]{1}}
  \proofstepwidestar[5,11]{\GraphStruct { V , E }}{%
      \FormulaRefAuto{UndirectedGraphAccess}[thm]{1}}
  \proofstepwidestar[5,11]{\DirectedWalk V E q n}{%
      \FormulaRefAuto{DirectedPathWalkAxiom}[ax]{%
            21,7}}
  \proofstepwidestar[5,11]{z\in p[\NatSeg{m+1}]\cap q[\NatSeg{n+1}]}{%
      \FormulaRefAuto{x\in A\setminus B\vdash x\in A}[thm]{11}}
  \proofstepwidestar[5,11]{\begin{aligned}[t]&n\in\mathbb N\land q\colon\NatSeg{n+1}\to V\land{}\\[-2pt]&\forall i\in\NatSeg{n}\,((q(i),q(i+1))\in E)\end{aligned}}{%
      \FormulaRefAuto{DirectedWalkCharacterization}[thm]{%
          20,
          22}}
  \proofstepwidestar[5,11]{z\in q[\NatSeg{n+1}]}{%
      \FormulaRefAuto{x\in A\cap B\vdash x\in B}[thm]{%
        23}}
  \proofstepwidestar[5,11]{q\colon\NatSeg{n+1}\to V}{%
      \FormulaRefAuto{P\land Q\land R\vdash Q}[thm]{%
        24}}
  \proofstepwidestarbreak[5,11]{%
    \exists j\in\NatSeg{n+1}\,(q(j)=z)}{%
    \FormulaRefAuto{F\colon A\to B\dsep y\in F[A]\vdash
      \exists x\in A\,y=F(x)}[thm]{%
      26,
      25}}
  \proofstepwidestarbreak[28]{i\in\NatSeg{m+1}\land p(i)=z}{\rA}
  \proofstepwidestarbreak[29]{j\in\NatSeg{n+1}\land q(j)=z}{\rA}
  \proofstepwidestar[6,28]{\GraphStruct { V , E }}{%
      \FormulaRefAuto{UndirectedGraphAccess}[thm]{1}}
  \proofstepwidestar[6,28]{\DirectedPath V E { p \restriction _ { \NatSeg { i + 1 } } } i \land p \restriction _ { \NatSeg { i + 1 } } ( 0 ) = p ( 0 ) \land p \restriction _ { \NatSeg { i + 1 } } ( i ) = p ( i )}{%
      \FormulaRefAuto{FinitePathInitialSegment}[thm]{%
        30,6,\rAEa{28}}}
  \proofstepwidestarbreak[6,28]{%
    \DirectedPath{V}{E}{p\restriction_{\NatSeg{i+1}}}{i}\land
    p\restriction_{\NatSeg{i+1}}(0)=p(0)\land
    p\restriction_{\NatSeg{i+1}}(i)=z}{%
    \rIE{\rAEb{28},
      31}}
  \proofstepwidestar[7,29]{\GraphStruct { V , E }}{%
      \FormulaRefAuto{UndirectedGraphAccess}[thm]{1}}
  \proofstepwidestar[7,29]{\DirectedPath V E { q \restriction _ { \NatSeg { j + 1 } } } j \land q \restriction _ { \NatSeg { j + 1 } } ( 0 ) = q ( 0 ) \land q \restriction _ { \NatSeg { j + 1 } } ( j ) = q ( j )}{%
      \FormulaRefAuto{FinitePathInitialSegment}[thm]{%
        33,7,\rAEa{29}}}
  \proofstepwidestarbreak[7,29]{%
    \DirectedPath{V}{E}{q\restriction_{\NatSeg{j+1}}}{j}\land
    q\restriction_{\NatSeg{j+1}}(0)=q(0)\land
    q\restriction_{\NatSeg{j+1}}(j)=z}{%
    \rIE{\rAEb{29},
      34}}
  \proofstepwidestar[1,2,4,6,28]{\GraphStruct { V , E }}{%
      \FormulaRefAuto{UndirectedGraphAccess}[thm]{1}}
  \proofstepwidestar[1,2,4,6,28]{\begin{aligned}[t]&\exists h\in\mathbb N\,\exists s\in\powerset(\mathbb N\times V)\,\bigl(\\[-2pt]&i+h=m\land\DirectedPath{V}{E}{s}{h}\land{}\\[-2pt]&s(0)=p(i)\land s(h)=p(m)\land{}\\[-2pt]&\forall k\in\NatSeg{h+1}\,(s(k)=p(i+k))\bigr)\end{aligned}}{%
      \FormulaRefAuto{FinitePathFinalSegment}[thm]{%
        36,2,4,6,\rAEa{28}}}
  \proofstepwidestarbreak[1,2,4,6,28]{%
    \exists h\in\mathbb N\,\exists s\in\powerset(\mathbb N\times V)\,
    \bigl(i+h=m\land\DirectedPath{V}{E}{s}{h}\land
      s(0)=z\land s(h)=p(m)\land
      \forall k\in\NatSeg{h+1}\,(s(k)=p(i+k))\bigr)}{%
    \rIE{\rAEb{28},
      37}}
  \proofstepwidestar[1,3,5,7,29]{\GraphStruct { V , E }}{%
      \FormulaRefAuto{UndirectedGraphAccess}[thm]{1}}
  \proofstepwidestar[1,3,5,7,29]{\begin{aligned}[t]&\exists k\in\mathbb N\,\exists t\in\powerset(\mathbb N\times V)\,\bigl(\\[-2pt]&j+k=n\land\DirectedPath{V}{E}{t}{k}\land{}\\[-2pt]&t(0)=q(j)\land t(k)=q(n)\land{}\\[-2pt]&\forall h\in\NatSeg{k+1}\,(t(h)=q(j+h))\bigr)\end{aligned}}{%
      \FormulaRefAuto{FinitePathFinalSegment}[thm]{%
        39,3,5,7,\rAEa{29}}}
  \proofstepwidestarbreak[1,3,5,7,29]{%
    \exists k\in\mathbb N\,\exists t\in\powerset(\mathbb N\times V)\,
    \bigl(j+k=n\land\DirectedPath{V}{E}{t}{k}\land
      t(0)=z\land t(k)=q(n)\land
      \forall h\in\NatSeg{k+1}\,(t(h)=q(j+h))\bigr)}{%
    \rIE{\rAEb{29},
      40}}
  \closeproofpartwide

  \proofpartwide{Strikte Verkürzung der Anfangsabschnitte}
  \setcounter{proofstepnr}{41}
  \proofstepwidestar[11]{z\notin\{p(0),p(m)\}\leftrightarrow(z\neq p(0)\land z\neq p(m))}{%
      \FormulaRefAuto{NotInPairSet}[thm]}
  \proofstepwidestar[11]{z\notin\{p(0),p(m)\}}{%
      \FormulaRefAuto{x\in A\setminus B\vdash x\notin B}[thm]{11}}
  \proofstepwidestar[11]{z\neq p(0)\land z\neq p(m)}{%
      \FormulaRefAuto{P\leftrightarrow Q\dsep P\vdash Q}[thm]{%
      42,
      43}}
  \proofstepwidestarbreak[11]{z\neq p(0)}{%
    \rAEa{44}}
  \proofstepwidestar[11]{z\notin\{p(0),p(m)\}\leftrightarrow(z\neq p(0)\land z\neq p(m))}{%
      \FormulaRefAuto{NotInPairSet}[thm]}
  \proofstepwidestar[11]{z\notin\{p(0),p(m)\}}{%
      \FormulaRefAuto{x\in A\setminus B\vdash x\notin B}[thm]{11}}
  \proofstepwidestar[11]{z\neq p(0)\land z\neq p(m)}{%
      \FormulaRefAuto{P\leftrightarrow Q\dsep P\vdash Q}[thm]{%
      46,
      47}}
  \proofstepwidestarbreak[11]{z\neq p(m)}{%
    \rAEb{48}}
  \proofstepwidestarbreak[6,28,45]{i\neq0}{%
    \rCI{\rA,\rBI{%
      \FormulaRefAuto{a=b\vdash b=a}[thm]{\rIE{\rA,\rAEb{28}}},45}}}
  \proofstepwidestarbreak[6,28,49]{i\neq m}{%
    \rCI{\rA,\rBI{%
      \FormulaRefAuto{a=b\vdash b=a}[thm]{\rIE{\rA,\rAEb{28}}},49}}}
  \proofstepwidestarbreak[7,8,29,45]{j\neq0}{%
    \rCI{\rA,\rBI{%
      \FormulaRefAuto{a=b\vdash b=a}[thm]{%
        \rIE{8,\rIE{\rA,\rAEb{29}}}},45}}}
  \proofstepwidestarbreak[7,9,29,49]{j\neq n}{%
    \rCI{\rA,\rBI{%
      \FormulaRefAuto{a=b\vdash b=a}[thm]{%
        \rIE{9,\rIE{\rA,\rAEb{29}}}},49}}}
  \proofstepwidestar[2,28]{i\in\mathbb N\land i\leq m\land\NatSeg{i+1}\subseteq\NatSeg{m+1}}{%
      \FormulaRefAuto{FinitePathPrefixIndexData}[thm]{2,\rAEa{28}}}
  \proofstepwidestar[2,28]{i\in\mathbb N\land i\leq m\land\NatSeg{i+1}\subseteq\NatSeg{m+1}}{%
      \FormulaRefAuto{FinitePathPrefixIndexData}[thm]{%
        2,\rAEa{28}}}
  \proofstepwidestarbreak[2,28]{i\in\mathbb N\land i\leq m}{%
    \rAI{%
      \rAEa{54},
      \rAEa{\rAEb{55}}}}
  \proofstepwidestar[3,29]{j\in\mathbb N\land j\leq n\land\NatSeg{j+1}\subseteq\NatSeg{n+1}}{%
      \FormulaRefAuto{FinitePathPrefixIndexData}[thm]{3,\rAEa{29}}}
  \proofstepwidestar[3,29]{j\in\mathbb N\land j\leq n\land\NatSeg{j+1}\subseteq\NatSeg{n+1}}{%
      \FormulaRefAuto{FinitePathPrefixIndexData}[thm]{%
        3,\rAEa{29}}}
  \proofstepwidestarbreak[3,29]{j\in\mathbb N\land j\leq n}{%
    \rAI{%
      \rAEa{57},
      \rAEa{\rAEb{58}}}}
  \proofstepwidestarbreak[2,28,51]{i<m}{%
    \FormulaRefAuto{PeanoLeqNeqStrict}[thm]{%
      \rAEa{56},2,\rAEb{56},51}}
  \proofstepwidestarbreak[3,29,53]{j<n}{%
    \FormulaRefAuto{PeanoLeqNeqStrict}[thm]{%
      \rAEa{59},3,\rAEb{59},53}}
  \proofstepwidestarbreak[2,3,28,29,51,53]{i+j<m+n}{%
    \FormulaRefAuto{PeanoAddStrictMonotoneBoth}[thm]{%
      \rAEa{56},2,\rAEa{59},3,60,61}}
  \closeproofpartwide

  \proofpartwide{Fall verschiedener Anfangsabschnitte}
  \setcounter{proofstepnr}{62}
  \proofstepwidestarbreak[]{%
    (i,p\restriction_{\NatSeg{i+1}})\neq
      (j,q\restriction_{\NatSeg{j+1}})\lor
    (i,p\restriction_{\NatSeg{i+1}})=
      (j,q\restriction_{\NatSeg{j+1}})}{%
    \FormulaRefAuto{P\lor\neg P}[thm]}
  \proofstepwidestarbreak[64]{%
    (i,p\restriction_{\NatSeg{i+1}})\neq
      (j,q\restriction_{\NatSeg{j+1}})}{\rA}
  \proofstepwidestarbreak[1,2,3,4,5,6,7,8,9,28,29,60,61,64]{%
    \exists h\in\mathbb N\,\exists k\in\mathbb N\,
    \exists s\in\powerset(\mathbb N\times V)\,
    \exists t\in\powerset(\mathbb N\times V)\,\Bigl(
      \DirectedPath{V}{E}{s}{h}\land\DirectedPath{V}{E}{t}{k}\land
      s(0)=t(0)\land s(h)=t(k)\land
      (h,s)\neq(k,t)\land h+k<m+n\Bigr)}{%
    \rEI{\rAI{\rAEa{56},\rEI{\rAI{\rAEa{59},
      \rEI{\rAI{4,\rEI{\rAI{5,\rAEa{32},\rAEa{35},
        \rIE{8,\rIE{\rAEa{\rAEn{32}},\rAEa{\rAEn{35}}}},
        \rIE{\rAEn{32},\rAEn{35}},64,62}}}}}}}}}
  \closeproofpartwide

  \proofpartwide{Fall gleicher Anfangsabschnitte}
  \setcounter{proofstepnr}{65}
  \proofstepwidestarbreak[66]{%
    (i,p\restriction_{\NatSeg{i+1}})=
      (j,q\restriction_{\NatSeg{j+1}})}{\rA}
  \proofstepwidestarbreak[38]{h\in\mathbb N\land
    s\in\powerset(\mathbb N\times V)\land i+h=m\land
    \DirectedPath{V}{E}{s}{h}\land s(0)=z\land s(h)=p(m)\land
    \forall a\in\NatSeg{h+1}\,(s(a)=p(i+a))}{\rA}
  \proofstepwidestarbreak[41]{k\in\mathbb N\land
    t\in\powerset(\mathbb N\times V)\land j+k=n\land
    \DirectedPath{V}{E}{t}{k}\land t(0)=z\land t(k)=q(n)\land
    \forall a\in\NatSeg{k+1}\,(t(a)=q(j+a))}{\rA}
  \proofstepwidestarbreak[69]{(h,s)=(k,t)}{\rA}
  \proofstepwidestar[1,2,3,6,7,66,67,68,69]{\GraphStruct { V , E }}{%
      \FormulaRefAuto{UndirectedGraphAccess}[thm]{1}}
  \proofstepwidestar[1,2,3,6,7,66,67,68,69]{\GraphStruct { V , E }}{%
      \FormulaRefAuto{UndirectedGraphAccess}[thm]{1}}
  \proofstepwidestar[1,2,3,6,7,66,67,68,69]{\GraphStruct { V , E }}{%
      \FormulaRefAuto{UndirectedGraphAccess}[thm]{1}}
  \proofstepwidestar[1,2,3,6,7,66,67,68,69]{\GraphStruct { V , E }}{%
      \FormulaRefAuto{UndirectedGraphAccess}[thm]{1}}
  \proofstepwidestar[1,2,3,6,7,66,67,68,69]{p \colon \NatSeg { m + 1 } \to V}{%
      \FormulaRefAuto{DirectedPathFunction}[thm]{%
        70,6}}
  \proofstepwidestar[1,2,3,6,7,66,67,68,69]{q \colon \NatSeg { n + 1 } \to V}{%
      \FormulaRefAuto{DirectedPathFunction}[thm]{%
        71,7}}
  \proofstepwidestar[1,2,3,6,7,66,67,68,69]{s \colon \NatSeg { h + 1 } \to V}{%
      \FormulaRefAuto{DirectedPathFunction}[thm]{%
        72,\rAEn{67}}}
  \proofstepwidestar[1,2,3,6,7,66,67,68,69]{t \colon \NatSeg { k + 1 } \to V}{%
      \FormulaRefAuto{DirectedPathFunction}[thm]{%
        73,\rAEn{68}}}
  \proofstepwidestarbreak[1,2,3,6,7,66,67,68,69]{(m,p)=(n,q)}{%
    \FormulaRefAuto{FinitePathSplitRigidity}[thm]{%
      2,3,\rAEa{56},\rAEa{59},\rAEa{67},\rAEa{68},
      \rAEb{\rAEb{67}},\rAEb{\rAEb{68}},
      74,
      75,66,
      76,
      77,
      \rAEn{67},\rAEn{68},69}}
  \proofstepwidestarbreak[10,66,67,68]{(h,s)\neq(k,t)}{%
    \rCI{69,\rBI{78,10}}}
  \proofstepwidestar[2,50,56,67]{i + h = h + i}{%
      \FormulaRefAuto{PeanoAddCommutative}[thm]{%
         \rAEa{56},\rAEa{67}}}
  \proofstepwidestarbreak[2,50,56,67]{h<m}{%
    \FormulaRefAuto{m\in\mathbb{N}\dsep n\in\mathbb{N}\dsep k\in\mathbb{N}\dsep m+k=n\dsep k\neq0\vdash m<n}[thm]{%
      \rAEa{67},2,\rAEa{56},
       \rIE{80,\rAEb{\rAEb{67}}},45}}
  \proofstepwidestar[3,52,59,68]{j + k = k + j}{%
      \FormulaRefAuto{PeanoAddCommutative}[thm]{%
         \rAEa{59},\rAEa{68}}}
  \proofstepwidestarbreak[3,52,59,68]{k<n}{%
    \FormulaRefAuto{m\in\mathbb{N}\dsep n\in\mathbb{N}\dsep k\in\mathbb{N}\dsep m+k=n\dsep k\neq0\vdash m<n}[thm]{%
      \rAEa{68},3,\rAEa{59},
       \rIE{82,\rAEb{\rAEb{68}}},50}}
  \proofstepwidestarbreak[2,3,67,68,81,83]{h+k<m+n}{%
    \FormulaRefAuto{PeanoAddStrictMonotoneBoth}[thm]{%
      \rAEa{67},2,\rAEa{68},3,81,83}}
  \proofstepwidestarbreak[1,2,3,4,5,6,7,8,9,10,66,67,68]{%
    \exists h\in\mathbb N\,\exists k\in\mathbb N\,
    \exists s\in\powerset(\mathbb N\times V)\,
    \exists t\in\powerset(\mathbb N\times V)\,\Bigl(
      \DirectedPath{V}{E}{s}{h}\land\DirectedPath{V}{E}{t}{k}\land
      s(0)=t(0)\land s(h)=t(k)\land
      (h,s)\neq(k,t)\land h+k<m+n\Bigr)}{%
    \rEI{\rAI{\rAEa{67},\rEI{\rAI{\rAEa{68},
      \rEI{\rAI{\rAEb{67},\rEI{\rAI{\rAEb{68},
        \rAEn{67},\rAEn{68},\rIE{\rAEn{67},\rAEn{68}},
        \rIE{9,\rIE{\rAEn{67},\rAEn{68}}},79,84}}}}}}}}}
  \proofstepwidestarbreak[1,2,3,4,5,6,7,8,9,10,66]{%
    \exists h\in\mathbb N\,\exists k\in\mathbb N\,
    \exists s\in\powerset(\mathbb N\times V)\,
    \exists t\in\powerset(\mathbb N\times V)\,\Bigl(
      \DirectedPath{V}{E}{s}{h}\land\DirectedPath{V}{E}{t}{k}\land
      s(0)=t(0)\land s(h)=t(k)\land
      (h,s)\neq(k,t)\land h+k<m+n\Bigr)}{%
    \rEE{38,67,\rEE{41,68,85}}}
  \closeproofpartwide

  \proofpartwide{Zusammenführung der Fälle}
  \setcounter{proofstepnr}{86}
  \proofstepwidestarbreak[1,2,3,4,5,6,7,8,9,10,28,29]{%
    \exists h\in\mathbb N\,\exists k\in\mathbb N\,
    \exists s\in\powerset(\mathbb N\times V)\,
    \exists t\in\powerset(\mathbb N\times V)\,\Bigl(
      \DirectedPath{V}{E}{s}{h}\land\DirectedPath{V}{E}{t}{k}\land
      s(0)=t(0)\land s(h)=t(k)\land
      (h,s)\neq(k,t)\land h+k<m+n\Bigr)}{%
    \rOE{63,64,65,66,86}}
  \proofstepwidestarbreak[1,2,3,4,5,6,7,8,9,10,11]{%
    \exists h\in\mathbb N\,\exists k\in\mathbb N\,
    \exists s\in\powerset(\mathbb N\times V)\,
    \exists t\in\powerset(\mathbb N\times V)\,\Bigl(
      \DirectedPath{V}{E}{s}{h}\land\DirectedPath{V}{E}{t}{k}\land
      s(0)=t(0)\land s(h)=t(k)\land
      (h,s)\neq(k,t)\land h+k<m+n\Bigr)}{%
    \rEE{19,28,\rEE{27,29,87}}}
  \closeproofpartwide
\end{tabproofsplitwide}

\FormulaThmDeltaK[Intern disjunkte Wegbögen tragen einen endlichen Wegzyklus]{\begin{aligned}[t]
  &\UndirectedGraphStruct{V,E}\dsep x,y\in V\dsep m,n\in\mathbb N\dsep{}\\[-2pt]
  &p,q\in\powerset(\mathbb N\times V)\dsep
    \DirectedPath{V}{E}{p}{m}\dsep\DirectedPath{V}{E}{q}{n}\dsep{}\\[-2pt]
  &p(0)=x\dsep p(m)=y\dsep q(0)=x\dsep q(n)=y\dsep{}\\[-2pt]
  &x\neq y\dsep(m,p)\neq(n,q)\dsep
    p[\NatSeg{m+1}]\cap q[\NatSeg{n+1}]=\{x,y\}\vdash{}\\[-2pt]
  &\quad U\subseteq V\land U\neq\varnothing\land{}\\[-2pt]
  &\quad\forall z\in U\,\exists w\in U\,
    \exists h\in\mathbb N\,\exists k\in\mathbb N\,{}\\[-2pt]
  &\quad\exists s\in\powerset(\mathbb N\times V)\,
    \exists t\in\powerset(\mathbb N\times V)\,\Bigl(z\neq w\land{}\\[-2pt]
  &\qquad\DirectedPath{V}{E}{s}{h}\land s(0)=z\land s(h)=w\land{}\\[-2pt]
  &\qquad\DirectedPath{V}{E}{t}{k}\land t(0)=z\land t(k)=w\land
    (h,s)\neq(k,t)\land{}\\[-2pt]
  &\qquad s[\NatSeg{h+1}]\cap t[\NatSeg{k+1}]=\{z,w\}\land{}\\[-2pt]
  &\qquad s[\NatSeg{h+1}]\cup t[\NatSeg{k+1}]\subseteq U\Bigr)
  \end{aligned}}{InternallyDisjointPathsFormCycleCore}{%
  \DeltaRow{Mengen}{V\dsep E\dsep x\dsep y\dsep z\dsep w\dsep
    m\dsep n\dsep h\dsep k\dsep p\dsep q\dsep s\dsep t}
  \DeltaRow{Abkürzung}{U\coloneqq p[\NatSeg{m+1}]\cup q[\NatSeg{n+1}]}
}
\paragraph{Begründung der Knotenbedingung.}\label{proof:internally-disjoint-cycle-node}
Schreibe \(U=p[\NatSeg{m+1}]\cup q[\NatSeg{n+1}]\).
Für jeden Knoten \(z\in U\) werden zwei verschiedene, bis auf ihre
Endpunkte disjunkte Wege von \(z\) zu einem anderen Knoten \(w\in U\)
konstruiert. Für \(z=x\) dienen \(p,q\) mit \(w=y\). Für \(z=y\)
werden beide Folgen umgekehrt und \(w=x\) gewählt. Die Umkehrung ist
injektiv, erhält das Bild und wegen der Ungerichtetheit auch die
Wegeigenschaft; zweimaliges Umkehren ergibt die ursprüngliche Folge,
also bleiben die beiden gespeicherten Wege verschieden.

Sei nun \(z\in p[\NatSeg{m+1}]\setminus\{x,y\}\). Die Injektivität
von \(p\) liefert einen eindeutigen Index \(i\) mit \(p(i)=z\) und
\(0<i<m\). Sei \(h\in\mathbb N\) durch \(i+h=m\) bestimmt. Setze
\(w=y\) und definiere die beiden Folgen durch
\[
  s(j)=p(i+j)\quad(j\in\NatSeg{h+1}),\qquad
  t(j)=
  \begin{cases}
    p(i-j),&0\leq j\leq i,\\
    q(j-i),&i<j\leq i+n.
  \end{cases}
\]
Die Differenzen bezeichnen jeweils die eindeutigen natürlichen
Additionsabstände. Das Endstück \(s\) ist ein Weg der Länge \(h\).
Die Folge \(t\) läuft den Anfang von \(p\) rückwärts bis \(x\) und
anschließend \(q\) vorwärts bis \(y\). Alle ihre aufeinanderfolgenden
Knoten sind daher durch Kanten verbunden. Innerhalb beider Stücke ist
\(t\) injektiv. Ein Knoten beider Stücke müsste wegen der inneren
Disjunktheit \(x\) oder \(y\) sein. Im Anfang von \(p\) tritt \(y\)
nicht auf, und \(x\) ist in der Definition von \(t\) nur am
Übergangsindex \(i\) enthalten. Damit ist auch \(t\) injektiv.

Es gilt \(s(0)=t(0)=z\), \(s(h)=t(i+n)=y\), und beide Bilder liegen
in \(U\). Anfang und Ende von \(p\) treffen sich wegen der
Injektivität genau in \(z\); das Endstück von \(p\) und \(q\)
treffen sich genau in \(y\). Somit
\[
  s[\NatSeg{h+1}]\cap t[\NatSeg{(i+n)+1}]=\{z,y\}.
\]
Da \(i>0\) und \(n>0\), besitzt \(t\) mindestens drei verschiedene
Knoten. Die beiden gespeicherten Wege können deshalb nicht gleich
sein: Sonst wäre ihr gemeinsames Bild die Zweiermenge \(\{z,y\}\).
Für \(z\in q[\NatSeg{n+1}]\setminus\{x,y\}\) gilt dieselbe
Konstruktion mit vertauschten \(p,m\) und \(q,n\). Diese Fälle
erschöpfen \(U\) und beweisen die in der Tabelle verwendete
Knotenbedingung. Die Anfangs-, End- und Umkehrstücke entsprechen
\FormulaRefAuto{FinitePathInitialSegment}[thm],
\FormulaRefAuto{FinitePathFinalSegment}[thm] und
\FormulaRefAuto{UndirectedFinitePathReversal}[thm].

\begin{tabproofwide}
  \proofstepwidestarbreak[1]{\UndirectedGraphStruct{V,E}}{\rA}
  \proofstepwidestarbreak[2]{x\in V}{\rA}
  \proofstepwidestarbreak[3]{y\in V}{\rA}
  \proofstepwidestarbreak[4]{m\in\mathbb N}{\rA}
  \proofstepwidestarbreak[5]{n\in\mathbb N}{\rA}
  \proofstepwidestarbreak[6]{p\in\powerset(\mathbb N\times V)}{\rA}
  \proofstepwidestarbreak[7]{q\in\powerset(\mathbb N\times V)}{\rA}
  \proofstepwidestarbreak[8]{\DirectedPath{V}{E}{p}{m}}{\rA}
  \proofstepwidestarbreak[9]{\DirectedPath{V}{E}{q}{n}}{\rA}
  \proofstepwidestarbreak[10]{p(0)=x\land p(m)=y\land
    q(0)=x\land q(n)=y}{\rA}
  \proofstepwidestarbreak[11]{x\neq y}{\rA}
  \proofstepwidestarbreak[12]{(m,p)\neq(n,q)}{\rA}
  \proofstepwidestarbreak[13]{%
    p[\NatSeg{m+1}]\cap q[\NatSeg{n+1}]=\{x,y\}}{\rA}
  \proofstepwidestar[1,4,5,6,7,8,9]{\GraphStruct { V , E }}{%
      \FormulaRefAuto{UndirectedGraphAccess}[thm]{1}}
  \proofstepwidestar[1,4,5,6,7,8,9]{\GraphStruct { V , E }}{%
      \FormulaRefAuto{UndirectedGraphAccess}[thm]{1}}
  \proofstepwidestar[1,4,5,6,7,8,9]{p \colon \NatSeg { m + 1 } \to V}{%
      \FormulaRefAuto{DirectedPathFunction}[thm]{%
          14,8}}
  \proofstepwidestar[1,4,5,6,7,8,9]{q \colon \NatSeg { n + 1 } \to V}{%
      \FormulaRefAuto{DirectedPathFunction}[thm]{%
          15,9}}
  \proofstepwidestar[1,4,5,6,7,8,9]{p [ { \NatSeg { m + 1 } } ] \subseteq V}{%
      \FormulaRefAuto{F\colon A\to B\vdash F[A]\subseteq B}[thm]{%
        16}}
  \proofstepwidestar[1,4,5,6,7,8,9]{q [ { \NatSeg { n + 1 } } ] \subseteq V}{%
      \FormulaRefAuto{F\colon A\to B\vdash F[A]\subseteq B}[thm]{%
        17}}
  \proofstepwidestarbreak[1,4,5,6,7,8,9]{%
    p[\NatSeg{m+1}]\cup q[\NatSeg{n+1}]\subseteq V}{%
    \FormulaRefAuto{A\subseteq C\dsep B\subseteq C\vdash
      A\cup B\subseteq C}[thm]{%
      18,
      19}}
  \proofstepwidestar[2,4,6,8,10]{0 \in \mathbb N}{%
      \FormulaRefAuto{PeanoZero}[ax]}
  \proofstepwidestar[2,4,6,8,10]{\GraphStruct { V , E }}{%
      \FormulaRefAuto{UndirectedGraphAccess}[thm]{1}}
  \proofstepwidestar[2,4,6,8,10]{p \colon \NatSeg { m + 1 } \to V}{%
      \FormulaRefAuto{DirectedPathFunction}[thm]{%
            22,8}}
  \proofstepwidestar[2,4,6,8,10]{\NatSeg{m+1}\subseteq\NatSeg{m+1}}{%
      \FormulaRefAuto{A\subseteq A}[thm]}
  \proofstepwidestar[2,4,6,8,10]{0\leq m}{%
      \FormulaRefAuto{PeanoZeroLeq}[thm]{4}}
  \proofstepwidestar[2,4,6,8,10]{p[\NatSeg{m+1}]\subseteq p[\NatSeg{m+1}]\cup q[\NatSeg{n+1}]}{%
      \FormulaRefAuto{A\subseteq A\cup B}[thm]}
  \proofstepwidestar[2,4,6,8,10]{0\leq m\leftrightarrow 0\in\NatSeg{m+1}}{%
      \FormulaRefAuto{PeanoNatSegLeqAddOneMembership}[thm]{%
              21,4}}
  \proofstepwidestar[2,4,6,8,10]{0\in\NatSeg{m+1}}{%
      \FormulaRefAuto{P\leftrightarrow Q\dsep P\vdash Q}[thm]{%
            27,
            25}}
  \proofstepwidestar[2,4,6,8,10]{p(0)\in p[\NatSeg{m+1}]}{%
      \FormulaRefAuto{C\subseteq A\dsep F\colon A\to B\dsep x\in C\vdash F(x)\in F[C]}[thm]{%
          24,
          23,
          28}}
  \proofstepwidestar[2,4,6,8,10]{p(0)\in p[\NatSeg{m+1}]\cup q[\NatSeg{n+1}]}{%
      \FormulaRefAuto{A\subseteq B\dsep x\in A\vdash x\in B}[thm]{%
        26,
        29}}
  \proofstepwidestarbreak[2,4,6,8,10]{%
    p[\NatSeg{m+1}]\cup q[\NatSeg{n+1}]\neq\varnothing}{%
    \FormulaRefAuto{a\in A\vdash A\neq\varnothing}[thm]{%
      30}}
  \proofstepwidestarbreak[32]{%
    z\in p[\NatSeg{m+1}]\cup q[\NatSeg{n+1}]}{\rA}
  \proofstepwidestar[1,2,3,4,5,6,7,8,9,10,11,12,13,32]{\GraphStruct { V , E }}{%
      \FormulaRefAuto{UndirectedGraphAccess}[thm]{1}}
  \proofstepwidestar[1,2,3,4,5,6,7,8,9,10,11,12,13,32]{\GraphStruct { V , E }}{%
      \FormulaRefAuto{UndirectedGraphAccess}[thm]{1}}
  \proofstepwidestar[1,2,3,4,5,6,7,8,9,10,11,12,13,32]{\GraphStruct { V , E }}{%
      \FormulaRefAuto{UndirectedGraphAccess}[thm]{1}}
  \proofstepwidestarbreak[1,2,3,4,5,6,7,8,9,10,11,12,13,32]{%
    \exists w\in p[\NatSeg{m+1}]\cup q[\NatSeg{n+1}]\,
    \exists h\in\mathbb N\,\exists k\in\mathbb N\,
    \exists s\in\powerset(\mathbb N\times V)\,
    \exists t\in\powerset(\mathbb N\times V)\,\Bigl(
      z\neq w\land
      \DirectedPath{V}{E}{s}{h}\land s(0)=z\land s(h)=w\land
      \DirectedPath{V}{E}{t}{k}\land t(0)=z\land t(k)=w\land
      (h,s)\neq(k,t)\land
      s[\NatSeg{h+1}]\cap t[\NatSeg{k+1}]=\{z,w\}\land
      s[\NatSeg{h+1}]\cup t[\NatSeg{k+1}]
        \subseteq p[\NatSeg{m+1}]\cup q[\NatSeg{n+1}]\Bigr)}{\text{\parbox[t]{\linewidth}{\raggedright \hyperref[proof:internally-disjoint-cycle-node]{Knotenbedingung oben}}}}
  \proofstepwidestarbreak[1,2,3,4,5,6,7,8,9,10,11,12,13]{%
    \forall z\in p[\NatSeg{m+1}]\cup q[\NatSeg{n+1}]\,
    \exists w\in p[\NatSeg{m+1}]\cup q[\NatSeg{n+1}]\,
    \exists h\in\mathbb N\,\exists k\in\mathbb N\,
    \exists s\in\powerset(\mathbb N\times V)\,
    \exists t\in\powerset(\mathbb N\times V)\,\Bigl(
      z\neq w\land
      \DirectedPath{V}{E}{s}{h}\land s(0)=z\land s(h)=w\land
      \DirectedPath{V}{E}{t}{k}\land t(0)=z\land t(k)=w\land
      (h,s)\neq(k,t)\land
      s[\NatSeg{h+1}]\cap t[\NatSeg{k+1}]=\{z,w\}\land
      s[\NatSeg{h+1}]\cup t[\NatSeg{k+1}]
        \subseteq p[\NatSeg{m+1}]\cup q[\NatSeg{n+1}]\Bigr)}{%
    \rUI{\rRI{32,36}}}
  \proofstepwidestarbreak[1,2,3,4,5,6,7,8,9,10,11,12,13]{%
    \FinitePathCycleCore{V}{E}{%
      p[\NatSeg{m+1}]\cup q[\NatSeg{n+1}]}}{%
    \rAI{20,31,37}}
\end{tabproofwide}

% Allgemeine Regeln nach den elementaren Pfadoperationen aus Band 26.
\subsection{Zugriff auf gespeicherte Wege und Erreichbarkeit}

\Needspace{10\baselineskip}
\FormulaThmDeltaK[Index- und Knotendaten eines Pfadschritts]{\begin{aligned}[t]
    &\GraphStruct{V,E}\dsep\DirectedPath{V}{E}{p}{n}\dsep i\in\NatSeg n\vdash\\[-2pt]
    &{} i\in\mathbb N\land i\in\NatSeg{n+1}\land i+1\in\NatSeg{n+1}\\[-2pt]
    &\land{} p(i)\in V\land p(i+1)\in V\\[-2pt]
    &\land(p(i),p(i+1))\in E
    \end{aligned}}{B27PathStepData}{\DeltaRow{Mengen}{V\dsep E\dsep p\dsep n\dsep i}}
\begin{tabproofwide}
  \proofstepwidestar[1]{\GraphStruct{V,E}}{\rA}
  \proofstepwidestar[2]{\DirectedPath{V}{E}{p}{n}}{\rA}
  \proofstepwidestar[3]{i\in\NatSeg n}{\rA}
  \proofstepwidestar[1,2]{\begin{aligned}[t]
    &n\in\mathbb N\land p\colon\NatSeg{n+1}\to V\\[-2pt]
    &\land p\colon\NatSeg{n+1}\inj V
    \end{aligned}}{\FormulaRefAuto{FinitePathTyping}[thm]{1,2}}
  \proofstepwidestar[1,2]{n\in\mathbb N}{\rAEa{4}}
  \proofstepwidestar[1,2]{p\colon\NatSeg{n+1}\to V}{\rAEa{\rAEb{4}}}
  \proofstepwidestar[1,2,3]{i\in\mathbb N}{\FormulaRefAuto{PeanoNatSegElementNatural}[thm]{5,3}}
  \proofstepwidestar[1,2]{\NatSeg n\subseteq\NatSeg{n+1}}{\FormulaRefAuto{PeanoNatSegSubsetAddOne}[thm]{5}}
  \proofstepwidestar[1,2,3]{i\in\NatSeg{n+1}}{\FormulaRefAuto{A\subseteq B\dsep x\in A\vdash x\in B}[thm]{8,3}}
  \proofstepwidestar[1,2,3]{i+1\in\NatSeg{n+1}}{\FormulaRefAuto{FinitePathIndexShiftOne}[thm]{5,3}}
  \proofstepwidestar[1,2,3]{p(i)\in V}{\FormulaRefAuto{F\colon A\to B\dsep x\in A\vdash F(x)\in B}[thm]{6,9}}
  \proofstepwidestar[1,2,3]{p(i+1)\in V}{\FormulaRefAuto{F\colon A\to B\dsep x\in A\vdash F(x)\in B}[thm]{6,10}}
  \proofstepwidestar[1,2]{\DirectedWalk{V}{E}{p}{n}}{\FormulaRefAuto{DirectedPathWalkAxiom}[ax]{1,2}}
  \proofstepwidestar[1,2,3]{(p(i),p(i+1))\in E}{\FormulaRefAuto{DirectedWalkStepEdge}[thm]{1,13,3}}
  \proofstepwidestar[1,2,3]{\begin{aligned}[t]
    &i\in\mathbb N\land i\in\NatSeg{n+1}\land i+1\in\NatSeg{n+1}\\[-2pt]
    &\land{}p(i)\in V\land p(i+1)\in V\\[-2pt]
    &\land(p(i),p(i+1))\in E
    \end{aligned}}{\rAIStar{7,9,10,11,12,14}}
\end{tabproofwide}

\Needspace{10\baselineskip}
\FormulaThmDeltaK[Einführung eines gespeicherten Pfades]{%
  \begin{aligned}[t]
    &\GraphStruct{V,E}\dsep n\in\mathbb N\dsep
      p\in\powerset(\mathbb N\times V)\dsep{}\\[-2pt]
    &\DirectedPath V E p n\dsep p(0)=u\dsep p(n)=v\vdash
      (n,p)\in\GraphPathsBetween V E u v
  \end{aligned}
}{B27SimplePathFamilyPairIntroduction}{%
  \DeltaRow{Mengen}{V\dsep E\dsep n\dsep p\dsep u\dsep v\dsep m\dsep s}
}
\begin{tabproofwide}
  \proofstepwidestar[1]{\GraphStruct{V,E}}{\rA}
  \proofstepwidestar[2]{n\in\mathbb N}{\rA}
  \proofstepwidestar[3]{p\in\powerset(\mathbb N\times V)}{\rA}
  \proofstepwidestar[4]{\DirectedPath V E p n}{\rA}
  \proofstepwidestar[5]{p(0)=u}{\rA}
  \proofstepwidestar[6]{p(n)=v}{\rA}
  \proofstepwidestar[]{(n,p)=(n,p)}{\rII}
  \proofstepwidestar[4,5,6]{\begin{aligned}[t]
    &(n,p)=(n,p)\land\DirectedPath V E p n\\[-2pt]
    &\land p(0)=u\land p(n)=v
    \end{aligned}}{%
    \rAI{7,\rAI{4,\rAI{5,6}}}}
  \proofstepwidestar[2,3,4,5,6]{\begin{aligned}[t]
    &\exists m\in\mathbb N\,\exists s\in\powerset(\mathbb N\times V)\,\\[-2pt]
    &\Bigl( \quad(n,p)=(m,s)\land\DirectedPath V E s m\\[-2pt]
    &\land s(0)=u\land s(m)=v\Bigr)
    \end{aligned}}{\rEIStar{2,3,8}}
  \proofstepwidestar[1]{\begin{aligned}[t]
    &(n,p)\in\GraphPathsBetween V E u v\leftrightarrow\\[-2pt]
    &{} \exists m\in\mathbb N\,\exists s\in\powerset(\mathbb N\times V)\,\\[-2pt]
    &\Bigl( \quad(n,p)=(m,s)\land\DirectedPath V E s m\\[-2pt]
    &\land s(0)=u\land s(m)=v\Bigr)
    \end{aligned}}{\FormulaRefAuto{FiniteSimplePathFamilyMembership}[thm]{1}}
  \proofstepwidestar[1,2,3,4,5,6]{(n,p)\in\GraphPathsBetween V E u v}{%
    \rRE{9,\rLREb{10}}}
\end{tabproofwide}

\Needspace{10\baselineskip}
\FormulaThmDeltaK[Die Bestandteile eines gespeicherten Pfades]{\begin{aligned}[t]
    &\GraphStruct{V,E}\dsep(n,p)\in\GraphPathsBetween V E u v\vdash\\[-2pt]
    &{} n\in\mathbb N\land p\in\powerset(\mathbb N\times V)\\[-2pt]
    &\land \DirectedPath V E p n\\[-2pt]
    &\land p(0)=u\land p(n)=v
    \end{aligned}}{B27SimplePathFamilyPairElimination}{%
  \DeltaRow{Mengen}{V\dsep E\dsep n\dsep p\dsep u\dsep v\dsep m\dsep s}
}
\begin{tabproofwide}
  \proofstepwidestar[1]{\GraphStruct{V,E}}{\rA}
  \proofstepwidestar[2]{(n,p)\in\GraphPathsBetween V E u v}{\rA}
  \proofstepwidestar[1]{\begin{aligned}[t]
    &(n,p)\in\GraphPathsBetween V E u v\leftrightarrow\\[-2pt]
    &{} \exists m\in\mathbb N\,\exists s\in\powerset(\mathbb N\times V)\,\\[-2pt]
    &\Bigl( \quad(n,p)=(m,s)\land\DirectedPath V E s m\\[-2pt]
    &\land s(0)=u\land s(m)=v\Bigr)
    \end{aligned}}{\FormulaRefAuto{FiniteSimplePathFamilyMembership}[thm]{1}}
  \proofstepwidestar[1,2]{\begin{aligned}[t]
    &\exists m\in\mathbb N\,\exists s\in\powerset(\mathbb N\times V)\,\\[-2pt]
    &\Bigl( \quad(n,p)=(m,s)\land\DirectedPath V E s m\\[-2pt]
    &\land s(0)=u\land s(m)=v\Bigr)
    \end{aligned}}{\rRE{2,\rLREa{3}}}
  \proofstepwidestar[5]{\begin{aligned}[t]
    &m\in\mathbb N\land s\in\powerset(\mathbb N\times V)\\[-2pt]
    &\land{} (n,p)=(m,s)\land\DirectedPath V E s m\\[-2pt]
    &\land s(0)=u\land s(m)=v
    \end{aligned}}{\rA}
  \proofstepwidestar[5]{(n,p)=(m,s)}{\rAEa{\rAEb{\rAEb{5}}}}
  \proofstepwidestar[]{(n,p)=(m,s)\leftrightarrow(n=m\land p=s)}{%
    \FormulaRefAuto{(a,b)=(c,d)\eqvdash a=c\land b=d}[thm]}
  \proofstepwidestar[5]{n=m\land p=s}{\rRE{6,\rLREa{7}}}
  \proofstepwidestar[5]{m=n}{\FormulaRefAuto{a=b\vdash b=a}[thm]{\rAEa{8}}}
  \proofstepwidestar[5]{s=p}{\FormulaRefAuto{a=b\vdash b=a}[thm]{\rAEb{8}}}
  \proofstepwidestar[5]{\begin{aligned}[t]
    &m\in\mathbb N\land s\in\powerset(\mathbb N\times V)\\[-2pt]
    &\land{} \DirectedPath V E s m\\[-2pt]
    &\land s(0)=u\land s(m)=v
    \end{aligned}}{%
    \rAI{\rAEa{5},\rAI{\rAEa{\rAEb{5}},\rAEb{\rAEb{\rAEb{5}}}}}}
  \proofstepwidestar[5]{\begin{aligned}[t]
    &n\in\mathbb N\land p\in\powerset(\mathbb N\times V)\\[-2pt]
    &\land{} \DirectedPath V E p n\\[-2pt]
    &\land p(0)=u\land p(n)=v
    \end{aligned}}{\rIE{9,10,11}}
  \proofstepwidestar[1,2]{\begin{aligned}[t]
    &n\in\mathbb N\land p\in\powerset(\mathbb N\times V)\\[-2pt]
    &\land{} \DirectedPath V E p n\\[-2pt]
    &\land p(0)=u\land p(n)=v
    \end{aligned}}{\rEEStar{4,5,12}}
\end{tabproofwide}

\Needspace{10\baselineskip}
\FormulaThmDeltaK[Der Weg aus einem einzigen Knoten]{%
  \GraphStruct{V,E}\dsep r\in V\vdash
  \exists z\,(z\in\GraphPathsBetween V E r r)
}{B27RootPathSingleton}{\DeltaRow{Mengen}{V\dsep E\dsep r\dsep i\dsep j\dsep z}
\DeltaRow{Abkürzung}{p=\{(0,r)\}}}
\begin{tabproofwide}
  \proofstepwidestar[1]{\GraphStruct{V,E}}{\rA}
  \proofstepwidestar[2]{r\in V}{\rA}
  \proofstepwidestar[]{0\in\mathbb N}{\FormulaRefAuto{PeanoZero}[ax]}
  \proofstepwidestar[2]{p\colon\{0\}\to V}{\FormulaRefAuto{SingletonGraphFunction}[thm]{2}}
  \proofstepwidestar[]{\NatSeg1=\{0\}}{\FormulaRefAuto{PeanoNatSegOneEquality}[thm]}
  \proofstepwidestar[]{0+1=1}{\FormulaRefAuto{PeanoAddLeftZero}[thm]{\FormulaRefAuto{PeanoOneInN}[thm]}}
  \proofstepwidestar[2]{p\colon\NatSeg{0+1}\to V}{%
    \rIE{\FormulaRefAuto{a=b\vdash b=a}[thm]{5},\FormulaRefAuto{a=b\vdash b=a}[thm]{6},4}}
  \proofstepwidestar[8]{i\in\NatSeg{0+1}}{\rA}
  \proofstepwidestar[9]{j\in\NatSeg{0+1}}{\rA}
  \proofstepwidestar[8]{i=0}{\FormulaRefAuto{x\in\{a\}\eqvdash x=a}[thm]{\rIE{6,5,8}}}
  \proofstepwidestar[9]{j=0}{\FormulaRefAuto{x\in\{a\}\eqvdash x=a}[thm]{\rIE{6,5,9}}}
  \proofstepwidestar[8,9]{i=j}{\FormulaRefAuto{a=b\dsep c=b\vdash a=c}[thm]{10,11}}
  \proofstepwidestar[]{\begin{aligned}[t]
    &\forall i\in\NatSeg{0+1}\,\forall j\in\NatSeg{0+1}\,\\[-2pt]
    &(p(i)=p(j)\rightarrow i=j)
    \end{aligned}}{\rUI{\rRI{8,\rUI{\rRI{9,\FormulaRefAuto{Q\vdash P\rightarrow Q}[thm]{12}}}}}}
  \proofstepwidestar[2]{p\colon\NatSeg{0+1}\inj V}{\FormulaRefAuto{Injektive Funktion}[def]{7,13}}
  \proofstepwidestar[15]{i\in\NatSeg0}{\rA}
  \proofstepwidestar[15]{\bot}{\rBI{\rIE{\FormulaRefAuto{PeanoNatSegZero}[thm],15},\FormulaRefAuto{x\not\in\varnothing}[thm]}}
  \proofstepwidestar[15]{(p(i),p(i+1))\in E}{\FormulaRefAuto{\bot\vdash P}[thm]{16}}
  \proofstepwidestar[]{\forall i\in\NatSeg0\,(p(i),p(i+1))\in E}{\rUI{\rRI{15,17}}}
  \proofstepwidestar[1,2]{\DirectedWalk V E p0}{\FormulaRefAuto{DirectedWalkDef}[def]{1,3,7,18}}
  \proofstepwidestar[1,2]{\DirectedPath V E p0}{\FormulaRefAuto{DirectedPathDef}[def]{1,19,14}}
  \proofstepwidestar[2]{p(0)=r}{\FormulaRefAuto{SingletonGraphValue}[thm]{2}}
  \proofstepwidestar[1,2]{p\in\powerset(\mathbb N\times V)}{\FormulaRefAuto{FinitePathPowerSetTyping}[thm]{1,20}}
  \proofstepwidestar[1,2]{(0,p)\in\GraphPathsBetween V E r r}{\FormulaRefAuto{B27SimplePathFamilyPairIntroduction}[thm]{1,3,22,20,21,21}}
  \proofstepwidestar[1,2]{\exists z\,(z\in\GraphPathsBetween V E r r)}{\rEI{23}}
\end{tabproofwide}

\Needspace{10\baselineskip}
\FormulaThmDeltaK[Ein erreichter Knoten führt über eine Kante weiter]{%
  \begin{aligned}[t]&\GraphStruct{V,E}\dsep
  \exists z\,(z\in\GraphPathsBetween V E r u)\dsep(u,v)\in E\vdash{}\\[-2pt]
  &\exists z\,(z\in\GraphPathsBetween V E r v)\end{aligned}
}{B27RootPathExtendExistence}{\DeltaRow{Mengen}{V\dsep E\dsep r\dsep u\dsep v\dsep z\dsep n\dsep p\dsep m\dsep q}}
\begin{tabproofwide}
  \proofstepwidestar[1]{\GraphStruct{V,E}}{\rA}
  \proofstepwidestar[2]{\exists z\,(z\in\GraphPathsBetween V E r u)}{\rA}
  \proofstepwidestar[3]{(u,v)\in E}{\rA}
  \proofstepwidestar[4]{z\in\GraphPathsBetween V E r u}{\rA}
  \proofstepwidestar[1]{\begin{aligned}[t]
    &z\in\GraphPathsBetween V E r u\leftrightarrow\\[-2pt]
    &\exists n\in\mathbb N\,\exists p\in\powerset(\mathbb N\times V)\\[-2pt]
    &\,(z=(n,p)\land\DirectedPath V E p n\\[-2pt]
    &\land p(0)=r\land p(n)=u)
    \end{aligned}}{\FormulaRefAuto{FiniteSimplePathFamilyMembership}[thm]{1}}
  \proofstepwidestar[1,4]{\begin{aligned}[t]
    &\exists n\in\mathbb N\,\exists p\in\powerset(\mathbb N\times V)\\[-2pt]
    &\,(z=(n,p)\land\DirectedPath V E p n\\[-2pt]
    &\land p(0)=r\land p(n)=u)
    \end{aligned}}{\rRE{4,\rLREa{5}}}
  \proofstepwidestar[7]{\begin{aligned}[t]
    &n\in\mathbb N\land p\in\powerset(\mathbb N\times V)\\[-2pt]
    &\land z=(n,p)\land\DirectedPath V E p n\\[-2pt]
    &\land p(0)=r\land p(n)=u
    \end{aligned}}{\rA}
  \proofstepwidestar[7]{\DirectedPath V E p n}{\rAEa{\rAEb{\rAEb{\rAEb{7}}}}}
  \proofstepwidestar[7]{p(0)=r}{\rAEa{\rAEb{\rAEb{\rAEb{\rAEb{7}}}}}}
  \proofstepwidestar[7]{p(n)=u}{\rAEb{\rAEb{\rAEb{\rAEb{\rAEb{7}}}}}}
  \proofstepwidestar[7]{u=p(n)}{\FormulaRefAuto{a=b\vdash b=a}[thm]{10}}
  \proofstepwidestar[3,7]{(p(n),v)\in E}{\rIE{11,3}}
  \proofstepwidestar[1,3,7]{\begin{aligned}[t]
    &\exists m\in\mathbb N\,\exists q\in\powerset(\mathbb N\times V)\\[-2pt]
    &\,(\DirectedPath V E q m\\[-2pt]
    &\land q(0)=p(0)\land q(m)=v)
    \end{aligned}}{\FormulaRefAuto{FinitePathEndpointStep}[thm]{1,8,12}}
  \proofstepwidestar[14]{\begin{aligned}[t]
    &m\in\mathbb N\land q\in\powerset(\mathbb N\times V)\\[-2pt]
    &\land\DirectedPath V E q m\\[-2pt]
    &\land q(0)=p(0)\land q(m)=v
    \end{aligned}}{\rA}
  \proofstepwidestar[7,14]{q(0)=r}{\rIE{9,\rAEa{\rAEb{\rAEb{\rAEb{14}}}}}}
  \proofstepwidestar[1,7,14]{(m,q)\in\GraphPathsBetween V E r v}{\FormulaRefAuto{B27SimplePathFamilyPairIntroduction}[thm]{1,\rAEa{14},\rAEa{\rAEb{14}},\rAEa{\rAEb{\rAEb{14}}},15,\rAEb{\rAEb{\rAEb{\rAEb{14}}}}}}
  \proofstepwidestar[1,7,14]{\exists z\,(z\in\GraphPathsBetween V E r v)}{\rEI{16}}
  \proofstepwidestar[1,3,7]{\exists z\,(z\in\GraphPathsBetween V E r v)}{\rEEStar{13,14,17}}
  \proofstepwidestar[1,3,4]{\exists z\,(z\in\GraphPathsBetween V E r v)}{\rEEStar{6,7,18}}
  \proofstepwidestar[1,2,3]{\exists z\,(z\in\GraphPathsBetween V E r v)}{\rEE{2,4,19}}
\end{tabproofwide}


\chapter{Wälder und Bäume}

\section{Wälder und Wegeindeutigkeit}

\FormulaDefDeltaK[Wald]{%
  \ForestStruct{V,E}%
}{ForestDef}{%
  \DeltaRow{Mengen}{V\dsep E\dsep u\dsep v\dsep a\dsep b}
  \DeltaRow{\textbf{Axiome}}{}
  \DeltaRow{\makecell[l]{Schlichter\\ungerichteter Graph}}{%
    \SimpleUndirectedGraphStruct{V,E}}
  \DeltaRow{Wegeindeutigkeit}{%
    \begin{aligned}[t]
      &\forall u,v\in V\,\forall a,b\,\Bigl(\\[-2pt]
      &\quad a\in\GraphPathsBetween{V}{E}{u}{v}\land{}\\[-2pt]
      &\quad b\in\GraphPathsBetween{V}{E}{u}{v}\rightarrow a=b\Bigr)
    \end{aligned}}
  \DeltaRow{\textbf{Neues Symbol}}{}
  \DeltaPrem{Struktur}{\ForestStruct{V,E}}
}

\FormulaAxiomDeltaK[Schlichter ungerichteter Graph]{%
  \begin{aligned}[t]
    &\ForestStruct{V,E}\vdash{}\\[-2pt]
    &\SimpleUndirectedGraphStruct{V,E}
  \end{aligned}%
}{ForestGraphAxiom}{%
  \DeltaRow{Mengen}{V\dsep E\dsep u\dsep v\dsep a\dsep b}
}

\FormulaAxiomDeltaK[Wegeindeutigkeit]{%
  \begin{aligned}[t]
    &\ForestStruct{V,E}\vdash{}\\[-2pt]
    &\begin{aligned}[t]
      &\forall u,v\in V\,\forall a,b\,\Bigl(\\[-2pt]
      &\quad a\in\GraphPathsBetween{V}{E}{u}{v}\land{}\\[-2pt]
      &\quad b\in\GraphPathsBetween{V}{E}{u}{v}\rightarrow a=b\Bigr)
    \end{aligned}
  \end{aligned}%
}{ForestUniquePathsAxiom}{%
  \DeltaRow{Mengen}{V\dsep E\dsep u\dsep v\dsep a\dsep b}
}

\FormulaThmDeltaK[Charakterisierung durch die Axiome]{%
  \begin{aligned}
  \ForestStruct{V,E}\eqvdash{}&
  \SimpleUndirectedGraphStruct{V,E}\land{}\\[-2pt]
  &\forall u\in V\,\forall v\in V\,\forall a\,\forall b\,\\[-2pt]
  &\qquad\Bigl(a\in\GraphPathsBetween{V}{E}{u}{v}\land
          b\in\GraphPathsBetween{V}{E}{u}{v}\rightarrow a=b\Bigr)
  \end{aligned}%
}{ForestCharacterization}{%
  \DeltaRow{Mengen}{V\dsep E\dsep u\dsep v\dsep a\dsep b}
}
\begin{tabproofsplitwide}
  \proofpartwide{Hinrichtung}
    \proofstepwidestar[1]{\ForestStruct{V,E}}{\rA}
    \proofstepwidestar[1]{\SimpleUndirectedGraphStruct{V,E}}{%
      \FormulaRefAuto{ForestGraphAxiom}[ax]{1}}
    \proofstepwidestar[1]{\begin{aligned}[t]
      &\forall u,v\in V\,\forall a,b\,\Bigl(\\[-2pt]
      &\quad a\in\GraphPathsBetween{V}{E}{u}{v}\land{}\\[-2pt]
      &\quad b\in\GraphPathsBetween{V}{E}{u}{v}\rightarrow a=b\Bigr)
    \end{aligned}}{%
      \FormulaRefAuto{ForestUniquePathsAxiom}[ax]{1}}
    \proofstepwidestar[1]{%
      \begin{aligned}[t]
      &\SimpleUndirectedGraphStruct{V,E}\land{}\\[-2pt]
      &\begin{aligned}[t]
      &\forall u,v\in V\,\forall a,b\,\Bigl(\\[-2pt]
      &\quad a\in\GraphPathsBetween{V}{E}{u}{v}\land{}\\[-2pt]
      &\quad b\in\GraphPathsBetween{V}{E}{u}{v}\rightarrow a=b\Bigr)
    \end{aligned}
    \end{aligned}}{\rAIStar{2,3}}
  \closeproofpartwide

  \proofpartwide{Rückrichtung}
    \proofstepwidestar[1]{%
      \begin{aligned}[t]
      &\SimpleUndirectedGraphStruct{V,E}\land{}\\[-2pt]
      &\begin{aligned}[t]
      &\forall u,v\in V\,\forall a,b\,\Bigl(\\[-2pt]
      &\quad a\in\GraphPathsBetween{V}{E}{u}{v}\land{}\\[-2pt]
      &\quad b\in\GraphPathsBetween{V}{E}{u}{v}\rightarrow a=b\Bigr)
    \end{aligned}
    \end{aligned}}{\rA}
    \proofstepwidestar[1]{\SimpleUndirectedGraphStruct{V,E}}{\rAEa{1}}
    \proofstepwidestar[1]{\begin{aligned}[t]
      &\forall u,v\in V\,\forall a,b\,\Bigl(\\[-2pt]
      &\quad a\in\GraphPathsBetween{V}{E}{u}{v}\land{}\\[-2pt]
      &\quad b\in\GraphPathsBetween{V}{E}{u}{v}\rightarrow a=b\Bigr)
    \end{aligned}}{\rAEb{1}}
    \proofstepwidestar[1]{\ForestStruct{V,E}}{%
      \FormulaRefAuto{ForestDef}[def]{2,3}}
  \closeproofpartwide
\end{tabproofsplitwide}

Die Leermenge ist als Träger eines Waldes zugelassen. Nichtleerheit wird erst
beim Baum gefordert.

\FormulaThmDeltaKR[Zugriff auf die Waldaxiome]{%
  \begin{aligned}[t]
    &\ForestStruct{V,E}\vdash{}\\[-2pt]
    &\text{(i)}\quad\SimpleUndirectedGraphStruct{V,E},\\[-2pt]
    &\text{(ii)}\quad
      \begin{aligned}[t]
        &u,v\in V\dsep
          a\in\GraphPathsBetween{V}{E}{u}{v}\dsep{}\\[-2pt]
        &b\in\GraphPathsBetween{V}{E}{u}{v}\vdash a=b.
      \end{aligned}
  \end{aligned}%
}{%
  \begin{aligned}[t]
  \ForestStruct{V,E}\vdash{}&
  \SimpleUndirectedGraphStruct{V,E}\dsep{}\\[-2pt]
  &u,v\in V\dsep
    a\in\GraphPathsBetween{V}{E}{u}{v}\dsep
    b\in\GraphPathsBetween{V}{E}{u}{v}\vdash a=b
  \end{aligned}%
}{ForestAccess}{%
  \DeltaRow{Mengen}{V\dsep E\dsep u\dsep v\dsep a\dsep b}
}
\begin{tabproofsplitwide}
  \proofpartwide{Schlichte Graphenstruktur}
  \proofstepwidestarbreak[1]{\ForestStruct{V,E}}{\rA}
  \proofstepwidestarbreak[1]{\SimpleUndirectedGraphStruct{V,E}}{%
    \FormulaRefAuto{ForestGraphAxiom}[ax]{1}}
  \closeproofpartwide

  \proofpartwide{Eindeutigkeit der Wege}
  \setcounter{proofstepnr}{2}
  \proofstepwidestarbreak[3]{u\in V}{\rA}
  \proofstepwidestarbreak[4]{v\in V}{\rA}
  \proofstepwidestarbreak[5]{a\in\GraphPathsBetween{V}{E}{u}{v}}{\rA}
  \proofstepwidestarbreak[6]{b\in\GraphPathsBetween{V}{E}{u}{v}}{\rA}
  \proofstepwidestarbreak[1,3,4,5,6]{a=b}{%
    \FormulaRefAuto{ForestUniquePathsAxiom}[ax]{1,3,4,5,6}}
  \closeproofpartwide
\end{tabproofsplitwide}

\FormulaDefDeltaK[Wegezyklus und Azyklizität]{%
  \PathAcyclicGraphStruct{V,E}%
}{PathAcyclicGraphDef}{%
  \DeltaRow{Mengen}{V\dsep E\dsep u\dsep v\dsep a\dsep b}
  \DeltaRow{\textbf{Axiome}}{}
  \DeltaRow{\makecell[l]{Schlichter\\ungerichteter Graph}}{%
    \SimpleUndirectedGraphStruct{V,E}}
  \DeltaRow{Wegezyklusfreiheit}{%
    \begin{aligned}[t]
      &\neg\exists u,v\in V\,\exists a,b\,\Bigl(\\[-2pt]
      &\quad a\in\GraphPathsBetween{V}{E}{u}{v}\land{}\\[-2pt]
      &\quad b\in\GraphPathsBetween{V}{E}{u}{v}\land a\neq b\Bigr)
    \end{aligned}}
  \DeltaRow{\textbf{Neues Symbol}}{}
  \DeltaPrem{Struktur}{\PathAcyclicGraphStruct{V,E}}
}

\FormulaAxiomDeltaK[Schlichter ungerichteter Graph]{%
  \begin{aligned}[t]
    &\PathAcyclicGraphStruct{V,E}\vdash{}\\[-2pt]
    &\SimpleUndirectedGraphStruct{V,E}
  \end{aligned}%
}{PathAcyclicGraphGraphAxiom}{%
  \DeltaRow{Mengen}{V\dsep E\dsep u\dsep v\dsep a\dsep b}
}

\FormulaAxiomDeltaK[Wegezyklusfreiheit]{%
  \begin{aligned}[t]
    &\PathAcyclicGraphStruct{V,E}\vdash{}\\[-2pt]
    &\begin{aligned}[t]
      &\neg\exists u,v\in V\,\exists a,b\,\Bigl(\\[-2pt]
      &\quad a\in\GraphPathsBetween{V}{E}{u}{v}\land{}\\[-2pt]
      &\quad b\in\GraphPathsBetween{V}{E}{u}{v}\land a\neq b\Bigr)
    \end{aligned}
  \end{aligned}%
}{PathAcyclicGraphNoCycleAxiom}{%
  \DeltaRow{Mengen}{V\dsep E\dsep u\dsep v\dsep a\dsep b}
}

\FormulaThmDeltaK[Charakterisierung durch die Axiome]{%
  \begin{aligned}
  \PathAcyclicGraphStruct{V,E}\eqvdash{}&
  \SimpleUndirectedGraphStruct{V,E}\land{}\\[-2pt]
  &\neg\exists u\in V\,\exists v\in V\,\exists a\,\exists b\,\\[-2pt]
  &\qquad\Bigl(a\in\GraphPathsBetween{V}{E}{u}{v}\land
          b\in\GraphPathsBetween{V}{E}{u}{v}\land a\neq b\Bigr).
  \end{aligned}%
}{PathAcyclicGraphCharacterization}{%
  \DeltaRow{Mengen}{V\dsep E\dsep u\dsep v\dsep a\dsep b}
}
\begin{tabproofsplitwide}
  \proofpartwide{Hinrichtung}
    \proofstepwidestar[1]{\PathAcyclicGraphStruct{V,E}}{\rA}
    \proofstepwidestar[1]{\SimpleUndirectedGraphStruct{V,E}}{%
      \FormulaRefAuto{PathAcyclicGraphGraphAxiom}[ax]{1}}
    \proofstepwidestar[1]{\begin{aligned}[t]
      &\neg\exists u,v\in V\,\exists a,b\,\Bigl(\\[-2pt]
      &\quad a\in\GraphPathsBetween{V}{E}{u}{v}\land{}\\[-2pt]
      &\quad b\in\GraphPathsBetween{V}{E}{u}{v}\land a\neq b\Bigr)
    \end{aligned}}{%
      \FormulaRefAuto{PathAcyclicGraphNoCycleAxiom}[ax]{1}}
    \proofstepwidestar[1]{%
      \begin{aligned}[t]
      &\SimpleUndirectedGraphStruct{V,E}\land{}\\[-2pt]
      &\begin{aligned}[t]
      &\neg\exists u,v\in V\,\exists a,b\,\Bigl(\\[-2pt]
      &\quad a\in\GraphPathsBetween{V}{E}{u}{v}\land{}\\[-2pt]
      &\quad b\in\GraphPathsBetween{V}{E}{u}{v}\land a\neq b\Bigr)
    \end{aligned}
    \end{aligned}}{\rAIStar{2,3}}
  \closeproofpartwide

  \proofpartwide{Rückrichtung}
    \proofstepwidestar[1]{%
      \begin{aligned}[t]
      &\SimpleUndirectedGraphStruct{V,E}\land{}\\[-2pt]
      &\begin{aligned}[t]
      &\neg\exists u,v\in V\,\exists a,b\,\Bigl(\\[-2pt]
      &\quad a\in\GraphPathsBetween{V}{E}{u}{v}\land{}\\[-2pt]
      &\quad b\in\GraphPathsBetween{V}{E}{u}{v}\land a\neq b\Bigr)
    \end{aligned}
    \end{aligned}}{\rA}
    \proofstepwidestar[1]{\SimpleUndirectedGraphStruct{V,E}}{\rAEa{1}}
    \proofstepwidestar[1]{\begin{aligned}[t]
      &\neg\exists u,v\in V\,\exists a,b\,\Bigl(\\[-2pt]
      &\quad a\in\GraphPathsBetween{V}{E}{u}{v}\land{}\\[-2pt]
      &\quad b\in\GraphPathsBetween{V}{E}{u}{v}\land a\neq b\Bigr)
    \end{aligned}}{\rAEb{1}}
    \proofstepwidestar[1]{\PathAcyclicGraphStruct{V,E}}{%
      \FormulaRefAuto{PathAcyclicGraphDef}[def]{2,3}}
  \closeproofpartwide
\end{tabproofsplitwide}

Ein Paar verschiedener einfacher Wege mit denselben Endpunkten heißt hier
ein \textbf{Wegezyklus}. In einem schlichten ungerichteten Graphen ist das
genau die für die Baumtheorie benötigte Form der Azyklizität: Die Vereinigung
zweier solcher Wege enthält einen üblichen einfachen Kreis, während ein
einfacher Kreis zwei verschiedene Wege zwischen zwei seiner Knoten liefert.

\FormulaThmDeltaK[Wälder sind genau die azyklischen schlichten Graphen]{%
  \ForestStruct{V,E}\eqvdash\PathAcyclicGraphStruct{V,E}%
}{ForestIffPathAcyclic}{%
  \DeltaRow{Mengen}{V\dsep E\dsep u\dsep v\dsep a\dsep b}
}
\begin{tabproofsplitwide}
  \proofpartwide{Hinrichtung}
    \proofstepwidestar[1]{\ForestStruct{V,E}}{\rA}
    \proofstepwidestar[1]{\SimpleUndirectedGraphStruct{V,E}}{%
      \FormulaRefAuto{ForestGraphAxiom}[ax]{1}}
    \proofstepwidestar[3]{%
      \begin{aligned}[t]
        &\exists u,v\in V\,\exists a,b\,\Bigl(\\[-2pt]
        &\quad a\in\GraphPathsBetween{V}{E}{u}{v}\land{}\\[-2pt]
        &\quad b\in\GraphPathsBetween{V}{E}{u}{v}\land a\neq b\Bigr)
      \end{aligned}}{\rA}
    \proofstepwidestar[4]{%
      \begin{aligned}[t]
        &u\in V\land v\in V\land{}\\[-2pt]
        &a\in\GraphPathsBetween{V}{E}{u}{v}\land{}\\[-2pt]
        &b\in\GraphPathsBetween{V}{E}{u}{v}\land a\neq b
      \end{aligned}}{\rA}
    \proofstepwidestar[4]{u\in V}{\rAEa{4}}
    \proofstepwidestar[4]{v\in V}{\rAEa{\rAEb{4}}}
    \proofstepwidestar[4]{a\in\GraphPathsBetween{V}{E}{u}{v}}{%
      \rAEa{\rAEb{\rAEb{4}}}}
    \proofstepwidestar[4]{b\in\GraphPathsBetween{V}{E}{u}{v}}{%
      \rAEa{\rAEb{\rAEb{\rAEb{4}}}}}
    \proofstepwidestar[4]{a\neq b}{\rAEb{\rAEb{\rAEb{\rAEb{4}}}}}
    \proofstepwidestar[1,4]{a=b}{%
      \FormulaRefAuto{ForestUniquePathsAxiom}[ax]{1,5,6,7,8}}
    \proofstepwidestar[1,4]{\bot}{\rBI{10,9}}
    \proofstepwidestar[1,3]{\bot}{\rEEStar{3,4,11}}
    \proofstepwidestar[1]{%
      \begin{aligned}[t]
        &\neg\exists u,v\in V\,\exists a,b\,\Bigl(\\[-2pt]
        &\quad a\in\GraphPathsBetween{V}{E}{u}{v}\land{}\\[-2pt]
        &\quad b\in\GraphPathsBetween{V}{E}{u}{v}\land a\neq b\Bigr)
      \end{aligned}}{\rCI{3,12}}
    \proofstepwidestar[1]{\PathAcyclicGraphStruct{V,E}}{%
      \FormulaRefAuto{PathAcyclicGraphDef}[def]{2,13}}
  \closeproofpartwide

  \proofpartwide{Rückrichtung}
    \proofstepwidestar[1]{\PathAcyclicGraphStruct{V,E}}{\rA}
    \proofstepwidestar[1]{\SimpleUndirectedGraphStruct{V,E}}{%
      \FormulaRefAuto{PathAcyclicGraphGraphAxiom}[ax]{1}}
    \proofstepwidestar[1]{%
      \begin{aligned}[t]
        &\neg\exists u,v\in V\,\exists a,b\,\Bigl(\\[-2pt]
        &\quad a\in\GraphPathsBetween{V}{E}{u}{v}\land{}\\[-2pt]
        &\quad b\in\GraphPathsBetween{V}{E}{u}{v}\land a\neq b\Bigr)
      \end{aligned}}{\FormulaRefAuto{PathAcyclicGraphNoCycleAxiom}[ax]{1}}
    \proofstepwidestar[4]{u\in V}{\rA}
    \proofstepwidestar[5]{v\in V}{\rA}
    \proofstepwidestar[6]{%
      \begin{aligned}[t]
        &a\in\GraphPathsBetween{V}{E}{u}{v}\land{}\\[-2pt]
        &b\in\GraphPathsBetween{V}{E}{u}{v}
      \end{aligned}}{\rA}
    \proofstepwidestar[7]{a\neq b}{\rA}
    \proofstepwidestar[4,5,6,7]{%
      \begin{aligned}[t]
        &\exists u,v\in V\,\exists a,b\,\Bigl(\\[-2pt]
        &\quad a\in\GraphPathsBetween{V}{E}{u}{v}\land{}\\[-2pt]
        &\quad b\in\GraphPathsBetween{V}{E}{u}{v}\land a\neq b\Bigr)
      \end{aligned}}{\rEIStar{\rAIStar{4,5,\rAEa{6},\rAEb{6},7}}}
    \proofstepwidestar[1,4,5,6,7]{\bot}{\rBI{8,3}}
    \proofstepwidestar[1,4,5,6]{\neg(a\neq b)}{\rCI{7,9}}
    \proofstepwidestar[1,4,5,6]{a=b}{\FormulaRefAuto{rDN}[thm]{10}}
    \proofstepwidestar[1,4,5]{%
      \begin{aligned}[t]
        &\bigl(a\in\GraphPathsBetween{V}{E}{u}{v}\land{}\\[-2pt]
        &\quad b\in\GraphPathsBetween{V}{E}{u}{v}\bigr)\rightarrow a=b
      \end{aligned}}{\rRI{6,11}}
    \proofstepwidestar[1,4,5]{%
      \begin{aligned}[t]
        &\forall a,b\,\Bigl(a\in\GraphPathsBetween{V}{E}{u}{v}\land{}\\[-2pt]
        &\quad b\in\GraphPathsBetween{V}{E}{u}{v}\rightarrow a=b\Bigr)
      \end{aligned}}{\rUI{12}}
    \proofstepwidestar[1,4]{%
      \begin{aligned}[t]
        &\forall v\in V\,\forall a,b\,\Bigl(\\[-2pt]
        &\quad a\in\GraphPathsBetween{V}{E}{u}{v}\land{}\\[-2pt]
        &\quad b\in\GraphPathsBetween{V}{E}{u}{v}\rightarrow a=b\Bigr)
      \end{aligned}}{\rUI{\rRI{5,13}}}
    \proofstepwidestar[1]{%
      \begin{aligned}[t]
        &\forall u,v\in V\,\forall a,b\,\Bigl(\\[-2pt]
        &\quad a\in\GraphPathsBetween{V}{E}{u}{v}\land{}\\[-2pt]
        &\quad b\in\GraphPathsBetween{V}{E}{u}{v}\rightarrow a=b\Bigr)
      \end{aligned}}{\rUI{\rRI{4,14}}}
    \proofstepwidestar[1]{\ForestStruct{V,E}}{%
      \FormulaRefAuto{ForestDef}[def]{2,15}}
  \closeproofpartwide
\end{tabproofsplitwide}

\section{Bäume}

\FormulaDefDeltaK[Baum]{%
  \TreeStruct{V,E}%
}{TreeDef}{%
  \DeltaRow{Mengen}{V\dsep E\dsep u\dsep v\dsep z}
  \DeltaRow{\textbf{Axiome}}{}
  \DeltaRow{\makecell[l]{Schlichter\\ungerichteter Graph}}{%
    \SimpleUndirectedGraphStruct{V,E}}
  \DeltaRow{Nichtleerer Träger}{%
    V\neq\varnothing}
  \DeltaRow{Eindeutige Verbindungswege}{%
    \forall u,v\in V\,\exists!z\,(z\in\GraphPathsBetween{V}{E}{u}{v})}
  \DeltaRow{\textbf{Neues Symbol}}{}
  \DeltaPrem{Struktur}{\TreeStruct{V,E}}
}

\FormulaAxiomDeltaK[Schlichter ungerichteter Graph]{%
  \begin{aligned}[t]
    &\TreeStruct{V,E}\vdash{}\\[-2pt]
    &\SimpleUndirectedGraphStruct{V,E}
  \end{aligned}%
}{TreeGraphAxiom}{%
  \DeltaRow{Mengen}{V\dsep E\dsep u\dsep v\dsep z}
}

\FormulaAxiomDeltaK[Nichtleerer Träger]{%
  \begin{aligned}[t]
    &\TreeStruct{V,E}\vdash{}\\[-2pt]
    &V\neq\varnothing
  \end{aligned}%
}{TreeNonemptyAxiom}{%
  \DeltaRow{Mengen}{V\dsep E\dsep u\dsep v\dsep z}
}

\FormulaAxiomDeltaK[Eindeutige Verbindungswege]{%
  \begin{aligned}[t]
    &\TreeStruct{V,E}\vdash{}\\[-2pt]
    &\forall u,v\in V\,\exists!z\,(z\in\GraphPathsBetween{V}{E}{u}{v})
  \end{aligned}%
}{TreeUniquePathsAxiom}{%
  \DeltaRow{Mengen}{V\dsep E\dsep u\dsep v\dsep z}
}

\FormulaThmDeltaK[Charakterisierung durch die Axiome]{%
  \begin{aligned}
  \TreeStruct{V,E}\eqvdash{}&
  \SimpleUndirectedGraphStruct{V,E}\land V\neq\varnothing\land{}\\[-2pt]
  &\forall u\in V\,\forall v\in V\,
    \exists!z\,\bigl(z\in\GraphPathsBetween{V}{E}{u}{v}\bigr)
  \end{aligned}%
}{TreeCharacterization}{%
  \DeltaRow{Mengen}{V\dsep E\dsep u\dsep v\dsep z}
}
\begin{tabproofsplitwide}
  \proofpartwide{Hinrichtung}
    \proofstepwidestar[1]{\TreeStruct{V,E}}{\rA}
    \proofstepwidestar[1]{\SimpleUndirectedGraphStruct{V,E}}{%
      \FormulaRefAuto{TreeGraphAxiom}[ax]{1}}
    \proofstepwidestar[1]{V\neq\varnothing}{%
      \FormulaRefAuto{TreeNonemptyAxiom}[ax]{1}}
    \proofstepwidestar[1]{\forall u,v\in V\,\exists!z\,(z\in\GraphPathsBetween{V}{E}{u}{v})}{%
      \FormulaRefAuto{TreeUniquePathsAxiom}[ax]{1}}
    \proofstepwidestar[1]{%
      \begin{aligned}[t]
      &\SimpleUndirectedGraphStruct{V,E}\land{}\\[-2pt]
      &V\neq\varnothing\land{}\\[-2pt]
      &\forall u,v\in V\,\exists!z\,(z\in\GraphPathsBetween{V}{E}{u}{v})
    \end{aligned}}{\rAIStar{2,3,4}}
  \closeproofpartwide

  \proofpartwide{Rückrichtung}
    \proofstepwidestar[1]{%
      \begin{aligned}[t]
      &\SimpleUndirectedGraphStruct{V,E}\land{}\\[-2pt]
      &V\neq\varnothing\land{}\\[-2pt]
      &\forall u,v\in V\,\exists!z\,(z\in\GraphPathsBetween{V}{E}{u}{v})
    \end{aligned}}{\rA}
    \proofstepwidestar[1]{\SimpleUndirectedGraphStruct{V,E}}{\rAEa{1}}
    \proofstepwidestar[1]{V\neq\varnothing}{\rAEa{\rAEb{1}}}
    \proofstepwidestar[1]{\forall u,v\in V\,\exists!z\,(z\in\GraphPathsBetween{V}{E}{u}{v})}{\rAEb{\rAEb{1}}}
    \proofstepwidestar[1]{\TreeStruct{V,E}}{%
      \FormulaRefAuto{TreeDef}[def]{2,3,4}}
  \closeproofpartwide
\end{tabproofsplitwide}

\FormulaThmDeltaKR[Zugriff auf die Baumaxiome]{%
  \begin{aligned}[t]
    &\TreeStruct{V,E}\vdash{}\\[-2pt]
    &\text{(i)}\quad \SimpleUndirectedGraphStruct{V,E},\\[-2pt]
    &\text{(ii)}\quad V\neq\varnothing,\\[-2pt]
    &\text{(iii)}\quad u,v\in V\vdash\exists!z\,\bigl(z\in\GraphPathsBetween{V}{E}{u}{v}\bigr).
  \end{aligned}%
}{%
  \begin{aligned}[t]
  \TreeStruct{V,E}\vdash{}&
    \SimpleUndirectedGraphStruct{V,E}\dsep V\neq\varnothing\dsep{}\\[-2pt]
  &u,v\in V\vdash
    \exists!z\,\bigl(z\in\GraphPathsBetween{V}{E}{u}{v}\bigr)
  \end{aligned}%
}{TreeAccess}{%
  \DeltaRow{Mengen}{V\dsep E\dsep u\dsep v\dsep z}
}
\begin{tabproofsplitwide}
  \proofpartwide{Graphenstruktur und Nichtleerheit}
  \proofstepwidestarbreak[1]{\TreeStruct{V,E}}{\rA}
  \proofstepwidestarbreak[1]{\SimpleUndirectedGraphStruct{V,E}}{%
    \FormulaRefAuto{TreeGraphAxiom}[ax]{1}}
  \proofstepwidestarbreak[1]{V\neq\varnothing}{%
    \FormulaRefAuto{TreeNonemptyAxiom}[ax]{1}}
  \closeproofpartwide

  \proofpartwide{Existenz und Eindeutigkeit der Wege}
  \setcounter{proofstepnr}{3}
  \proofstepwidestarbreak[4]{u\in V}{\rA}
  \proofstepwidestarbreak[5]{v\in V}{\rA}
  \proofstepwidestarbreak[1,4,5]{%
    \exists!z\,\bigl(z\in\GraphPathsBetween{V}{E}{u}{v}\bigr)}{%
    \FormulaRefAuto{TreeUniquePathsAxiom}[ax]{1,4,5}}
  \closeproofpartwide
\end{tabproofsplitwide}

\FormulaThmDeltaK[Wegeindeutigkeit im Baum]{%
  \begin{aligned}[t]
  &\TreeStruct{V,E}\dsep u,v\in V\dsep
    a\in\GraphPathsBetween{V}{E}{u}{v}\dsep
    b\in\GraphPathsBetween{V}{E}{u}{v}\vdash a=b
  \end{aligned}%
}{TreePathUnique}{%
  \DeltaRow{Mengen}{V\dsep E\dsep u\dsep v\dsep a\dsep b}
}
\begin{tabproofwide}
  \proofstepwidestarbreak[1]{\TreeStruct{V,E}}{\rA}
  \proofstepwidestarbreak[2]{u\in V}{\rA}
  \proofstepwidestarbreak[3]{v\in V}{\rA}
  \proofstepwidestarbreak[4]{a\in\GraphPathsBetween{V}{E}{u}{v}}{\rA}
  \proofstepwidestarbreak[5]{b\in\GraphPathsBetween{V}{E}{u}{v}}{\rA}
  \proofstepwidestarbreak[1,2,3]{%
    \exists!z\,\bigl(z\in\GraphPathsBetween{V}{E}{u}{v}\bigr)}{%
    \FormulaRefAuto{TreeAccess}[thm]{1,2,3}}
  \proofstepwidestarbreak[1,2,3,4,5]{a=b}{%
    \FormulaRefAuto{\exists! x\,P(x),\; P(a),\; P(b) \vdash a = b}[thm]{6,4,5}}
\end{tabproofwide}

\FormulaThmDeltaK[Jeder Baum ist ein nichtleerer Wald]{%
  \TreeStruct{V,E}\vdash\ForestStruct{V,E}\land V\neq\varnothing%
}{TreeIsNonemptyForest}{%
  \DeltaRow{Mengen}{V\dsep E\dsep u\dsep v\dsep a\dsep b}
}
\begin{tabproofsplitwide}
  \proofpartwide{Graphenstruktur und Nichtleerheit}
  \proofstepwidestarbreak[1]{\TreeStruct{V,E}}{\rA}
  \proofstepwidestarbreak[1]{\SimpleUndirectedGraphStruct{V,E}}{%
    \FormulaRefAuto{TreeAccess}[thm]{1}}
  \proofstepwidestarbreak[1]{V\neq\varnothing}{%
    \FormulaRefAuto{TreeAccess}[thm]{1}}
  \closeproofpartwide

  \proofpartwide{Waldeigenschaft und Schluss}
  \setcounter{proofstepnr}{3}
  \proofstepwidestarbreak[4]{u\in V}{\rA}
  \proofstepwidestarbreak[5]{v\in V}{\rA}
  \proofstepwidestarbreak[6]{a\in\GraphPathsBetween{V}{E}{u}{v}}{\rA}
  \proofstepwidestarbreak[7]{b\in\GraphPathsBetween{V}{E}{u}{v}}{\rA}
  \proofstepwidestarbreak[1,4,5,6,7]{a=b}{%
    \FormulaRefAuto{TreePathUnique}[thm]{1,4,5,6,7}}
  \proofstepwidestarbreak[1]{\ForestStruct{V,E}}{%
    \FormulaRefAuto{ForestCharacterization}[thm]{2,4,5,6,7,8}}
  \proofstepwidestarbreak[1]{\ForestStruct{V,E}\land V\neq\varnothing}{%
    \rAI{9,3}}
  \closeproofpartwide
\end{tabproofsplitwide}

\FormulaThmDeltaK[Bäume sind zusammenhängend]{%
  \TreeStruct{V,E}\vdash\ConnectedGraphStruct{V,E}%
}{TreeConnected}{%
  \DeltaRow{Mengen}{V\dsep E\dsep u\dsep v\dsep z\dsep n\dsep p}
}
\begin{tabproofwide}
  \proofstepwidestarbreak[1]{\TreeStruct{V,E}}{\rA}
  \proofstepwidestarbreak[1]{\SimpleUndirectedGraphStruct{V,E}}{%
    \FormulaRefAuto{TreeAccess}[thm]{1}}
  \proofstepwidestar[1]{\UndirectedGraphStruct { V , E } \land \GraphStruct { V , E }}{%
      \FormulaRefAuto{SimpleUndirectedGraphAccess}[thm]{2}}
  \proofstepwidestarbreak[1]{\UndirectedGraphStruct{V,E}}{%
    \rAEa{3}}
  \proofstepwidestarbreak[5]{u\in V}{\rA}
  \proofstepwidestarbreak[6]{v\in V}{\rA}
  \proofstepwidestarbreak[1,5,6]{%
    \exists!z\,\bigl(z\in\GraphPathsBetween{V}{E}{u}{v}\bigr)}{%
    \FormulaRefAuto{TreeAccess}[thm]{1,5,6}}
  \proofstepwidestarbreak[1,5,6]{%
    \exists z\,\bigl(z\in\GraphPathsBetween{V}{E}{u}{v}\bigr)}{%
    \FormulaRefAuto{\exists! x\,P(x)\vdash \exists x\,P(x)}[thm]{7}}
  \proofstepwidestarbreak[9]{z\in\GraphPathsBetween{V}{E}{u}{v}}{\rA}
  \proofstepwidestarbreak[9]{%
    \exists n\in\mathbb N\,\exists p\in\powerset(\mathbb N\times V)\,
    \bigl(z=(n,p)\land\DirectedPath{V}{E}{p}{n}\land
      p(0)=u\land p(n)=v\bigr)}{%
    \FormulaRefAuto{FiniteSimplePathFamilyMembership}[thm]{9}}
  \proofstepwidestarbreak[11]{n\in\mathbb N\land
    \DirectedPath{V}{E}{p}{n}\land p(0)=u\land p(n)=v}{\rA}
  \proofstepwidestar[11]{\UndirectedGraphStruct { V , E } \land \GraphStruct { V , E }}{%
      \FormulaRefAuto{SimpleUndirectedGraphAccess}[thm]{2}}
  \proofstepwidestarbreak[11]{\DirectedWalk{V}{E}{p}{n}}{%
    \FormulaRefAuto{DirectedPathWalkAxiom}[ax]{%
      \rAEb{12},
      \rAEn{11}}}
  \proofstepwidestar[11]{\UndirectedGraphStruct { V , E } \land \GraphStruct { V , E }}{%
      \FormulaRefAuto{SimpleUndirectedGraphAccess}[thm]{2}}
  \proofstepwidestarbreak[11]{u\mathrel{\GraphReach{E}}v}{%
    \FormulaRefAuto{DirectedReachabilityDef}[def]{%
      \rAEb{14},
      5,6,11,13}}
  \proofstepwidestarbreak[1,5,6]{u\mathrel{\GraphReach{E}}v}{%
    \rEE{8,9,\rEE{10,11,15}}}
  \proofstepwidestarbreak[1]{%
    \forall u\in V\,\forall v\in V\,
      u\mathrel{\GraphReach{E}}v}{\rUI{\rRI{5,\rUI{\rRI{6,16}}}}}
  \proofstepwidestarbreak[1]{\ConnectedGraphStruct{V,E}}{%
    \FormulaRefAuto{ConnectedGraphDef}[def]{4,17}}
\end{tabproofwide}

\FormulaThmDeltaK[Bäume sind zusammenhängend und azyklisch]{%
  \TreeStruct{V,E}\vdash
    \ConnectedGraphStruct{V,E}\land\PathAcyclicGraphStruct{V,E}%
}{TreeConnectedAcyclic}{%
  \DeltaRow{Mengen}{V\dsep E}
}
\begin{tabproofwide}
  \proofstepwidestarbreak[1]{\TreeStruct{V,E}}{\rA}
  \proofstepwidestarbreak[1]{\ConnectedGraphStruct{V,E}}{%
    \FormulaRefAuto{TreeConnected}[thm]{1}}
  \proofstepwidestar[1]{\ForestStruct { V , E } \land V \neq \varnothing}{%
      \FormulaRefAuto{TreeIsNonemptyForest}[thm]{1}}
  \proofstepwidestarbreak[1]{\ForestStruct{V,E}}{%
    \rAEa{3}}
  \proofstepwidestarbreak[1]{\PathAcyclicGraphStruct{V,E}}{%
    \FormulaRefAuto{ForestIffPathAcyclic}[thm]{4}}
  \proofstepwidestarbreak[1]{%
    \ConnectedGraphStruct{V,E}\land\PathAcyclicGraphStruct{V,E}}{%
    \rAI{2,5}}
\end{tabproofwide}

\section{Übertragung eindeutiger Wurzelwege}

Für konkrete Bäume ist es oft wesentlich kürzer, nur die Wege von einem
festen Knoten aus zu untersuchen. Die beiden folgenden endlichen
Wegargumente bereiten dieses Kriterium vor.

\FormulaThmDeltaK[Wurzelwege erzeugen Wege zwischen beliebigen Endpunkten]{%
  \begin{aligned}[t]
  &\UndirectedGraphStruct{V,E}\dsep r,u,v\in V\dsep
    a\in\GraphPathsBetween{V}{E}{r}{u}\dsep
    b\in\GraphPathsBetween{V}{E}{r}{v}\vdash{}\\[-2pt]
  &\qquad\exists z\,\bigl(z\in\GraphPathsBetween{V}{E}{u}{v}\bigr)
  \end{aligned}%
}{RootPathsGivePairPath}{%
  \DeltaRow{Mengen}{V\dsep E\dsep r\dsep u\dsep v\dsep a\dsep b\dsep z}
}

Die Umkehrung des ersten Wurzelwegs liefert einen Pfad von \(u\) nach \(r\).
Der kantenweise Transport dieses Pfads entlang des zweiten Wurzelwegs
schneidet jede entstehende Wiederholung sofort ab und endet daher mit einem
einfachen Pfad von \(u\) nach \(v\).
\begin{tabproofwide}
  \proofstepwidestarbreak[1]{\UndirectedGraphStruct{V,E}}{\rA}
  \proofstepwidestarbreak[2]{r\in V}{\rA}
  \proofstepwidestarbreak[3]{u\in V}{\rA}
  \proofstepwidestarbreak[4]{v\in V}{\rA}
  \proofstepwidestarbreak[5]{a\in\GraphPathsBetween{V}{E}{r}{u}}{\rA}
  \proofstepwidestarbreak[6]{b\in\GraphPathsBetween{V}{E}{r}{v}}{\rA}
  \proofstepwidestarbreak[5]{%
    \exists m\in\mathbb N\,\exists p\in\powerset(\mathbb N\times V)\,
    \bigl(a=(m,p)\land\DirectedPath{V}{E}{p}{m}\land
      p(0)=r\land p(m)=u\bigr)}{%
    \FormulaRefAuto{FiniteSimplePathFamilyMembership}[thm]{5}}
  \proofstepwidestarbreak[6]{%
    \exists n\in\mathbb N\,\exists q\in\powerset(\mathbb N\times V)\,
    \bigl(b=(n,q)\land\DirectedPath{V}{E}{q}{n}\land
      q(0)=r\land q(n)=v\bigr)}{%
    \FormulaRefAuto{FiniteSimplePathFamilyMembership}[thm]{6}}
  \proofstepwidestarbreak[9]{m\in\mathbb N\land
    p\in\powerset(\mathbb N\times V)\land\DirectedPath{V}{E}{p}{m}\land
    p(0)=r\land p(m)=u}{\rA}
  \proofstepwidestarbreak[10]{n\in\mathbb N\land
    q\in\powerset(\mathbb N\times V)\land\DirectedPath{V}{E}{q}{n}\land
    q(0)=r\land q(n)=v}{\rA}
  \proofstepwidestarbreak[1,9]{%
    \exists s\in\powerset(\mathbb N\times V)\,
      \bigl(\DirectedPath{V}{E}{s}{m}\land s(0)=u\land s(m)=r\bigr)}{%
    \FormulaRefAuto{UndirectedFinitePathReversal}[thm]{1,\rAEn{9}}}
  \proofstepwidestarbreak[12]{s\in\powerset(\mathbb N\times V)\land
    \DirectedPath{V}{E}{s}{m}\land s(0)=u\land s(m)=r}{\rA}
  \proofstepwidestar[1,10]{\GraphStruct { V , E }}{%
      \FormulaRefAuto{UndirectedGraphAccess}[thm]{1}}
  \proofstepwidestarbreak[1,10]{\DirectedWalk{V}{E}{q}{n}}{%
    \FormulaRefAuto{DirectedPathWalkAxiom}[ax]{%
      13,\rAEn{10}}}
  \proofstepwidestar[1,9,10,12]{\GraphStruct { V , E }}{%
      \FormulaRefAuto{UndirectedGraphAccess}[thm]{1}}
  \proofstepwidestarbreak[1,9,10,12]{%
    \exists l\in\mathbb N\,\exists t\in\powerset(\mathbb N\times V)\,
    \bigl(\DirectedPath{V}{E}{t}{l}\land t(0)=u\land t(l)=v\bigr)}{%
    \FormulaRefAuto{FinitePathTransportAlongWalk}[thm]{%
      15,
      \rAEa{10},\rAEa{9},14,\rAEb{12},
      \rIE{\rAEa{\rAEn{10}},\rIE{\rAEn{12},\rAEb{\rAEn{9}}}}}}
  \proofstepwidestarbreak[17]{l\in\mathbb N\land
    t\in\powerset(\mathbb N\times V)\land
    \DirectedPath{V}{E}{t}{l}\land t(0)=u\land t(l)=v}{\rA}
  \proofstepwidestar[1,3,4,9,10,12,17]{(l,t)\in\GraphPathsBetween{V}{E}{u}{v}}{%
      \FormulaRefAuto{FiniteSimplePathFamilyMembership}[thm]{%
      \rAEa{17},\rAEb{17},\rAEn{17},\rAEn{17}}}
  \proofstepwidestarbreak[1,3,4,9,10,12,17]{%
    \exists z\,\bigl(z\in\GraphPathsBetween{V}{E}{u}{v}\bigr)}{%
    \rEI{18}}
  \proofstepwidestarbreak[1,3,4,9,10,12]{%
    \exists z\,\bigl(z\in\GraphPathsBetween{V}{E}{u}{v}\bigr)}{%
    \rEE{16,17,19}}
  \proofstepwidestarbreak[1,3,4,9,10]{%
    \exists z\,\bigl(z\in\GraphPathsBetween{V}{E}{u}{v}\bigr)}{%
    \rEE{11,12,20}}
  \proofstepwidestarbreak[1,3,4,9]{%
    \exists z\,\bigl(z\in\GraphPathsBetween{V}{E}{u}{v}\bigr)}{%
    \rEE{8,10,21}}
  \proofstepwidestarbreak[1,3,4,5,6]{%
    \exists z\,\bigl(z\in\GraphPathsBetween{V}{E}{u}{v}\bigr)}{%
    \rEE{7,9,22}}
\end{tabproofwide}

\FormulaThmDeltaK[Eindeutigkeit von Wurzelwegen überträgt sich]{%
  \begin{aligned}[t]
  &\SimpleUndirectedGraphStruct{V,E}\dsep r\in V\dsep
    \forall x\in V\,\exists!z\,
      \bigl(z\in\GraphPathsBetween{V}{E}{r}{x}\bigr)\dsep{}\\[-2pt]
  &u,v\in V\dsep
    a\in\GraphPathsBetween{V}{E}{u}{v}\dsep
    b\in\GraphPathsBetween{V}{E}{u}{v}\vdash a=b
  \end{aligned}%
}{RootPathUniquenessTransfers}{%
  \DeltaRow{Mengen}{V\dsep E\dsep r\dsep x\dsep u\dsep v\dsep a\dsep b\dsep c\dsep d}
}

Der entscheidende Widerspruch ist endlich. Angenommen, \(a\) und \(b\) seien
verschieden. Vom ersten Auseinanderlaufen bis zum ersten erneuten
Zusammentreffen enthalten ihre Knotenfolgen zwei intern disjunkte Wege und
damit einen einfachen Kreis. Unter den Wurzelwegen zu den endlich vielen
Kreisknoten wähle man einen mit kleinster Länge. Er trifft den Kreis vor
seinem Endpunkt nicht, denn sein entsprechender Anfangsabschnitt wäre sonst
ein kürzerer Wurzelweg zu einem Kreisknoten. Vom Endpunkt dieses Wurzelwegs zu
einem anderen Kreisknoten gibt es auf dem Kreis zwei intern disjunkte Bögen.
Das Anfügen beider Bögen erzeugt zwei verschiedene einfache Wurzelwege und
widerspricht der vorausgesetzten Eindeutigkeit.
\begin{tabproofwide}
  \proofstepwidestarbreak[1]{\SimpleUndirectedGraphStruct{V,E}}{\rA}
  \proofstepwidestarbreak[2]{r\in V}{\rA}
  \proofstepwidestarbreak[3]{%
    \forall x\in V\,\exists!z\,
      \bigl(z\in\GraphPathsBetween{V}{E}{r}{x}\bigr)}{\rA}
  \proofstepwidestarbreak[4]{u\in V}{\rA}
  \proofstepwidestarbreak[5]{v\in V}{\rA}
  \proofstepwidestarbreak[6]{a\in\GraphPathsBetween{V}{E}{u}{v}}{\rA}
  \proofstepwidestarbreak[7]{b\in\GraphPathsBetween{V}{E}{u}{v}}{\rA}
  \proofstepwidestarbreak[8]{a\neq b}{\rA}
  \proofstepwidestarbreak[1,2,3,4,5,6,7,8]{%
    \exists x\in V\,\exists c\,\exists d\,\bigl(
      c\in\GraphPathsBetween{V}{E}{r}{x}\land
      d\in\GraphPathsBetween{V}{E}{r}{x}\land c\neq d\bigr)}{%
    \text{\parbox[t]{\linewidth}{\raggedright Erster Kreisabschnitt und kürzester Wurzelweg\\
      zu diesem endlichen Kreis}}}
  \proofstepwidestarbreak[10]{x\in V\land
      c\in\GraphPathsBetween{V}{E}{r}{x}\land
      d\in\GraphPathsBetween{V}{E}{r}{x}\land c\neq d}{\rA}
  \proofstepwidestarbreak[3,10]{%
    \exists!z\,\bigl(z\in\GraphPathsBetween{V}{E}{r}{x}\bigr)}{%
    \rUE{3,\rAEa{10}}}
  \proofstepwidestarbreak[3,10]{c=d}{%
    \FormulaRefAuto{\exists! x\,P(x),\; P(a),\; P(b) \vdash a = b}[thm]{%
      11,\rAEn{10},\rAEn{10}}}
  \proofstepwidestarbreak[10]{c\neq d}{\rAEn{10}}
  \proofstepwidestarbreak[1,2,3,4,5,6,7,8,10]{\bot}{\rBI{12,13}}
  \proofstepwidestarbreak[1,2,3,4,5,6,7]{a=b}{%
    \FormulaRefAuto{rDN}[thm]{%
      \rCI{8,\rEE{9,10,14}}}}
\end{tabproofwide}

\chapter{Verwurzelte Bäume}

\section{Wurzel und Wurzelwege}

\FormulaDefDeltaK[Verwurzelter Baum]{%
  \RootedTreeStruct{V,E,r}%
}{RootedTreeDef}{%
  \DeltaRow{Mengen}{V\dsep E\dsep r}
  \DeltaRow{\textbf{Axiome}}{}
  \DeltaRow{Baumstruktur}{%
    \TreeStruct{V,E}}
  \DeltaRow{Wurzel}{%
    r\in V}
  \DeltaRow{\textbf{Neues Symbol}}{}
  \DeltaPrem{Struktur}{\RootedTreeStruct{V,E,r}}
}

\FormulaAxiomDeltaK[Baumstruktur]{%
  \begin{aligned}[t]
    &\RootedTreeStruct{V,E,r}\vdash{}\\[-2pt]
    &\TreeStruct{V,E}
  \end{aligned}%
}{RootedTreeTreeAxiom}{%
  \DeltaRow{Mengen}{V\dsep E\dsep r}
}

\FormulaAxiomDeltaK[Wurzel]{%
  \begin{aligned}[t]
    &\RootedTreeStruct{V,E,r}\vdash{}\\[-2pt]
    &r\in V
  \end{aligned}%
}{RootedTreeRootAxiom}{%
  \DeltaRow{Mengen}{V\dsep E\dsep r}
}

\FormulaThmDeltaK[Charakterisierung durch die Axiome]{%
  \RootedTreeStruct{V,E,r}\eqvdash\TreeStruct{V,E}\land r\in V%
}{RootedTreeCharacterization}{%
  \DeltaRow{Mengen}{V\dsep E\dsep r}
}
\begin{tabproofsplitwide}
  \proofpartwide{Hinrichtung}
    \proofstepwidestar[1]{\RootedTreeStruct{V,E,r}}{\rA}
    \proofstepwidestar[1]{\TreeStruct{V,E}}{%
      \FormulaRefAuto{RootedTreeTreeAxiom}[ax]{1}}
    \proofstepwidestar[1]{r\in V}{%
      \FormulaRefAuto{RootedTreeRootAxiom}[ax]{1}}
    \proofstepwidestar[1]{%
      \begin{aligned}[t]
      &\TreeStruct{V,E}\land{}\\[-2pt]
      &r\in V
    \end{aligned}}{\rAIStar{2,3}}
  \closeproofpartwide

  \proofpartwide{Rückrichtung}
    \proofstepwidestar[1]{%
      \begin{aligned}[t]
      &\TreeStruct{V,E}\land{}\\[-2pt]
      &r\in V
    \end{aligned}}{\rA}
    \proofstepwidestar[1]{\TreeStruct{V,E}}{\rAEa{1}}
    \proofstepwidestar[1]{r\in V}{\rAEb{1}}
    \proofstepwidestar[1]{\RootedTreeStruct{V,E,r}}{%
      \FormulaRefAuto{RootedTreeDef}[def]{2,3}}
  \closeproofpartwide
\end{tabproofsplitwide}

\FormulaThmDeltaK[Zugriff auf einen verwurzelten Baum]{%
  \RootedTreeStruct{V,E,r}\vdash\TreeStruct{V,E}\land r\in V%
}{RootedTreeAccess}{%
  \DeltaRow{Mengen}{V\dsep E\dsep r}
}
\begin{tabproof}
  \proofstep{1}{\RootedTreeStruct{V,E,r}}{\rA}
  \proofstep{1}{\TreeStruct { V , E }}{%
      \FormulaRefAuto{RootedTreeTreeAxiom}[ax]{1}}
  \proofstep{1}{r \in V}{%
      \FormulaRefAuto{RootedTreeRootAxiom}[ax]{1}}
  \proofstep{1}{\TreeStruct{V,E}\land r\in V}{%
    \rAI{2,
      3}}
\end{tabproof}

\FormulaThmDeltaK[Eindeutiger Wurzelweg]{%
  \RootedTreeStruct{V,E,r}\dsep x\in V\vdash
  \exists!z\,\bigl(z\in\GraphPathsBetween{V}{E}{r}{x}\bigr)%
}{RootPathUnique}{%
  \DeltaRow{Mengen}{V\dsep E\dsep r\dsep x\dsep z}
}
\begin{tabproofwide}
  \proofstepwidestarbreak[1]{\RootedTreeStruct{V,E,r}}{\rA}
  \proofstepwidestarbreak[2]{x\in V}{\rA}
  \proofstepwidestar[1]{\TreeStruct { V , E } \land r \in V}{%
      \FormulaRefAuto{RootedTreeAccess}[thm]{1}}
  \proofstepwidestarbreak[1]{\TreeStruct{V,E}}{%
    \rAEa{3}}
  \proofstepwidestar[1]{\TreeStruct { V , E } \land r \in V}{%
      \FormulaRefAuto{RootedTreeAccess}[thm]{1}}
  \proofstepwidestarbreak[1]{r\in V}{%
    \rAEb{5}}
  \proofstepwidestarbreak[1,2]{%
    \exists!z\,\bigl(z\in\GraphPathsBetween{V}{E}{r}{x}\bigr)}{%
    \FormulaRefAuto{TreeAccess}[thm]{4,6,2}}
\end{tabproofwide}

\FormulaThmDeltaK[Baumkriterium durch eindeutige Wurzelwege]{%
  \begin{aligned}[t]
  &\SimpleUndirectedGraphStruct{V,E}\dsep r\in V\dsep
    \forall x\in V\,\exists!z\,
      \bigl(z\in\GraphPathsBetween{V}{E}{r}{x}\bigr)\vdash{}\\[-2pt]
  &\qquad\RootedTreeStruct{V,E,r}
  \end{aligned}%
}{RootedTreeFromUniqueRootPaths}{%
  \DeltaRow{Mengen}{V\dsep E\dsep r\dsep x\dsep u\dsep v\dsep a\dsep b\dsep z}
}
\begin{tabproofwide}
  \proofstepwidestarbreak[1]{\SimpleUndirectedGraphStruct{V,E}}{\rA}
  \proofstepwidestarbreak[2]{r\in V}{\rA}
  \proofstepwidestarbreak[3]{%
    \forall x\in V\,\exists!z\,
      \bigl(z\in\GraphPathsBetween{V}{E}{r}{x}\bigr)}{\rA}
  \proofstepwidestarbreak[4]{u\in V}{\rA}
  \proofstepwidestarbreak[5]{v\in V}{\rA}
  \proofstepwidestarbreak[2,3,4]{%
    \exists a\,\bigl(a\in\GraphPathsBetween{V}{E}{r}{u}\bigr)}{%
    \FormulaRefAuto{\exists! x\,P(x)\vdash \exists x\,P(x)}[thm]{\rUE{3,4}}}
  \proofstepwidestarbreak[2,3,5]{%
    \exists b\,\bigl(b\in\GraphPathsBetween{V}{E}{r}{v}\bigr)}{%
    \FormulaRefAuto{\exists! x\,P(x)\vdash \exists x\,P(x)}[thm]{\rUE{3,5}}}
  \proofstepwidestarbreak[8]{a\in\GraphPathsBetween{V}{E}{r}{u}}{\rA}
  \proofstepwidestarbreak[9]{b\in\GraphPathsBetween{V}{E}{r}{v}}{\rA}
  \proofstepwidestar[1]{\UndirectedGraphStruct { V , E } \land \GraphStruct { V , E }}{%
      \FormulaRefAuto{SimpleUndirectedGraphAccess}[thm]{1}}
  \proofstepwidestarbreak[1]{\UndirectedGraphStruct{V,E}}{%
    \rAEa{10}}
  \proofstepwidestarbreak[1,2,3,4,5,8,9]{%
    \exists z\,\bigl(z\in\GraphPathsBetween{V}{E}{u}{v}\bigr)}{%
    \FormulaRefAuto{RootPathsGivePairPath}[thm]{11,2,4,5,8,9}}
  \proofstepwidestarbreak[13]{a\in\GraphPathsBetween{V}{E}{u}{v}}{\rA}
  \proofstepwidestarbreak[14]{b\in\GraphPathsBetween{V}{E}{u}{v}}{\rA}
  \proofstepwidestarbreak[1,2,3,4,5,13,14]{a=b}{%
    \FormulaRefAuto{RootPathUniquenessTransfers}[thm]{1,2,3,4,5,13,14}}
  \proofstepwidestarbreak[1,2,3,4,5]{%
    \exists!z\,\bigl(z\in\GraphPathsBetween{V}{E}{u}{v}\bigr)}{%
    \UEI{\rEE{6,8,\rEE{7,9,12}},13,14,15}}
  \proofstepwidestarbreak[1,2,3]{%
    \forall u\in V\,\forall v\in V\,
      \exists!z\,\bigl(z\in\GraphPathsBetween{V}{E}{u}{v}\bigr)}{%
    \rUI{\rRI{4,\rUI{\rRI{5,16}}}}}
  \proofstepwidestarbreak[2]{V\neq\varnothing}{%
    \FormulaRefAuto{a \in A \vdash A \neq \varnothing}[thm]{2}}
  \proofstepwidestarbreak[1,2,3]{\TreeStruct{V,E}}{%
    \FormulaRefAuto{TreeCharacterization}[thm]{1,17,18}}
  \proofstepwidestarbreak[1,2,3]{\RootedTreeStruct{V,E,r}}{%
    \FormulaRefAuto{RootedTreeCharacterization}[thm]{19,2}}
\end{tabproofwide}

\FormulaThmDeltaK[Verwurzelte Bäume und eindeutige Wurzelwege]{%
  \begin{aligned}[t]
  \RootedTreeStruct{V,E,r}\eqvdash{}&
    \SimpleUndirectedGraphStruct{V,E}\land r\in V\land{}\\[-2pt]
  &\forall x\in V\,\exists!z\,
      \bigl(z\in\GraphPathsBetween{V}{E}{r}{x}\bigr)
  \end{aligned}%
}{RootedTreeIffUniqueRootPaths}{%
  \DeltaRow{Mengen}{V\dsep E\dsep r\dsep x\dsep z}
}
\begin{tabproofwide}
  \proofstepwidestarbreak[1]{\RootedTreeStruct{V,E,r}}{\rA}
  \proofstepwidestar[1]{\TreeStruct { V , E } \land r \in V}{%
      \FormulaRefAuto{RootedTreeAccess}[thm]{1}}
  \proofstepwidestarbreak[1]{\TreeStruct{V,E}}{%
    \rAEa{2}}
  \proofstepwidestar[1]{\TreeStruct { V , E } \land r \in V}{%
      \FormulaRefAuto{RootedTreeAccess}[thm]{1}}
  \proofstepwidestarbreak[1]{r\in V}{%
    \rAEb{4}}
  \proofstepwidestarbreak[1]{\SimpleUndirectedGraphStruct{V,E}}{%
    \FormulaRefAuto{TreeAccess}[thm]{3}}
  \proofstepwidestar[1]{x\in V\vdash\exists!z\,(z\in\GraphPathsBetween{V}{E}{r}{x})}{%
      \FormulaRefAuto{RootPathUnique}[thm]{1}}
  \proofstepwidestarbreak[1]{%
    \forall x\in V\,\exists!z\,
      \bigl(z\in\GraphPathsBetween{V}{E}{r}{x}\bigr)}{%
    \rUI{\rRI{7}}}
  \proofstepwidestarbreak[1]{%
    \SimpleUndirectedGraphStruct{V,E}\land r\in V\land
    \forall x\in V\,\exists!z\,
      \bigl(z\in\GraphPathsBetween{V}{E}{r}{x}\bigr)}{%
    \rAIStar{6,5,8}}
  \proofstepwidestarbreak[10]{%
    \SimpleUndirectedGraphStruct{V,E}\land r\in V\land
    \forall x\in V\,\exists!z\,
      \bigl(z\in\GraphPathsBetween{V}{E}{r}{x}\bigr)}{\rA}
  \proofstepwidestarbreak[10]{\RootedTreeStruct{V,E,r}}{%
    \FormulaRefAuto{RootedTreeFromUniqueRootPaths}[thm]{10}}
  \proofstepwidestarbreak[]{%
    \RootedTreeStruct{V,E,r}\leftrightarrow
    \bigl(\SimpleUndirectedGraphStruct{V,E}\land r\in V\land
    \forall x\in V\,\exists!z\,
      (z\in\GraphPathsBetween{V}{E}{r}{x})\bigr)}{%
    \rLRI{\rRI{1,9},\rRI{10,11}}}
\end{tabproofwide}

\section{Eltern und Kinder}

\FormulaDefDeltaK[Elternrelation]{%
  \begin{aligned}
  &\TreeParentRel{V}{E}{r}{a}{x}\\[-2pt]
  &\quad\coloneqq a\in V\land x\in V\setminus\{r\}\land{}\\[-2pt]
  &\qquad\exists k\in\mathbb N\,\exists p\in\powerset(\mathbb N\times V)\,\\[-2pt]
  &\qquad\quad
    \Bigl((k+1,p)\in\GraphPathsBetween{V}{E}{r}{x}
          \land p(k)=a\Bigr).
  \end{aligned}%
}{TreeParentDef}{%
  \DeltaRow{Mengen}{V\dsep E\dsep r\dsep a\dsep x\dsep k\dsep p}
  \DeltaPrem{Verwurzelter Baum}{\RootedTreeStruct{V,E,r}}
  \DeltaRow{Neue Relation}{\TreeParentRel{V}{E}{r}{a}{x}}
}

Die Definition liest \(a\) als den vorletzten Knoten des eindeutigen
Wurzelwegs nach \(x\). Der Ausschluss \(x\neq r\) verhindert einen
künstlichen Elternknoten der Wurzel.

\FormulaThmDeltaKR[Zugriff auf die Elternrelation]{%
  \begin{aligned}[t]
    &\TreeParentRel{V}{E}{r}{a}{x}\vdash{}\\[-2pt]
    &\text{(i)}\quad a\in V,\\[-2pt]
    &\text{(ii)}\quad x\in V\setminus\{r\},\\[-2pt]
    &\text{(iii)}\quad
      \begin{aligned}[t]
        &\exists k\in\mathbb N\,\exists p\in\powerset(\mathbb N\times V)\,\\[-2pt]
        &\quad\Bigl((k+1,p)\in\GraphPathsBetween{V}{E}{r}{x}
          \land p(k)=a\Bigr).
      \end{aligned}
  \end{aligned}%
}{%
  \begin{aligned}[t]
  \TreeParentRel{V}{E}{r}{a}{x}\vdash{}&
    a\in V\dsep x\in V\setminus\{r\}\dsep{}\\[-2pt]
  &\exists k\in\mathbb N\,\exists p\in\powerset(\mathbb N\times V)\,\\[-2pt]
  &\quad
    \Bigl((k+1,p)\in\GraphPathsBetween{V}{E}{r}{x}
          \land p(k)=a\Bigr)
  \end{aligned}%
}{TreeParentAccess}{%
  \DeltaRow{Mengen}{V\dsep E\dsep r\dsep a\dsep x\dsep k\dsep p}
  \DeltaPrem{Verwurzelter Baum}{\RootedTreeStruct{V,E,r}}
}
\begin{tabproofwide}
  \proofstepwidestarbreak[1]{\TreeParentRel{V}{E}{r}{a}{x}}{\rA}
  \proofstepwidestarbreak[1]{a\in V}{\FormulaRefAuto{TreeParentDef}[def]{1}}
  \proofstepwidestarbreak[1]{x\in V\setminus\{r\}}{%
    \FormulaRefAuto{TreeParentDef}[def]{1}}
  \proofstepwidestarbreak[1]{%
    \exists k\in\mathbb N\,\exists p\in\powerset(\mathbb N\times V)\,
    \bigl((k+1,p)\in\GraphPathsBetween{V}{E}{r}{x}
          \land p(k)=a\bigr)}{%
    \FormulaRefAuto{TreeParentDef}[def]{1}}
\end{tabproofwide}

\FormulaThmDeltaK[Ein nichttrivialer Wurzelweg hat positive Länge]{%
  \begin{aligned}[t]
  &\RootedTreeStruct{V,E,r}\dsep x\in V\setminus\{r\}\dsep
    (n,p)\in\GraphPathsBetween{V}{E}{r}{x}\vdash n\neq0
  \end{aligned}%
}{NonrootPathLengthNonzero}{%
  \DeltaRow{Mengen}{V\dsep E\dsep r\dsep x\dsep n\dsep p}
}
\begin{tabproofwide}
  \proofstepwidestarbreak[1]{\RootedTreeStruct{V,E,r}}{\rA}
  \proofstepwidestarbreak[2]{x\in V\setminus\{r\}}{\rA}
  \proofstepwidestarbreak[3]{(n,p)\in\GraphPathsBetween{V}{E}{r}{x}}{\rA}
  \proofstepwidestarbreak[3]{p(0)=r\land p(n)=x}{%
    \FormulaRefAuto{FiniteSimplePathFamilyMembership}[thm]{3}}
  \proofstepwidestarbreak[5]{n=0}{\rA}
  \proofstepwidestarbreak[3,5]{p(n)=p(0)}{\rIE{5}}
  \proofstepwidestarbreak[3,5]{x=r}{\rIE{\rAEb{4},\rIE{6,\rAEa{4}}}}
  \proofstepwidestar[2]{x \notin { \{ r \} }}{%
      \FormulaRefAuto{x\in A\setminus B\vdash x\notin B}[thm]{2}}
  \proofstepwidestarbreak[2]{x\neq r}{%
    \FormulaRefAuto{x \notin \{a\} \eqvdash x \neq a}[thm]{%
      8}}
  \proofstepwidestarbreak[2,3,5]{\bot}{\rBI{7,9}}
  \proofstepwidestarbreak[1,2,3]{n\neq0}{\rCI{5,10}}
\end{tabproofwide}

\FormulaThmDeltaK[Jeder Nichtwurzelknoten besitzt einen Elternknoten]{%
  \RootedTreeStruct{V,E,r}\dsep x\in V\setminus\{r\}\vdash
  \exists a\in V\,\TreeParentRel{V}{E}{r}{a}{x}%
}{TreeParentExists}{%
  \DeltaRow{Mengen}{V\dsep E\dsep r\dsep x\dsep a\dsep n\dsep p\dsep k}
}
\begin{tabproofwide}
  \proofstepwidestarbreak[1]{\RootedTreeStruct{V,E,r}}{\rA}
  \proofstepwidestarbreak[2]{x\in V\setminus\{r\}}{\rA}
  \proofstepwidestarbreak[2]{x\in V}{%
    \FormulaRefAuto{x\in A\setminus B\vdash x\in A}[thm]{2}}
  \proofstepwidestarbreak[1,2]{%
    \exists!z\,\bigl(z\in\GraphPathsBetween{V}{E}{r}{x}\bigr)}{%
    \FormulaRefAuto{RootPathUnique}[thm]{1,3}}
  \proofstepwidestarbreak[1,2]{%
    \exists z\,\bigl(z\in\GraphPathsBetween{V}{E}{r}{x}\bigr)}{%
    \FormulaRefAuto{\exists! x\,P(x)\vdash \exists x\,P(x)}[thm]{4}}
  \proofstepwidestarbreak[6]{z\in\GraphPathsBetween{V}{E}{r}{x}}{\rA}
  \proofstepwidestarbreak[6]{%
    \exists n\in\mathbb N\,\exists p\in\powerset(\mathbb N\times V)\,
    \bigl(z=(n,p)\land\DirectedPath{V}{E}{p}{n}\land
          p(0)=r\land p(n)=x\bigr)}{%
    \FormulaRefAuto{FiniteSimplePathFamilyMembership}[thm]{6}}
  \proofstepwidestarbreak[8]{n\in\mathbb N\land
    p\in\powerset(\mathbb N\times V)\land
    (n,p)\in\GraphPathsBetween{V}{E}{r}{x}}{\rA}
  \proofstepwidestarbreak[1,2,8]{n\neq0}{%
    \FormulaRefAuto{NonrootPathLengthNonzero}[thm]{1,2,\rAEn{8}}}
  \proofstepwidestarbreak[1,2,8]{%
    \exists k\in\mathbb N\,(n=k+1)}{%
    \FormulaRefAuto{PeanoIndexSplitAtOne}[thm]{\rAEa{8},9}}
  \proofstepwidestarbreak[11]{k\in\mathbb N\land n=k+1}{\rA}
  \proofstepwidestar[11]{k + 1 \in \mathbb N}{%
      \FormulaRefAuto{PeanoAddOneClosure}[thm]{\rAEa{11}}}
  \proofstepwidestar[11]{k\leq k+1}{%
      \FormulaRefAuto{PeanoLeqAddOne}[thm]{\rAEa{11}}}
  \proofstepwidestar[11]{k\leq (k+1)\leftrightarrow k\in\NatSeg{(k+1)+1}}{%
      \FormulaRefAuto{PeanoNatSegLeqAddOneMembership}[thm]{%
        \rAEa{11},12}}
  \proofstepwidestarbreak[11]{k\in\NatSeg{(k+1)+1}}{%
    \FormulaRefAuto{P\leftrightarrow Q\dsep P\vdash Q}[thm]{%
      14,
      13}}
  \proofstepwidestar[1,8,11]{\TreeStruct { V , E } \land r \in V}{%
      \FormulaRefAuto{RootedTreeAccess}[thm]{1}}
  \proofstepwidestar[1,8,11]{\SimpleUndirectedGraphStruct{V,E}}{%
      \FormulaRefAuto{TreeAccess}[thm]{%
            \rAEa{16}}}
  \proofstepwidestar[1,8,11]{\UndirectedGraphStruct{V,E}\land\GraphStruct{V,E}}{%
      \FormulaRefAuto{SimpleUndirectedGraphAccess}[thm]{%
          17}}
  \proofstepwidestar[1,8,11]{p\colon\NatSeg{n+1}\to V}{%
      \FormulaRefAuto{FiniteSimplePathSequence}[thm]{%
        \rAEb{18},
        \rAEn{8}}}
  \proofstepwidestarbreak[1,8,11]{p\colon\NatSeg{(k+1)+1}\to V}{%
    \rIE{\rAEb{11},
      19}}
  \proofstepwidestarbreak[1,8,11]{p(k)\in V}{%
    \FormulaRefAuto{F\colon A\to B\dsep x\in A\vdash F(x)\in B}[thm]{20,15}}
  \proofstepwidestarbreak[1,2,8,11]{%
    \TreeParentRel{V}{E}{r}{p(k)}{x}}{%
    \FormulaRefAuto{TreeParentDef}[def]{21,2,\rAEa{11},\rAEn{8},\rII}}
  \proofstepwidestarbreak[1,2,8,11]{%
    \exists a\in V\,\TreeParentRel{V}{E}{r}{a}{x}}{\rEI{21,22}}
  \proofstepwidestarbreak[1,2]{%
    \exists a\in V\,\TreeParentRel{V}{E}{r}{a}{x}}{%
    \rEE{5,6,\rEE{7,8,\rEE{10,11,23}}}}
\end{tabproofwide}

\FormulaThmDeltaK[Höchstens ein Elternknoten]{%
  \begin{aligned}[t]
  &\RootedTreeStruct{V,E,r}\dsep
    \TreeParentRel{V}{E}{r}{a}{x}\dsep
    \TreeParentRel{V}{E}{r}{b}{x}\vdash a=b
  \end{aligned}%
}{TreeParentAtMostOne}{%
  \DeltaRow{Mengen}{V\dsep E\dsep r\dsep a\dsep b\dsep x\dsep k\dsep l\dsep p\dsep q}
}
\begin{tabproofwide}
  \proofstepwidestarbreak[1]{\RootedTreeStruct{V,E,r}}{\rA}
  \proofstepwidestarbreak[2]{\TreeParentRel{V}{E}{r}{a}{x}}{\rA}
  \proofstepwidestarbreak[3]{\TreeParentRel{V}{E}{r}{b}{x}}{\rA}
  \proofstepwidestar[1]{\TreeStruct { V , E } \land r \in V}{%
      \FormulaRefAuto{RootedTreeAccess}[thm]{1}}
  \proofstepwidestarbreak[1]{\TreeStruct{V,E}}{%
    \rAEa{4}}
  \proofstepwidestar[1]{\TreeStruct { V , E } \land r \in V}{%
      \FormulaRefAuto{RootedTreeAccess}[thm]{1}}
  \proofstepwidestarbreak[1]{r\in V}{%
    \rAEb{6}}
  \proofstepwidestar[2]{x\in V\setminus\{r\}}{%
      \FormulaRefAuto{TreeParentAccess}[thm]{2}}
  \proofstepwidestarbreak[2]{x\in V}{%
    \FormulaRefAuto{x\in A\setminus B\vdash x\in A}[thm]{%
      8}}
  \proofstepwidestarbreak[2]{%
    \exists k\in\mathbb N\,\exists p\in\powerset(\mathbb N\times V)\,
    \bigl((k+1,p)\in\GraphPathsBetween{V}{E}{r}{x}\land p(k)=a\bigr)}{%
    \FormulaRefAuto{TreeParentAccess}[thm]{2}}
  \proofstepwidestarbreak[3]{%
    \exists l\in\mathbb N\,\exists q\in\powerset(\mathbb N\times V)\,
    \bigl((l+1,q)\in\GraphPathsBetween{V}{E}{r}{x}\land q(l)=b\bigr)}{%
    \FormulaRefAuto{TreeParentAccess}[thm]{3}}
  \proofstepwidestarbreak[12]{k\in\mathbb N\land p\in\powerset(\mathbb N\times V)\land
    (k+1,p)\in\GraphPathsBetween{V}{E}{r}{x}\land p(k)=a}{\rA}
  \proofstepwidestarbreak[13]{l\in\mathbb N\land q\in\powerset(\mathbb N\times V)\land
    (l+1,q)\in\GraphPathsBetween{V}{E}{r}{x}\land q(l)=b}{\rA}
  \proofstepwidestarbreak[1,2,3,12,13]{(k+1,p)=(l+1,q)}{%
    \FormulaRefAuto{TreePathUnique}[thm]{5,7,9,\rAEn{12},\rAEn{13}}}
  \proofstepwidestarbreak[1,2,3,12,13]{k+1=l+1\land p=q}{%
    \FormulaRefAuto{(a,b)=(c,d)\eqvdash a=c\land b=d}[thm]{14}}
  \proofstepwidestar[1,2,3,12,13]{1 \in \mathbb N}{%
      \FormulaRefAuto{PeanoOneInN}[thm]}
  \proofstepwidestarbreak[1,2,3,12,13]{k=l}{%
    \FormulaRefAuto{PeanoAddRightCancellation}[thm]{%
      \rAEa{12},\rAEa{13},16,\rAEa{15}}}
  \proofstepwidestarbreak[1,2,3,12,13]{p(k)=q(l)}{%
    \rIE{\rAEb{15},\rIE{17}}}
  \proofstepwidestarbreak[1,2,3,12,13]{a=b}{%
    \rIE{\rAEn{12},\rIE{18,\rAEn{13}}}}
  \proofstepwidestarbreak[1,2,3]{a=b}{%
    \rEE{10,12,\rEE{11,13,19}}}
\end{tabproofwide}

\FormulaThmDeltaK[Eindeutiger Elternknoten]{%
  \RootedTreeStruct{V,E,r}\dsep x\in V\setminus\{r\}\vdash
  \exists!a\in V\,\TreeParentRel{V}{E}{r}{a}{x}%
}{TreeParentUnique}{%
  \DeltaRow{Mengen}{V\dsep E\dsep r\dsep x\dsep a\dsep b}
}
\begin{tabproofwide}
  \proofstepwidestarbreak[1]{\RootedTreeStruct{V,E,r}}{\rA}
  \proofstepwidestarbreak[2]{x\in V\setminus\{r\}}{\rA}
  \proofstepwidestarbreak[1,2]{%
    \exists a\in V\,\TreeParentRel{V}{E}{r}{a}{x}}{%
    \FormulaRefAuto{TreeParentExists}[thm]{1,2}}
  \proofstepwidestarbreak[4]{a\in V\land\TreeParentRel{V}{E}{r}{a}{x}}{\rA}
  \proofstepwidestarbreak[5]{b\in V\land\TreeParentRel{V}{E}{r}{b}{x}}{\rA}
  \proofstepwidestarbreak[1,4,5]{a=b}{%
    \FormulaRefAuto{TreeParentAtMostOne}[thm]{1,\rAEb{4},\rAEb{5}}}
  \proofstepwidestarbreak[1,2]{%
    \exists!a\in V\,\TreeParentRel{V}{E}{r}{a}{x}}{%
    \UEI{3,4,5,6}}
\end{tabproofwide}

\FormulaThmDeltaK[Elternknoten sind benachbart]{%
  \RootedTreeStruct{V,E,r}\dsep\TreeParentRel{V}{E}{r}{a}{x}
  \vdash(a,x)\in E%
}{TreeParentAdjacent}{%
  \DeltaRow{Mengen}{V\dsep E\dsep r\dsep a\dsep x\dsep k\dsep p}
}
\begin{tabproofwide}
  \proofstepwidestarbreak[1]{\RootedTreeStruct{V,E,r}}{\rA}
  \proofstepwidestarbreak[2]{\TreeParentRel{V}{E}{r}{a}{x}}{\rA}
  \proofstepwidestar[1]{\TreeStruct { V , E } \land r \in V}{%
      \FormulaRefAuto{RootedTreeAccess}[thm]{1}}
  \proofstepwidestarbreak[1]{\TreeStruct{V,E}}{%
    \rAEa{3}}
  \proofstepwidestarbreak[1]{\SimpleUndirectedGraphStruct{V,E}}{%
    \FormulaRefAuto{TreeAccess}[thm]{4}}
  \proofstepwidestar[1]{\UndirectedGraphStruct { V , E } \land \GraphStruct { V , E }}{%
      \FormulaRefAuto{SimpleUndirectedGraphAccess}[thm]{5}}
  \proofstepwidestarbreak[1]{\GraphStruct{V,E}}{%
    \rAEb{6}}
  \proofstepwidestarbreak[2]{%
    \exists k\in\mathbb N\,\exists p\in\powerset(\mathbb N\times V)\,
    \bigl((k+1,p)\in\GraphPathsBetween{V}{E}{r}{x}\land p(k)=a\bigr)}{%
    \FormulaRefAuto{TreeParentAccess}[thm]{2}}
  \proofstepwidestarbreak[9]{k\in\mathbb N\land
    (k+1,p)\in\GraphPathsBetween{V}{E}{r}{x}\land p(k)=a}{\rA}
  \proofstepwidestarbreak[9]{\DirectedPath{V}{E}{p}{k+1}\land
    p(k+1)=x}{%
    \FormulaRefAuto{FiniteSimplePathFamilyMembership}[thm]{\rAEn{9}}}
  \proofstepwidestarbreak[1,9]{\DirectedWalk{V}{E}{p}{k+1}}{%
    \FormulaRefAuto{DirectedPathWalkAxiom}[ax]{7,\rAEa{10}}}
  \proofstepwidestar[9]{k\leq k}{%
      \FormulaRefAuto{PeanoLeqReflexive}[thm]{\rAEa{9}}}
  \proofstepwidestar[9]{k\leq k\leftrightarrow k\in\NatSeg{k+1}}{%
      \FormulaRefAuto{PeanoNatSegLeqAddOneMembership}[thm]{%
        \rAEa{9},\rAEa{9}}}
  \proofstepwidestarbreak[9]{k\in\NatSeg{k+1}}{%
    \FormulaRefAuto{P\leftrightarrow Q\dsep P\vdash Q}[thm]{%
      13,
      12}}
  \proofstepwidestarbreak[1,9]{(p(k),p(k+1))\in E}{%
    \FormulaRefAuto{DirectedWalkStepEdge}[thm]{7,11,14}}
  \proofstepwidestarbreak[1,9]{(a,x)\in E}{%
    \rIE{\rAEn{9},\rIE{\rAEb{10},15}}}
  \proofstepwidestarbreak[1,2]{(a,x)\in E}{\rEE{8,9,16}}
\end{tabproofwide}

\FormulaThmDeltaK[Die Wurzel besitzt keinen Elternknoten]{%
  \RootedTreeStruct{V,E,r}\vdash
  \neg\exists a\in V\,\TreeParentRel{V}{E}{r}{a}{r}%
}{RootHasNoParent}{%
  \DeltaRow{Mengen}{V\dsep E\dsep r\dsep a}
}
\begin{tabproofwide}
  \proofstepwidestarbreak[1]{\RootedTreeStruct{V,E,r}}{\rA}
  \proofstepwidestarbreak[2]{\exists a\in V\,\TreeParentRel{V}{E}{r}{a}{r}}{\rA}
  \proofstepwidestarbreak[3]{a\in V\land\TreeParentRel{V}{E}{r}{a}{r}}{\rA}
  \proofstepwidestarbreak[3]{r\in V\setminus\{r\}}{%
    \FormulaRefAuto{TreeParentAccess}[thm]{\rAEb{3}}}
  \proofstepwidestarbreak[]{r\in\{r\}}{\FormulaRefAuto{a\in\{a\}}[thm]}
  \proofstepwidestarbreak[3]{r\notin\{r\}}{%
    \FormulaRefAuto{x\in A\setminus B\vdash x\notin B}[thm]{4}}
  \proofstepwidestarbreak[3]{\bot}{\rBI{5,6}}
  \proofstepwidestarbreak[1]{%
    \neg\exists a\in V\,\TreeParentRel{V}{E}{r}{a}{r}}{%
    \rCI{2,\rEE{2,3,7}}}
\end{tabproofwide}

\FormulaDefDeltaK[Kindrelation]{%
  \TreeChildRel{V}{E}{r}{x}{a}\coloneqq
  \TreeParentRel{V}{E}{r}{a}{x}%
}{TreeChildDef}{%
  \DeltaRow{Mengen}{V\dsep E\dsep r\dsep a\dsep x}
  \DeltaPrem{Verwurzelter Baum}{\RootedTreeStruct{V,E,r}}
  \DeltaRow{Neue Relation}{\TreeChildRel{V}{E}{r}{x}{a}}
}

\FormulaDefDeltaK[Kindermenge eines Knotens]{%
  \TreeChildren{V}{E}{r}{a}\coloneqq
  \bigl\{x\in V\bigm|\TreeChildRel{V}{E}{r}{x}{a}\bigr\}%
}{TreeChildrenDef}{%
  \DeltaRow{Mengen}{V\dsep E\dsep r\dsep a\dsep x}
  \DeltaPrem{Verwurzelter Baum}{\RootedTreeStruct{V,E,r}}
  \DeltaRow{Neue Menge}{\TreeChildren{V}{E}{r}{a}}
}

\FormulaThmDeltaK[Elementkriterium der Kindermenge]{%
  \RootedTreeStruct{V,E,r}\dsep a\in V\vdash
  x\in\TreeChildren{V}{E}{r}{a}
  \leftrightarrow\TreeParentRel{V}{E}{r}{a}{x}%
}{TreeChildrenMembership}{%
  \DeltaRow{Mengen}{V\dsep E\dsep r\dsep a\dsep x}
}
\begin{tabproofwide}
  \proofstepwidestarbreak[1]{\RootedTreeStruct{V,E,r}}{\rA}
  \proofstepwidestarbreak[2]{a\in V}{\rA}
  \proofstepwidestarbreak[3]{x\in\TreeChildren{V}{E}{r}{a}}{\rA}
  \proofstepwidestarbreak[3]{x\in V\land\TreeChildRel{V}{E}{r}{x}{a}}{%
    \FormulaRefAuto{TreeChildrenDef}[def]{3}}
  \proofstepwidestarbreak[3]{\TreeParentRel{V}{E}{r}{a}{x}}{%
    \FormulaRefAuto{TreeChildDef}[def]{\rAEb{4}}}
  \proofstepwidestarbreak[6]{\TreeParentRel{V}{E}{r}{a}{x}}{\rA}
  \proofstepwidestar[6]{x\in V\setminus\{r\}}{%
      \FormulaRefAuto{TreeParentAccess}[thm]{6}}
  \proofstepwidestarbreak[6]{x\in V}{%
    \FormulaRefAuto{x\in A\setminus B\vdash x\in A}[thm]{%
      7}}
  \proofstepwidestarbreak[6]{\TreeChildRel{V}{E}{r}{x}{a}}{%
    \FormulaRefAuto{TreeChildDef}[def]{6}}
  \proofstepwidestarbreak[6]{x\in V\land\TreeChildRel{V}{E}{r}{x}{a}}{%
    \rAI{8,9}}
  \proofstepwidestarbreak[6]{x\in\TreeChildren{V}{E}{r}{a}}{%
    \FormulaRefAuto{TreeChildrenDef}[def]{10}}
  \proofstepwidestarbreak[1,2]{%
    x\in\TreeChildren{V}{E}{r}{a}\leftrightarrow
    \TreeParentRel{V}{E}{r}{a}{x}}{%
    \rLRI{\rRI{3,5},\rRI{6,11}}}
\end{tabproofwide}

% Allgemeine endliche Elternsysteme, nach Wurzelwegen sowie Eltern und Kindern.
\section{Endliche Wurzelbäume und gerichtete Elternkanten}

\providecommand{\BXXVIIFiniteRooted}[3]{%
  \mathsf{EndWBaum}(#1,#2,#3)}
\providecommand{\BXXVIIParentSystem}[4]{%
  \mathsf{ElterSys}(#1,#2,#3,#4)}
\providecommand{\BXXVIISymmetricEdges}[2]{%
  \operatorname{sym}_{#1}(#2)}

\noindent\begin{minipage}{\linewidth}
\FormulaDefDeltaK[Endlicher verwurzelter Baum]{%
  \BXXVIIFiniteRooted VEr\coloneqq
  \Finite V\land\RootedTreeStruct{V,E,r}
}{B27FiniteRootedTreeDef}{%
  \DeltaRow{Mengen}{V\dsep E\dsep r}
  \DeltaRow{Neue Struktur}{\BXXVIIFiniteRooted VEr}
}
\end{minipage}\par

Die Endlichkeit ist genau die Endlichkeit des Knotenträgers aus Band~20.
Eltern und Kinder behalten die zuvor eingeführten Definitionen. Eine Beschriftung ist eine zusätzliche Abbildung auf den
betreffenden Knoten; sie verändert weder deren Identität noch die Kanten.

\noindent\begin{minipage}{\linewidth}
\FormulaThmDeltaK[Die Kindermengen eines endlichen Baums sind endlich]{%
  \BXXVIIFiniteRooted VEr\dsep a\in V\vdash
  \Finite{\TreeChildren VEr a}
}{B27FiniteRootedChildrenFinite}{%
  \DeltaRow{Mengen}{V\dsep E\dsep r\dsep a\dsep x}
}
\end{minipage}\par
\begin{tabproofwide}
  \proofstepwidestar[1]{\BXXVIIFiniteRooted VEr}{\rA}
  \proofstepwidestar[2]{a\in V}{\rA}
  \proofstepwidestar[1]{\Finite V\land\RootedTreeStruct{V,E,r}}{%
    \FormulaRefAuto{B27FiniteRootedTreeDef}[def]{1}}
  \proofstepwidestar[4]{x\in\TreeChildren VEr a}{\rA}
  \proofstepwidestar[1,2]{x\in\TreeChildren VEr a\leftrightarrow
    \TreeParentRel VEr a x}{%
    \FormulaRefAuto{TreeChildrenMembership}[thm]{\rAEb{3},2}}
  \proofstepwidestar[1,2,4]{\TreeParentRel VEr a x}{%
    \rRE{4,\rLREa{5}}}
  \proofstepwidestar[1,2,4]{x\in V\setminus\{r\}}{%
    \FormulaRefAuto{TreeParentAccess}[thm]{\rAEb{3},6}}
  \proofstepwidestar[1,2,4]{x\in V}{%
    \FormulaRefAuto{x\in A\setminus B\vdash x\in A}[thm]{7}}
  \proofstepwidestar[1,2]{\TreeChildren VEr a\subseteq V}{%
    \FormulaRefAuto{SubsetIntroductionRule}[thm]{[4]\vdots8}}
  \proofstepwidestar[1,2]{\Finite{\TreeChildren VEr a}}{%
    \FormulaRefAuto{FiniteSubset}[thm]{9,\rAEa{3}}}
\end{tabproofwide}

\subsection{Ein gerichtetes Einführungskriterium}

Eine gerichtete Kante \((a,x)\) soll von einem Elternknoten zu seinem Kind
zeigen. Eine natürliche Höhenmarke verhindert, dass man durch solche
Kanten zu einem früheren Knoten zurückkehrt. Die Marken müssen nicht die
Abstände zur Wurzel sein; es genügt, dass sie an jeder Kante streng wachsen.

\noindent\begin{minipage}{\linewidth}
\FormulaDefDeltaK[Endliches Elternsystem mit Höhenmarken]{%
  \begin{aligned}[t]
    \BXXVIIParentSystem VDrh\coloneqq{}&
      \Finite V\land D\subseteq V\times V\land r\in V\\[-2pt]
    &\land h\colon V\to\mathbb N\land h(r)=0\\[-2pt]
    &\land\bigl(\forall x\in V\setminus\{r\}\,
       \exists!a\in V\,((a,x)\in D)\bigr)\\[-2pt]
    &\land\forall a,x\in V\,
       ((a,x)\in D\rightarrow h(a)<h(x)).
  \end{aligned}
}{B27RankedParentSystemDef}{%
  \DeltaRow{Mengen}{V\dsep D\dsep r\dsep h\dsep a\dsep x}
  \DeltaRow{Neue Struktur}{\BXXVIIParentSystem VDrh}
}
\end{minipage}\par

\noindent\begin{minipage}{\linewidth}
\FormulaDefDeltaK[Vergessen der Kantenrichtung]{%
  \BXXVIISymmetricEdges VD\coloneqq
  \{(a,x)\in V\times V\mid(a,x)\in D\lor(x,a)\in D\}
}{B27SymmetricParentEdgesDef}{%
  \DeltaRow{Mengen}{V\dsep D\dsep a\dsep x}
  \DeltaRow{Neue Menge}{\BXXVIISymmetricEdges VD}
}
\end{minipage}\par

% Zitierbare Zwischenergebnisse des allgemeinen Elternsystem-Kriteriums.
\FormulaThmDeltaK[Graphdaten eines Elternsystems]{%
  \begin{aligned}[t]
    \BXXVIIParentSystem VDrh\vdash{}&\GraphStruct{V,D}\land
      \SimpleUndirectedGraphStruct{V,\BXXVIISymmetricEdges VD}\land{}\\[-2pt]
    &r\in V\land\forall a\in V\;((a,r)\notin D)
  \end{aligned}
}{B27RankedParentGraphData}{\DeltaRow{Mengen}{V\dsep D\dsep r\dsep h\dsep a\dsep x}}
\begin{proof}
\emergencystretch=2em
\FormulaRefAuto{B27RankedParentSystemDef}[def] liefert
\(D\subseteq V\times V\), \(r\in V\), \(h\colon V\to\mathbb N\),
\(h(r)=0\) und die strikte Höhenzunahme an jeder Kante.
Nach \FormulaRefAuto{GraphDef}[def] ist \((V,D)\) ein Graph.
Aus \FormulaRefAuto{B27SymmetricParentEdgesDef}[def] und
\FormulaRefAuto{ProductComprehensionMembership}[thm] folgt für
\(a,x\in V\)
\[
 (a,x)\in\BXXVIISymmetricEdges VD
 \leftrightarrow ((a,x)\in D\lor(x,a)\in D).
\]
Die definierte Relation liegt in \(V\times V\); die rechte Seite ist
unter Vertauschung von \(a,x\) invariant. Daher gilt
\FormulaRefAuto{UndirectedGraphIntroduction}[thm]. Eine Schleife würde
in beiden Disjunktionsfällen \(h(x)<h(x)\) ergeben. Dies widerspricht
\FormulaRefAuto{PeanoStrictEqualityImpossible}[thm] mit der reflexiven
Gleichheit \(h(x)=h(x)\). Mit
\FormulaRefAuto{SimpleUndirectedGraphIntroduction}[thm] folgt die
Schlichtheit. Für \(a\in V\) würde \((a,r)\in D\) entsprechend
\(h(a)<h(r)=0\) liefern. Der Funktionswert liegt in \(\mathbb N\)
nach \FormulaRefAuto{F\colon A\to B\dsep x\in A\vdash F(x)\in B}[thm];
\FormulaRefAuto{PeanoNoLtZero}[thm] schließt die Kante aus.
\end{proof}

\FormulaThmDeltaK[Höhenzunahme entlang eines gerichteten Elternwegs]{%
  \begin{aligned}[t]
    &\BXXVIIParentSystem VDrh\dsep\DirectedPath VDp n\dsep
      i,j\in\NatSeg{n+1}\dsep i<j\\[-2pt]
    &\vdash h(p(i))<h(p(j))
  \end{aligned}
}{B27RankedPathHeight}{\DeltaRow{Mengen}{V\dsep D\dsep r\dsep h\dsep p\dsep n\dsep i\dsep j}}
\begin{proof}
\emergencystretch=2em
Die Typdaten von \(p\) folgen aus
\FormulaRefAuto{B27RankedParentGraphData}[thm] und
\FormulaRefAuto{FinitePathTyping}[thm]. Für jeden Schrittindex
\(k\in\NatSeg n\) liefern
\FormulaRefAuto{DirectedPathWalkAxiom}[ax] und
\FormulaRefAuto{DirectedWalkStepEdge}[thm] die Kante
\((p(k),p(k+1))\in D\). Nach
\FormulaRefAuto{B27RankedParentSystemDef}[def] gilt deshalb
\(h(p(k))<h(p(k+1))\).

Für die festen Parameter \(i,n\) verwenden wir das Prädikat
\[
 Q(d)\coloneqq\bigl(i+(d+1)\in\NatSeg{n+1}\rightarrow
   h(p(i))<h(p(i+(d+1)))\bigr).
\]
Induktion über \(d\) mit
\FormulaRefAuto{PeanoInductionAddOne}[thm] verkettet die Schrittungleichungen.
Bei \(d=0\) ist die Behauptung unter ihrer Indexvoraussetzung die
soeben bewiesene Kantenungleichung. Für \(Q(d+1)\) nehmen wir
\(i+((d+1)+1)\in\NatSeg{n+1}\) an. Dann ist auch der Zwischenindex
\(i+(d+1)\) zulässig; die Anfangsabschnittskriterien
\FormulaRefAuto{PeanoNatSegStrictMembership}[thm] und
\FormulaRefAuto{FinitePathPrefixIndexData}[thm] sichern alle benötigten
Indexzugehörigkeiten. Die Induktionsungleichung und die letzte
Schrittungleichung ergeben \(Q(d+1)\) durch
\FormulaRefAuto{PeanoLtTransitive}[thm]; die Schreibweisen der Indizes
stimmen nach \FormulaRefAuto{PeanoAddAssociative}[thm] überein.
Aus dem ursprünglich gegebenen \(i<j\) liefert
\FormulaRefAuto{PeanoStrictAddOneDistance}[thm] schließlich ein
\(d\in\mathbb N\) mit \(j=i+(d+1)\). Die nun bewiesene Aussage
\(Q(d)\) liefert wegen \(j\in\NatSeg{n+1}\) die Behauptung.
\end{proof}

\FormulaThmDeltaK[Gerichtete Erreichbarkeit vom Wurzelknoten]{%
  \BXXVIIParentSystem VDrh\dsep x\in V\vdash
    \exists z\;(z\in\GraphPathsBetween VDrx)
}{B27RankedDirectedReachability}{\DeltaRow{Mengen}{V\dsep D\dsep r\dsep h\dsep x\dsep n}}
\begin{proof}
\emergencystretch=2em
Wir beweisen für jedes \(n\in\mathbb N\): Jeder Knoten \(x\in V\)
mit \(h(x)\leq n\) ist durch einen gerichteten einfachen Weg erreichbar.
Das Induktionsprinzip ist \FormulaRefAuto{PeanoInductionAddOne}[thm].
Für \(n=0\) gilt \(h(x)=0\), denn
\FormulaRefAuto{PeanoZeroLeq}[thm] und
\FormulaRefAuto{PeanoLeqAntisymmetric}[thm] liefern beide Richtungen.
Wäre \(x\ne r\), so gäbe es nach
\FormulaRefAuto{B27RankedParentSystemDef}[def] einen Elternknoten
\(a\in V\) mit \(h(a)<h(x)=0\), entgegen
\FormulaRefAuto{PeanoNoLtZero}[thm]. Also ist \(x=r\).
Der Graph \(p=\{(0,r)\}\) ist nach
\FormulaRefAuto{SingletonGraphFunction}[thm],
\FormulaRefAuto{SingletonGraphValue}[thm] und
\FormulaRefAuto{PeanoNatSegOneEquality}[thm] eine injektive Folge auf
\(\NatSeg1\) mit Anfang und Ende \(r\). Die Kantenbedingung ist auf
\(\NatSeg0=\varnothing\) leer. Somit gelten
\FormulaRefAuto{DirectedPathCharacterization}[thm] und
\FormulaRefAuto{FiniteSimplePathFamilyMembership}[thm].

Im Induktionsschritt sei \(h(x)\leq n+1\). Falls \(h(x)\leq n\),
gilt die Induktionsannahme. Andernfalls liefert
\FormulaRefAuto{PeanoLeqSplit}[thm], zusammen mit
\FormulaRefAuto{PeanoLeqAddOneStrictBackward}[thm], die Gleichheit \(h(x)=n+1\).
Wegen \(h(r)=0\) und \FormulaRefAuto{PeanoZeroNeAddOne}[thm] ist
\(x\ne r\). Wähle den durch
\FormulaRefAuto{B27RankedParentSystemDef}[def] gegebenen Elternknoten
\(a\in V\). Aus \(h(a)<n+1\) folgt \(h(a)\leq n\) nach
\FormulaRefAuto{PeanoLeqAddOneStrictBackward}[thm]. Die Induktionsannahme und
\FormulaRefAuto{FiniteSimplePathFamilyMembership}[thm] liefern einen
einfachen Weg \(p\) mit \(p(0)=r\) und \(p(m)=a\).
Kein Wert von \(p\) kann \(x\) sein: Am Endindex wäre sonst
\(h(a)=h(x)\); an einem früheren Index ergäbe
\FormulaRefAuto{B27RankedPathHeight}[thm] sogar \(h(x)<h(a)\).
Im ersten Fall widerspricht dies \(h(a)<h(x)\) nach
\FormulaRefAuto{PeanoStrictEqualityImpossible}[thm]. Im zweiten Fall
liefert \FormulaRefAuto{PeanoLtTransitive}[thm] die unmögliche
Ungleichung \(h(x)<h(x)\).
Das Bildmengenkriterium
\FormulaRefAuto{F\colon A\to B\dsep y\in F[A]\vdash\exists x\in A\;y=F(x)}[thm]
liefert folglich \(x\notin p[\NatSeg{m+1}]\).
\FormulaRefAuto{FinitePathFreshEndpointExtension}[thm] hängt die Kante
\((a,x)\) an; \FormulaRefAuto{FiniteSimplePathFamilyMembership}[thm]
liefert das neue gespeicherte Wegobjekt. Damit ist der Induktionsschritt
bewiesen. Für den gegebenen Knoten setzen wir schließlich \(n=h(x)\);
Natürlichkeit und \FormulaRefAuto{PeanoLeqReflexive}[thm] sind erfüllt.
\end{proof}

\FormulaThmDeltaK[Jeder einfache Wurzelweg folgt der Elternrichtung]{%
  \begin{aligned}[t]
    &\BXXVIIParentSystem VDrh\dsep
      \DirectedPath V{\BXXVIISymmetricEdges VD}p n\dsep p(0)=r\vdash
      \DirectedPath VDp n
  \end{aligned}
}{B27RankedRootPathForward}{\DeltaRow{Mengen}{V\dsep D\dsep r\dsep h\dsep p\dsep n\dsep i}}
\begin{proof}
\emergencystretch=2em
Nach \FormulaRefAuto{FinitePathTyping}[thm] ist \(p\) eine injektive
Folge auf \(\NatSeg{n+1}\). Für das feste \(n\) setzen wir
\[
 Q(i)\coloneqq\bigl(i\in\NatSeg n\rightarrow(p(i),p(i+1))\in D\bigr).
\]
Wir beweisen \(Q(i)\) über den natürlichen Schrittindex; das Prinzip ist
\FormulaRefAuto{PeanoInductionAddOne}[thm]. Das Kantenkriterium aus
\FormulaRefAuto{B27SymmetricParentEdgesDef}[def] und
\FormulaRefAuto{ProductComprehensionMembership}[thm] liefert jeweils
die Vorwärts- oder Rückwärtskante.

Ist der jeweils betrachtete Index nicht in \(\NatSeg n\), ist
\(Q(i)\) wegen der falschen Voraussetzung erfüllt. Dies erfasst
insbesondere den Fall \(n=0\). Beim ersten zulässigen Schritt würde
die Rückwärtskante nach \(p(0)=r\)
führen. Dies widerspricht
\FormulaRefAuto{B27RankedParentGraphData}[thm]. Angenommen, der Schritt
von \(p(i)\) nach \(p(i+1)\) liegt in \(D\), und der nächste Schritt
läge rückwärts: \((p(i+2),p(i+1))\in D\).
Der Knoten \(p(i+1)\) ist von \(r=p(0)\) verschieden, da \(p\)
injektiv und \(i+1\ne0\) ist
(\FormulaRefAuto{PeanoZeroNeAddOne}[thm]). Die eindeutige
Elternexistenz in \FormulaRefAuto{B27RankedParentSystemDef}[def]
und \FormulaRefAuto{\exists! x\,P(x),\; P(a),\; P(b) \vdash a = b}[thm]
ergeben \(p(i)=p(i+2)\). Injektivität liefert \(i=i+2\).
Andererseits liefern zwei Anwendungen von
\FormulaRefAuto{PeanoLtAddOne}[thm] die Ungleichungen
\(i<i+1\) und \(i+1<i+2\). Nach
\FormulaRefAuto{PeanoLtTransitive}[thm] gilt also \(i<i+2\), im
Widerspruch zu \FormulaRefAuto{PeanoStrictEqualityImpossible}[thm].
Also ist auch der nächste
Schritt vorwärts gerichtet. Die Indexzugehörigkeiten in diesem Schritt
folgen aus \FormulaRefAuto{FinitePathPrefixIndexData}[thm] und
\FormulaRefAuto{FinitePathIndexShiftOne}[thm].
Die Folge und ihre Injektivität bleiben unverändert; die bewiesenen
Kantenbedingungen liefern mit
\FormulaRefAuto{DirectedPathCharacterization}[thm] den gerichteten Pfad.
\end{proof}

\Needspace{14\baselineskip}
\FormulaThmDeltaK[Eindeutigkeit gerichteter Wurzelwege]{%
  \begin{aligned}[t]
    &\BXXVIIParentSystem VDrh\dsep x\in V\dsep
      z,z'\in\GraphPathsBetween VDrx\vdash z=z'
  \end{aligned}
}{B27RankedDirectedPathUnique}{\DeltaRow{Mengen}{V\dsep D\dsep r\dsep h\dsep x\dsep z\dsep z'}}
\begin{proof}
\emergencystretch=2em
\FormulaRefAuto{FiniteSimplePathFamilyMembership}[thm] schreibt
\(z=(m,p)\), \(z'=(n,q)\) mit natürlichen Längen und einfachen
gerichteten Wegen, beide von \(r\) nach \(x\).
Für die festen Parameter \(V,D,r,h\) betrachten wir
\[
\begin{aligned}
 \mathcal C=\{s\in\mathbb N\mid{}&\exists x\in V\,\exists m,n\in\mathbb N\,
   \exists p,q\in\powerset(\mathbb N\times V)\\[-2pt]
 &((m,p),(n,q)\in\GraphPathsBetween VDrx\land{}\\[-2pt]
 &\quad s=m+n\land(m,p)\ne(n,q))\}.
\end{aligned}
\]
Der Endpunkt wird hier ausdrücklich mitquantifiziert. Die Menge ist
durch Aussonderung definiert. Wäre sie nicht leer, gäbe
\FormulaRefAuto{NaturalWellOrderingPrinciple}[thm] eine kleinste
Gegenbeispiellänge. Wähle zu ihr die angezeigten Endpunkt- und
Wegzeugen. Sie hätten folgende Eigenschaften.

Ist eine Länge null, so ist \(x=r\). Eine positive andere Länge
hätte eine letzte eingehende Kante nach \(r\), was
\FormulaRefAuto{B27RankedParentGraphData}[thm] ausschließt.
Also sind beide Längen null. Die Folgen sind dann auf
\(\NatSeg1=\{0\}\) gleich; dies folgt aus
\FormulaRefAuto{PeanoNatSegOneEquality}[thm] und
\FormulaRefAuto{FunctionExtensionality}[thm].

Sind beide Längen positiv, liefert
\FormulaRefAuto{PeanoNonzeroUniquePredecessor}[thm]
\(m=k+1\) und \(n=l+1\). Der gemeinsame Endpunkt ist wegen
\FormulaRefAuto{B27RankedParentGraphData}[thm] keine Wurzel.
Die beiden letzten Kanten geben daher dieselben Eltern:
\(p(k)=q(l)\), nach \FormulaRefAuto{B27RankedParentSystemDef}[def]
und der Eindeutigkeitsregel
\FormulaRefAuto{\exists! x\,P(x),\; P(a),\; P(b) \vdash a = b}[thm].
\FormulaRefAuto{FinitePathInitialSegment}[thm] liefert die beiden
verkürzten Wege \(p\restriction_{\NatSeg{k+1}}\) und
\(q\restriction_{\NatSeg{l+1}}\) zu diesem gemeinsamen Elternknoten.
Ihre Längensumme ist kleiner: \(k+l<m+n\) folgt aus
\FormulaRefAuto{PeanoLtAddOne}[thm] und
\FormulaRefAuto{PeanoAddStrictMonotoneBoth}[thm]. Die Minimalität
der gewählten Gegenbeispiellänge erzwingt deshalb die Gleichheit
der verkürzten Wegobjekte. Nach
\FormulaRefAuto{(a,b)=(c,d)\eqvdash a=c\land b=d}[thm]
sind \(k=l\) und die beiden Einschränkungen gleich.
Somit sind \(m=n\), alle früheren Koordinaten gleich und beide
letzten Koordinaten gleich \(x\). Die Aufteilung in den früheren
Anfangsabschnitt und den letzten Index folgt aus
\FormulaRefAuto{PeanoNatSegAddOneMembershipIff}[thm].
\FormulaRefAuto{FunctionExtensionality}[thm] ergibt \(p=q\),
und das Paargleichheitskriterium liefert \(z=z'\).
Es gibt daher kein kleinstes und folglich kein Gegenbeispiel.
\end{proof}

\FormulaThmDeltaK[Eindeutige Wege im symmetrisierten Elternsystem]{%
  \BXXVIIParentSystem VDrh\vdash
    \forall x\in V\,\exists!z\;(z\in\GraphPathsBetween V{\BXXVIISymmetricEdges VD}rx)
}{B27RankedSymmetricRootPaths}{\DeltaRow{Mengen}{V\dsep D\dsep r\dsep h\dsep x\dsep z}}
\begin{proof}
\emergencystretch=2em
Für \(x\in V\) liefert
\FormulaRefAuto{B27RankedDirectedReachability}[thm] einen gerichteten
Wurzelweg. Jede Kante aus \(D\) liegt nach
\FormulaRefAuto{B27SymmetricParentEdgesDef}[def] auch in der
symmetrisierten Relation. Folglich bleibt die gleiche injektive Folge
ein Pfad, nach \FormulaRefAuto{DirectedPathCharacterization}[thm].
Mit \FormulaRefAuto{FiniteSimplePathFamilyMembership}[thm] ist also
ein gespeicherter symmetrisierter Wurzelweg vorhanden.

Sind zwei solche Wegobjekte gegeben, so liefern
\FormulaRefAuto{FiniteSimplePathFamilyMembership}[thm] und
\FormulaRefAuto{B27RankedRootPathForward}[thm] gerichtete Wege
mit denselben Folgen und Längen. Ihre Gleichheit folgt aus
\FormulaRefAuto{B27RankedDirectedPathUnique}[thm]. Die Regel
\UEI{} verbindet den vorhandenen Zeugen mit dieser Eindeutigkeit;
\rUI{} schließt über den beliebigen Knoten \(x\).
\end{proof}


\noindent\begin{minipage}{\linewidth}
\FormulaThmDeltaK[Ein Elternsystem trägt einen endlichen Wurzelbaum]{%
  \BXXVIIParentSystem VDrh\vdash
    \BXXVIIFiniteRooted V{\BXXVIISymmetricEdges VD}r
}{B27RankedParentSystemTree}{%
  \DeltaRow{Mengen}{V\dsep D\dsep r\dsep h}
}
\end{minipage}\par

\begin{tabproofwide}
  \proofstepwidestar[1]{\BXXVIIParentSystem VDrh}{\rA}
  \proofstepwidestar[1]{\Finite V}{\FormulaRefAuto{B27RankedParentSystemDef}[def]{1}}
  \proofstepwidestar[1]{\begin{aligned}[t]
    &\GraphStruct{V,D}\land
      \SimpleUndirectedGraphStruct{V,\BXXVIISymmetricEdges VD}\land{}\\[-2pt]
    &r\in V\land\forall a\in V\;((a,r)\notin D)
  \end{aligned}}{\FormulaRefAuto{B27RankedParentGraphData}[thm]{1}}
  \proofstepwidestar[1]{\SimpleUndirectedGraphStruct{V,\BXXVIISymmetricEdges VD}}{\rAEa{\rAEb{3}}}
  \proofstepwidestar[1]{r\in V}{\rAEa{\rAEb{\rAEb{3}}}}
  \proofstepwidestar[1]{\forall x\in V\,\exists!z\;
    (z\in\GraphPathsBetween V{\BXXVIISymmetricEdges VD}rx)}{\FormulaRefAuto{B27RankedSymmetricRootPaths}[thm]{1}}
  \proofstepwidestar[1]{\RootedTreeStruct{V,\BXXVIISymmetricEdges VD,r}}{\FormulaRefAuto{RootedTreeFromUniqueRootPaths}[thm]{4,5,6}}
  \proofstepwidestar[1]{\BXXVIIFiniteRooted V{\BXXVIISymmetricEdges VD}r}{\FormulaRefAuto{B27FiniteRootedTreeDef}[def]{\rAI{2,7}}}
\end{tabproofwide}

\noindent\begin{minipage}{\linewidth}
\FormulaThmDeltaK[Die gerichteten Kanten sind genau die Baumeltern]{%
  \begin{aligned}[t]
    &\BXXVIIParentSystem VDrh\dsep a,x\in V\vdash\\[-2pt]
    &(a,x)\in D\leftrightarrow
      \TreeParentRel V{\BXXVIISymmetricEdges VD}r a x
  \end{aligned}
}{B27RankedParentEdgesAgree}{%
  \DeltaRow{Mengen}{V\dsep D\dsep r\dsep h\dsep a\dsep x}
}
\end{minipage}\par
\noindent\textit{Beweis.}
Nach \FormulaRefAuto{B27RankedParentSystemTree}[thm] ist die
zuvor definierte Elternrelation anwendbar. Gilt \((a,x)\in D\), so verlängert
diese Kante den nach \FormulaRefAuto{B27RankedDirectedReachability}[thm]
vorhandenen gerichteten Wurzelweg nach \(a\) zu einem einfachen
Wurzelweg nach \(x\): Nach \FormulaRefAuto{B27RankedPathHeight}[thm] ist die Höhe von \(x\) größer als die Höhe aller
Knoten des bisherigen Weges; \FormulaRefAuto{FinitePathFreshEndpointExtension}[thm] liefert den verlängerten Weg. Ferner ist \(x\neq r\). Der vorletzte
Knoten dieses Weges ist \(a\); nach
\FormulaRefAuto{TreeParentDef}[def] gilt die behauptete Elternrelation.
Umgekehrt ist ein Baumelter der vorletzte Knoten des eindeutigen
Wurzelwegs. Dieser Weg läuft nach
\FormulaRefAuto{B27RankedRootPathForward}[thm] vollständig in Richtung von \(D\). Seine letzte Kante ist also
\((a,x)\in D\).


\section{Vorfahren}

\FormulaDefDeltaK[Vorfahrrelation]{%
  \begin{aligned}
  &\TreeAncestorRel{V}{E}{r}{a}{x}\\[-2pt]
  &\quad\coloneqq a\in V\land x\in V\land{}\\[-2pt]
  &\qquad\exists n\in\mathbb N\,\exists p\in\powerset(\mathbb N\times V)\,\\[-2pt]
  &\qquad\quad\exists k\in\NatSeg{n+1}\,\Bigl(
      (n,p)\in\GraphPathsBetween{V}{E}{r}{x}\land p(k)=a\Bigr).
  \end{aligned}%
}{TreeAncestorDef}{%
  \DeltaRow{Mengen}{V\dsep E\dsep r\dsep a\dsep x\dsep n\dsep p\dsep k}
  \DeltaPrem{Verwurzelter Baum}{\RootedTreeStruct{V,E,r}}
  \DeltaRow{Neue Relation}{\TreeAncestorRel{V}{E}{r}{a}{x}}
}

Der Knoten selbst ist als Vorfahr zugelassen. Ein \textbf{echter Vorfahr}
von \(x\) ist ein Vorfahr \(a\) mit \(a\neq x\).

\FormulaThmDeltaK[Jeder Knoten ist sein eigener Vorfahr]{%
  \RootedTreeStruct{V,E,r}\dsep x\in V\vdash
  \TreeAncestorRel{V}{E}{r}{x}{x}%
}{TreeAncestorReflexive}{%
  \DeltaRow{Mengen}{V\dsep E\dsep r\dsep x\dsep z\dsep n\dsep p}
}
\begin{tabproofwide}
  \proofstepwidestarbreak[1]{\RootedTreeStruct{V,E,r}}{\rA}
  \proofstepwidestarbreak[2]{x\in V}{\rA}
  \proofstepwidestar[1]{\TreeStruct { V , E } \land r \in V}{%
      \FormulaRefAuto{RootedTreeAccess}[thm]{1}}
  \proofstepwidestarbreak[1]{r\in V}{%
    \rAEb{3}}
  \proofstepwidestar[1,2]{\exists!z\,(z\in\GraphPathsBetween{V}{E}{r}{x})}{%
      \FormulaRefAuto{RootPathUnique}[thm]{1,2}}
  \proofstepwidestarbreak[1,2]{%
    \exists z\,\bigl(z\in\GraphPathsBetween{V}{E}{r}{x}\bigr)}{%
    \FormulaRefAuto{\exists! x\,P(x)\vdash \exists x\,P(x)}[thm]{%
      5}}
  \proofstepwidestarbreak[7]{z\in\GraphPathsBetween{V}{E}{r}{x}}{\rA}
  \proofstepwidestarbreak[7]{%
    \exists n\in\mathbb N\,\exists p\in\powerset(\mathbb N\times V)\,
    \bigl(z=(n,p)\land(n,p)\in\GraphPathsBetween{V}{E}{r}{x}
      \land p(n)=x\bigr)}{%
    \FormulaRefAuto{FiniteSimplePathFamilyMembership}[thm]{7}}
  \proofstepwidestarbreak[9]{n\in\mathbb N\land
    (n,p)\in\GraphPathsBetween{V}{E}{r}{x}\land p(n)=x}{\rA}
  \proofstepwidestar[9]{n\leq n}{%
      \FormulaRefAuto{PeanoLeqReflexive}[thm]{\rAEa{9}}}
  \proofstepwidestar[9]{n\leq n\leftrightarrow n\in\NatSeg{n+1}}{%
      \FormulaRefAuto{PeanoNatSegLeqAddOneMembership}[thm]{%
        \rAEa{9},\rAEa{9}}}
  \proofstepwidestarbreak[9]{n\in\NatSeg{n+1}}{%
    \FormulaRefAuto{P\leftrightarrow Q\dsep P\vdash Q}[thm]{%
      11,
      10}}
  \proofstepwidestarbreak[1,2,9]{\TreeAncestorRel{V}{E}{r}{x}{x}}{%
    \FormulaRefAuto{TreeAncestorDef}[def]{2,2,\rAEa{9},\rAEn{9},12,\rAEn{9}}}
  \proofstepwidestarbreak[1,2]{\TreeAncestorRel{V}{E}{r}{x}{x}}{%
    \rEE{6,7,\rEE{8,9,13}}}
\end{tabproofwide}

\FormulaThmDeltaK[Die Wurzel ist Vorfahr jedes Knotens]{%
  \RootedTreeStruct{V,E,r}\dsep x\in V\vdash
  \TreeAncestorRel{V}{E}{r}{r}{x}%
}{RootAncestorEveryVertex}{%
  \DeltaRow{Mengen}{V\dsep E\dsep r\dsep x\dsep z\dsep n\dsep p}
}
\begin{tabproofwide}
  \proofstepwidestarbreak[1]{\RootedTreeStruct{V,E,r}}{\rA}
  \proofstepwidestarbreak[2]{x\in V}{\rA}
  \proofstepwidestar[1]{\TreeStruct { V , E } \land r \in V}{%
      \FormulaRefAuto{RootedTreeAccess}[thm]{1}}
  \proofstepwidestarbreak[1]{r\in V}{%
    \rAEb{3}}
  \proofstepwidestar[1,2]{\exists!z\,(z\in\GraphPathsBetween{V}{E}{r}{x})}{%
      \FormulaRefAuto{RootPathUnique}[thm]{1,2}}
  \proofstepwidestarbreak[1,2]{%
    \exists z\,\bigl(z\in\GraphPathsBetween{V}{E}{r}{x}\bigr)}{%
    \FormulaRefAuto{\exists! x\,P(x)\vdash \exists x\,P(x)}[thm]{%
      5}}
  \proofstepwidestarbreak[7]{z\in\GraphPathsBetween{V}{E}{r}{x}}{\rA}
  \proofstepwidestarbreak[7]{%
    \exists n\in\mathbb N\,\exists p\in\powerset(\mathbb N\times V)\,
    \bigl(z=(n,p)\land(n,p)\in\GraphPathsBetween{V}{E}{r}{x}
      \land p(0)=r\bigr)}{%
    \FormulaRefAuto{FiniteSimplePathFamilyMembership}[thm]{7}}
  \proofstepwidestarbreak[9]{n\in\mathbb N\land
    (n,p)\in\GraphPathsBetween{V}{E}{r}{x}\land p(0)=r}{\rA}
  \proofstepwidestarbreak[9]{0\leq n}{%
    \FormulaRefAuto{PeanoZeroLeq}[thm]{\rAEa{9}}}
  \proofstepwidestar[9]{0 \in \mathbb N}{%
      \FormulaRefAuto{PeanoZero}[ax]}
  \proofstepwidestar[9]{0\leq n\leftrightarrow 0\in\NatSeg{n+1}}{%
      \FormulaRefAuto{PeanoNatSegLeqAddOneMembership}[thm]{%
        11,\rAEa{9}}}
  \proofstepwidestarbreak[9]{0\in\NatSeg{n+1}}{%
    \FormulaRefAuto{P\leftrightarrow Q\dsep P\vdash Q}[thm]{%
      12,10}}
  \proofstepwidestarbreak[1,2,9]{\TreeAncestorRel{V}{E}{r}{r}{x}}{%
    \FormulaRefAuto{TreeAncestorDef}[def]{4,2,\rAEa{9},\rAEn{9},13,\rAEn{9}}}
  \proofstepwidestarbreak[1,2]{\TreeAncestorRel{V}{E}{r}{r}{x}}{%
    \rEE{6,7,\rEE{8,9,14}}}
\end{tabproofwide}

\chapter{Blätter und binäre Bäume}

\section{Blätter und innere Knoten}

\FormulaDefDeltaK[Blatt und innerer Knoten]{%
  \begin{aligned}
  \TreeLeafPred{V}{E}{r}{x}&\coloneqq
    x\in V\land\TreeChildren{V}{E}{r}{x}=\varnothing,\\
  \TreeInnerPred{V}{E}{r}{x}&\coloneqq
    x\in V\land\TreeChildren{V}{E}{r}{x}\neq\varnothing.
  \end{aligned}%
}{TreeLeafInnerDef}{%
  \DeltaRow{Mengen}{V\dsep E\dsep r\dsep x}
  \DeltaPrem{Verwurzelter Baum}{\RootedTreeStruct{V,E,r}}
  \DeltaRow{Neue Prädikate}{%
    \TreeLeafPred{V}{E}{r}{x}\dsep\TreeInnerPred{V}{E}{r}{x}}
}

\FormulaThmDeltaK[Blätter und innere Knoten zerlegen den Träger]{%
  \begin{aligned}[t]
  &\RootedTreeStruct{V,E,r}\dsep x\in V\vdash{}\\[-2pt]
  &\quad
    \bigl(\TreeLeafPred{V}{E}{r}{x}\lor
          \TreeInnerPred{V}{E}{r}{x}\bigr)\land
    \neg\bigl(\TreeLeafPred{V}{E}{r}{x}\land
              \TreeInnerPred{V}{E}{r}{x}\bigr)
  \end{aligned}%
}{TreeLeafInnerPartition}{%
  \DeltaRow{Mengen}{V\dsep E\dsep r\dsep x}
}
\begin{tabproofsplitwide}
  \proofpartwide{Vollständigkeit der Fallunterscheidung}
  \proofstepwidestarbreak[1]{\RootedTreeStruct{V,E,r}}{\rA}
  \proofstepwidestarbreak[2]{x\in V}{\rA}
  \proofstepwidestarbreak[]{%
    \TreeChildren{V}{E}{r}{x}=\varnothing\lor
    \TreeChildren{V}{E}{r}{x}\neq\varnothing}{%
    \FormulaRefAuto{P\lor\neg P}[thm]}
  \proofstepwidestarbreak[4]{\TreeChildren{V}{E}{r}{x}=\varnothing}{\rA}
  \proofstepwidestarbreak[2,4]{\TreeLeafPred{V}{E}{r}{x}}{%
    \FormulaRefAuto{TreeLeafInnerDef}[def]{2,4}}
  \proofstepwidestarbreak[6]{\TreeChildren{V}{E}{r}{x}\neq\varnothing}{\rA}
  \proofstepwidestarbreak[2,6]{\TreeInnerPred{V}{E}{r}{x}}{%
    \FormulaRefAuto{TreeLeafInnerDef}[def]{2,6}}
  \proofstepwidestarbreak[2]{%
    \TreeLeafPred{V}{E}{r}{x}\lor\TreeInnerPred{V}{E}{r}{x}}{%
    \rOE{3,4,\rOIa{5},6,\rOIb{7}}}
  \closeproofpartwide

  \proofpartwide{Disjunktheit von Blättern und inneren Knoten}
  \setcounter{proofstepnr}{8}
  \proofstepwidestarbreak[9]{%
    \TreeLeafPred{V}{E}{r}{x}\land\TreeInnerPred{V}{E}{r}{x}}{\rA}
  \proofstepwidestarbreak[9]{\TreeChildren{V}{E}{r}{x}=\varnothing}{%
    \FormulaRefAuto{TreeLeafInnerDef}[def]{\rAEa{9}}}
  \proofstepwidestarbreak[9]{\TreeChildren{V}{E}{r}{x}\neq\varnothing}{%
    \FormulaRefAuto{TreeLeafInnerDef}[def]{\rAEb{9}}}
  \proofstepwidestarbreak[9]{\bot}{\rBI{10,11}}
  \proofstepwidestarbreak[]{%
    \neg\bigl(\TreeLeafPred{V}{E}{r}{x}\land
              \TreeInnerPred{V}{E}{r}{x}\bigr)}{\rCI{9,12}}
  \proofstepwidestarbreak[1,2]{%
    \bigl(\TreeLeafPred{V}{E}{r}{x}\lor
          \TreeInnerPred{V}{E}{r}{x}\bigr)\land
    \neg\bigl(\TreeLeafPred{V}{E}{r}{x}\land
              \TreeInnerPred{V}{E}{r}{x}\bigr)}{\rAI{8,13}}
  \closeproofpartwide
\end{tabproofsplitwide}

\FormulaDefDeltaK[Menge der inneren Knoten]{%
  \TreeInnerVertices{V}{E}{r}\coloneqq
  \bigl\{x\in V\bigm|\TreeInnerPred{V}{E}{r}{x}\bigr\}%
}{TreeInnerVertexSetDef}{%
  \DeltaRow{Mengen}{V\dsep E\dsep r\dsep x}
  \DeltaPrem{Verwurzelter Baum}{\RootedTreeStruct{V,E,r}}
  \DeltaRow{Neue Menge}{\TreeInnerVertices{V}{E}{r}}
}

\section{Binäre und volle Binärbäume}

\FormulaDefDeltaK[Binärer Baum]{%
  \BinaryTreeStruct{V,E,r}%
}{BinaryTreeDef}{%
  \DeltaRow{Mengen}{V\dsep E\dsep r\dsep x\dsep a\dsep b\dsep c}
  \DeltaRow{\textbf{Axiome}}{}
  \DeltaRow{Verwurzelter Baum}{%
    \RootedTreeStruct{V,E,r}}
  \DeltaRow{\makecell[l]{Höchstens\\zwei Kinder}}{%
    \begin{aligned}[t]
      &\forall x\in V\,\forall a,b,c\in\TreeChildren{V}{E}{r}{x}\,\\[-2pt]
      &\quad(a=b\lor a=c\lor b=c)
    \end{aligned}}
  \DeltaRow{\textbf{Neues Symbol}}{}
  \DeltaPrem{Struktur}{\BinaryTreeStruct{V,E,r}}
}

\FormulaAxiomDeltaK[Verwurzelter Baum]{%
  \begin{aligned}[t]
    &\BinaryTreeStruct{V,E,r}\vdash{}\\[-2pt]
    &\RootedTreeStruct{V,E,r}
  \end{aligned}%
}{BinaryTreeRootedAxiom}{%
  \DeltaRow{Mengen}{V\dsep E\dsep r\dsep x\dsep a\dsep b\dsep c}
}

\FormulaAxiomDeltaK[Höchstens zwei Kinder]{%
  \begin{aligned}[t]
    &\BinaryTreeStruct{V,E,r}\vdash{}\\[-2pt]
    &\forall x\in V\,\forall a,b,c\in\TreeChildren{V}{E}{r}{x}\,(a=b\lor a=c\lor b=c)
  \end{aligned}%
}{BinaryTreeAtMostTwoChildrenAxiom}{%
  \DeltaRow{Mengen}{V\dsep E\dsep r\dsep x\dsep a\dsep b\dsep c}
}

\FormulaThmDeltaK[Charakterisierung durch die Axiome]{%
  \begin{aligned}
  \BinaryTreeStruct{V,E,r}\eqvdash{}&
    \RootedTreeStruct{V,E,r}\land{}\\[-2pt]
  &\forall x\in V\,\forall a,b,c\in\TreeChildren{V}{E}{r}{x}\,
    (a=b\lor a=c\lor b=c).
  \end{aligned}%
}{BinaryTreeCharacterization}{%
  \DeltaRow{Mengen}{V\dsep E\dsep r\dsep x\dsep a\dsep b\dsep c}
}
\begin{tabproofsplitwide}
  \proofpartwide{Hinrichtung}
    \proofstepwidestar[1]{\BinaryTreeStruct{V,E,r}}{\rA}
    \proofstepwidestar[1]{\RootedTreeStruct{V,E,r}}{%
      \FormulaRefAuto{BinaryTreeRootedAxiom}[ax]{1}}
    \proofstepwidestar[1]{\forall x\in V\,\forall a,b,c\in\TreeChildren{V}{E}{r}{x}\,(a=b\lor a=c\lor b=c)}{%
      \FormulaRefAuto{BinaryTreeAtMostTwoChildrenAxiom}[ax]{1}}
    \proofstepwidestar[1]{%
      \begin{aligned}[t]
      &\RootedTreeStruct{V,E,r}\land{}\\[-2pt]
      &\forall x\in V\,\forall a,b,c\in\TreeChildren{V}{E}{r}{x}\,{}\\[-2pt]
      &\quad(a=b\lor a=c\lor b=c)
    \end{aligned}}{\rAIStar{2,3}}
  \closeproofpartwide

  \proofpartwide{Rückrichtung}
    \proofstepwidestar[1]{%
      \begin{aligned}[t]
      &\RootedTreeStruct{V,E,r}\land{}\\[-2pt]
      &\forall x\in V\,\forall a,b,c\in\TreeChildren{V}{E}{r}{x}\,{}\\[-2pt]
      &\quad(a=b\lor a=c\lor b=c)
    \end{aligned}}{\rA}
    \proofstepwidestar[1]{\RootedTreeStruct{V,E,r}}{\rAEa{1}}
    \proofstepwidestar[1]{\forall x\in V\,\forall a,b,c\in\TreeChildren{V}{E}{r}{x}\,(a=b\lor a=c\lor b=c)}{\rAEb{1}}
    \proofstepwidestar[1]{\BinaryTreeStruct{V,E,r}}{%
      \FormulaRefAuto{BinaryTreeDef}[def]{2,3}}
  \closeproofpartwide
\end{tabproofsplitwide}

Die zweite Zeile besagt ohne Kardinalzahlschreibweise, dass kein Knoten drei
paarweise verschiedene Kinder besitzt.

\FormulaDefDeltaK[Voller Binärbaum]{%
  \FullBinaryTreeStruct{V,E,r}%
}{FullBinaryTreeDef}{%
  \DeltaRow{Mengen}{V\dsep E\dsep r\dsep x\dsep a\dsep b}
  \DeltaRow{\textbf{Axiome}}{}
  \DeltaRow{Binärer Baum}{%
    \BinaryTreeStruct{V,E,r}}
  \DeltaRow{Null oder zwei Kinder}{%
    \begin{aligned}[t]
      &\forall x\in V\,\Bigl(\TreeChildren{V}{E}{r}{x}=\varnothing\lor{}\\[-2pt]
      &\quad\exists a,b\in\TreeChildren{V}{E}{r}{x}\,\bigl(a\neq b\land{}\\[-2pt]
      &\qquad\TreeChildren{V}{E}{r}{x}=\{a,b\}\bigr)\Bigr)
    \end{aligned}}
  \DeltaRow{\textbf{Neues Symbol}}{}
  \DeltaPrem{Struktur}{\FullBinaryTreeStruct{V,E,r}}
}

\FormulaAxiomDeltaK[Binärer Baum]{%
  \begin{aligned}[t]
    &\FullBinaryTreeStruct{V,E,r}\vdash{}\\[-2pt]
    &\BinaryTreeStruct{V,E,r}
  \end{aligned}%
}{FullBinaryTreeBinaryAxiom}{%
  \DeltaRow{Mengen}{V\dsep E\dsep r\dsep x\dsep a\dsep b}
}

\FormulaAxiomDeltaK[Null oder zwei Kinder]{%
  \begin{aligned}[t]
    &\FullBinaryTreeStruct{V,E,r}\vdash{}\\[-2pt]
    &\begin{aligned}[t]
      &\forall x\in V\,\Bigl(\TreeChildren{V}{E}{r}{x}=\varnothing\lor{}\\[-2pt]
      &\quad\exists a,b\in\TreeChildren{V}{E}{r}{x}\,\bigl(a\neq b\land{}\\[-2pt]
      &\qquad\TreeChildren{V}{E}{r}{x}=\{a,b\}\bigr)\Bigr)
    \end{aligned}
  \end{aligned}%
}{FullBinaryTreeChildrenAxiom}{%
  \DeltaRow{Mengen}{V\dsep E\dsep r\dsep x\dsep a\dsep b}
}

\FormulaThmDeltaK[Charakterisierung durch die Axiome]{%
  \begin{aligned}
  &\FullBinaryTreeStruct{V,E,r}\\[-2pt]
  &\quad\eqvdash \BinaryTreeStruct{V,E,r}\land{}\\[-2pt]
  &\qquad\forall x\in V\,\Bigl(\\[-2pt]
  &\qquad\quad\TreeChildren{V}{E}{r}{x}=\varnothing\ \lor{}\\[-2pt]
  &\qquad\quad\exists a,b\in\TreeChildren{V}{E}{r}{x}\,
      \bigl(a\neq b\land
        \TreeChildren{V}{E}{r}{x}=\{a,b\}\bigr)\Bigr).
  \end{aligned}%
}{FullBinaryTreeCharacterization}{%
  \DeltaRow{Mengen}{V\dsep E\dsep r\dsep x\dsep a\dsep b}
}
\begin{tabproofsplitwide}
  \proofpartwide{Hinrichtung}
    \proofstepwidestar[1]{\FullBinaryTreeStruct{V,E,r}}{\rA}
    \proofstepwidestar[1]{\BinaryTreeStruct{V,E,r}}{%
      \FormulaRefAuto{FullBinaryTreeBinaryAxiom}[ax]{1}}
    \proofstepwidestar[1]{\begin{aligned}[t]
      &\forall x\in V\,\Bigl(\TreeChildren{V}{E}{r}{x}=\varnothing\lor{}\\[-2pt]
      &\quad\exists a,b\in\TreeChildren{V}{E}{r}{x}\,\bigl(a\neq b\land{}\\[-2pt]
      &\qquad\TreeChildren{V}{E}{r}{x}=\{a,b\}\bigr)\Bigr)
    \end{aligned}}{%
      \FormulaRefAuto{FullBinaryTreeChildrenAxiom}[ax]{1}}
    \proofstepwidestar[1]{%
      \begin{aligned}[t]
      &\BinaryTreeStruct{V,E,r}\land{}\\[-2pt]
      &\begin{aligned}[t]
      &\forall x\in V\,\Bigl(\TreeChildren{V}{E}{r}{x}=\varnothing\lor{}\\[-2pt]
      &\quad\exists a,b\in\TreeChildren{V}{E}{r}{x}\,\bigl(a\neq b\land{}\\[-2pt]
      &\qquad\TreeChildren{V}{E}{r}{x}=\{a,b\}\bigr)\Bigr)
    \end{aligned}
    \end{aligned}}{\rAIStar{2,3}}
  \closeproofpartwide

  \proofpartwide{Rückrichtung}
    \proofstepwidestar[1]{%
      \begin{aligned}[t]
      &\BinaryTreeStruct{V,E,r}\land{}\\[-2pt]
      &\begin{aligned}[t]
      &\forall x\in V\,\Bigl(\TreeChildren{V}{E}{r}{x}=\varnothing\lor{}\\[-2pt]
      &\quad\exists a,b\in\TreeChildren{V}{E}{r}{x}\,\bigl(a\neq b\land{}\\[-2pt]
      &\qquad\TreeChildren{V}{E}{r}{x}=\{a,b\}\bigr)\Bigr)
    \end{aligned}
    \end{aligned}}{\rA}
    \proofstepwidestar[1]{\BinaryTreeStruct{V,E,r}}{\rAEa{1}}
    \proofstepwidestar[1]{\begin{aligned}[t]
      &\forall x\in V\,\Bigl(\TreeChildren{V}{E}{r}{x}=\varnothing\lor{}\\[-2pt]
      &\quad\exists a,b\in\TreeChildren{V}{E}{r}{x}\,\bigl(a\neq b\land{}\\[-2pt]
      &\qquad\TreeChildren{V}{E}{r}{x}=\{a,b\}\bigr)\Bigr)
    \end{aligned}}{\rAEb{1}}
    \proofstepwidestar[1]{\FullBinaryTreeStruct{V,E,r}}{%
      \FormulaRefAuto{FullBinaryTreeDef}[def]{2,3}}
  \closeproofpartwide
\end{tabproofsplitwide}

\FormulaThmDeltaK[Innere Knoten eines vollen Binärbaums haben zwei Kinder]{%
  \begin{aligned}[t]
  &\FullBinaryTreeStruct{V,E,r}\dsep
    \TreeInnerPred{V}{E}{r}{x}\vdash{}\\[-2pt]
  &\quad\exists a,b\in\TreeChildren{V}{E}{r}{x}\,
    \bigl(a\neq b\land\TreeChildren{V}{E}{r}{x}=\{a,b\}\bigr)
  \end{aligned}%
}{FullBinaryInnerHasTwoChildren}{%
  \DeltaRow{Mengen}{V\dsep E\dsep r\dsep x\dsep a\dsep b}
}
\begin{tabproofwide}
  \proofstepwidestarbreak[1]{\FullBinaryTreeStruct{V,E,r}}{\rA}
  \proofstepwidestarbreak[2]{\TreeInnerPred{V}{E}{r}{x}}{\rA}
  \proofstepwidestarbreak[2]{x\in V}{%
    \FormulaRefAuto{TreeLeafInnerDef}[def]{2}}
  \proofstepwidestarbreak[2]{\TreeChildren{V}{E}{r}{x}\neq\varnothing}{%
    \FormulaRefAuto{TreeLeafInnerDef}[def]{2}}
  \proofstepwidestarbreak[1,3]{%
    \TreeChildren{V}{E}{r}{x}=\varnothing\lor
    \exists a,b\in\TreeChildren{V}{E}{r}{x}\,
      \bigl(a\neq b\land\TreeChildren{V}{E}{r}{x}=\{a,b\}\bigr)}{%
    \FormulaRefAuto{FullBinaryTreeCharacterization}[thm]{1,3}}
  \proofstepwidestarbreak[6]{\TreeChildren{V}{E}{r}{x}=\varnothing}{\rA}
  \proofstepwidestarbreak[2,6]{\bot}{\rBI{6,4}}
  \proofstepwidestarbreak[8]{%
    \exists a,b\in\TreeChildren{V}{E}{r}{x}\,
      \bigl(a\neq b\land\TreeChildren{V}{E}{r}{x}=\{a,b\}\bigr)}{\rA}
  \proofstepwidestarbreak[1,2]{%
    \exists a,b\in\TreeChildren{V}{E}{r}{x}\,
      \bigl(a\neq b\land\TreeChildren{V}{E}{r}{x}=\{a,b\}\bigr)}{%
    \rOE{5,6,\FormulaRefAuto{\bot\vdash P}[thm]{7},8,8}}
\end{tabproofwide}

\section{Geordnete volle Binärbäume}

\FormulaDefDeltaK[Geordneter voller Binärbaum]{%
  \OrderedFullBinaryTreeStruct{V,E,r,\lambda,\rho}%
}{OrderedFullBinaryTreeDef}{%
  \DeltaRow{Mengen}{V\dsep E\dsep r\dsep x}
  \DeltaRow{Funktionen}{\lambda\dsep\rho}
  \DeltaRow{\textbf{Axiome}}{}
  \DeltaRow{Voller Binärbaum}{%
    \FullBinaryTreeStruct{V,E,r}}
  \DeltaRow{Kindoperatoren}{%
    \lambda,\rho\colon\TreeInnerVertices{V}{E}{r}\to V}
  \DeltaRow{Verschiedene Kinder}{%
    \forall x\in\TreeInnerVertices{V}{E}{r}\,(\lambda(x)\neq\rho(x))}
  \DeltaRow{Alle Kinder}{%
    \begin{aligned}[t]
      &\forall x\in\TreeInnerVertices{V}{E}{r}\,\\[-2pt]
      &\quad\TreeChildren{V}{E}{r}{x}=\{\lambda(x),\rho(x)\}
    \end{aligned}}
  \DeltaRow{\textbf{Neues Symbol}}{}
  \DeltaPrem{Struktur}{\OrderedFullBinaryTreeStruct{V,E,r,\lambda,\rho}}
}

\FormulaAxiomDeltaK[Voller Binärbaum]{%
  \begin{aligned}[t]
    &\OrderedFullBinaryTreeStruct{V,E,r,\lambda,\rho}\vdash{}\\[-2pt]
    &\FullBinaryTreeStruct{V,E,r}
  \end{aligned}%
}{OrderedFullBinaryTreeFullAxiom}{%
  \DeltaRow{Mengen}{V\dsep E\dsep r\dsep x}
  \DeltaRow{Funktionen}{\lambda\dsep\rho}
}

\FormulaAxiomDeltaK[Kindoperatoren]{%
  \begin{aligned}[t]
    &\OrderedFullBinaryTreeStruct{V,E,r,\lambda,\rho}\vdash{}\\[-2pt]
    &\lambda,\rho\colon\TreeInnerVertices{V}{E}{r}\to V
  \end{aligned}%
}{OrderedFullBinaryTreeFunctionsAxiom}{%
  \DeltaRow{Mengen}{V\dsep E\dsep r\dsep x}
  \DeltaRow{Funktionen}{\lambda\dsep\rho}
}

\FormulaAxiomDeltaK[Verschiedene Kinder]{%
  \begin{aligned}[t]
    &\OrderedFullBinaryTreeStruct{V,E,r,\lambda,\rho}\vdash{}\\[-2pt]
    &\forall x\in\TreeInnerVertices{V}{E}{r}\,(\lambda(x)\neq\rho(x))
  \end{aligned}%
}{OrderedFullBinaryTreeDistinctAxiom}{%
  \DeltaRow{Mengen}{V\dsep E\dsep r\dsep x}
  \DeltaRow{Funktionen}{\lambda\dsep\rho}
}

\FormulaAxiomDeltaK[Alle Kinder]{%
  \begin{aligned}[t]
    &\OrderedFullBinaryTreeStruct{V,E,r,\lambda,\rho}\vdash{}\\[-2pt]
    &\begin{aligned}[t]
      &\forall x\in\TreeInnerVertices{V}{E}{r}\,\\[-2pt]
      &\quad\TreeChildren{V}{E}{r}{x}=\{\lambda(x),\rho(x)\}
    \end{aligned}
  \end{aligned}%
}{OrderedFullBinaryTreeChildrenAxiom}{%
  \DeltaRow{Mengen}{V\dsep E\dsep r\dsep x}
  \DeltaRow{Funktionen}{\lambda\dsep\rho}
}

\FormulaThmDeltaK[Charakterisierung durch die Axiome]{%
  \begin{aligned}[t]
    &\OrderedFullBinaryTreeStruct{V,E,r,\lambda,\rho}\eqvdash{}\\[-2pt]
    &\quad\FullBinaryTreeStruct{V,E,r}\land{}\\[-2pt]
    &\quad\lambda,\rho\colon\TreeInnerVertices{V}{E}{r}\to V\land{}\\[-2pt]
    &\quad\forall x\in\TreeInnerVertices{V}{E}{r}\,(\lambda(x)\neq\rho(x))\land{}\\[-2pt]
    &\quad\forall x\in\TreeInnerVertices{V}{E}{r}\,
      (\TreeChildren{V}{E}{r}{x}=\{\lambda(x),\rho(x)\})
  \end{aligned}%
}{OrderedFullBinaryTreeCharacterization}{%
  \DeltaRow{Mengen}{V\dsep E\dsep r\dsep x}
  \DeltaRow{Funktionen}{\lambda\dsep\rho}
}
\begin{tabproofsplitwide}
  \proofpartwide{Hinrichtung}
    \proofstepwidestar[1]{\OrderedFullBinaryTreeStruct{V,E,r,\lambda,\rho}}{\rA}
    \proofstepwidestar[1]{\FullBinaryTreeStruct{V,E,r}}{%
      \FormulaRefAuto{OrderedFullBinaryTreeFullAxiom}[ax]{1}}
    \proofstepwidestar[1]{\lambda,\rho\colon\TreeInnerVertices{V}{E}{r}\to V}{%
      \FormulaRefAuto{OrderedFullBinaryTreeFunctionsAxiom}[ax]{1}}
    \proofstepwidestar[1]{\forall x\in\TreeInnerVertices{V}{E}{r}\,(\lambda(x)\neq\rho(x))}{%
      \FormulaRefAuto{OrderedFullBinaryTreeDistinctAxiom}[ax]{1}}
    \proofstepwidestar[1]{\begin{aligned}[t]
      &\forall x\in\TreeInnerVertices{V}{E}{r}\,\\[-2pt]
      &\quad\TreeChildren{V}{E}{r}{x}=\{\lambda(x),\rho(x)\}
    \end{aligned}}{%
      \FormulaRefAuto{OrderedFullBinaryTreeChildrenAxiom}[ax]{1}}
    \proofstepwidestar[1]{%
      \begin{aligned}[t]
      &\FullBinaryTreeStruct{V,E,r}\land{}\\[-2pt]
      &\lambda,\rho\colon\TreeInnerVertices{V}{E}{r}\to V\land{}\\[-2pt]
      &\forall x\in\TreeInnerVertices{V}{E}{r}\,(\lambda(x)\neq\rho(x))\land{}\\[-2pt]
      &\begin{aligned}[t]
      &\forall x\in\TreeInnerVertices{V}{E}{r}\,\\[-2pt]
      &\quad\TreeChildren{V}{E}{r}{x}=\{\lambda(x),\rho(x)\}
    \end{aligned}
    \end{aligned}}{\rAIStar{2,3,4,5}}
  \closeproofpartwide

  \proofpartwide{Rückrichtung}
    \proofstepwidestar[1]{%
      \begin{aligned}[t]
      &\FullBinaryTreeStruct{V,E,r}\land{}\\[-2pt]
      &\lambda,\rho\colon\TreeInnerVertices{V}{E}{r}\to V\land{}\\[-2pt]
      &\forall x\in\TreeInnerVertices{V}{E}{r}\,(\lambda(x)\neq\rho(x))\land{}\\[-2pt]
      &\begin{aligned}[t]
      &\forall x\in\TreeInnerVertices{V}{E}{r}\,\\[-2pt]
      &\quad\TreeChildren{V}{E}{r}{x}=\{\lambda(x),\rho(x)\}
    \end{aligned}
    \end{aligned}}{\rA}
    \proofstepwidestar[1]{\FullBinaryTreeStruct{V,E,r}}{\rAEa{1}}
    \proofstepwidestar[1]{\lambda,\rho\colon\TreeInnerVertices{V}{E}{r}\to V}{\rAEa{\rAEb{1}}}
    \proofstepwidestar[1]{\forall x\in\TreeInnerVertices{V}{E}{r}\,(\lambda(x)\neq\rho(x))}{\rAEa{\rAEb{\rAEb{1}}}}
    \proofstepwidestar[1]{\begin{aligned}[t]
      &\forall x\in\TreeInnerVertices{V}{E}{r}\,\\[-2pt]
      &\quad\TreeChildren{V}{E}{r}{x}=\{\lambda(x),\rho(x)\}
    \end{aligned}}{\rAEb{\rAEb{\rAEb{1}}}}
    \proofstepwidestar[1]{\OrderedFullBinaryTreeStruct{V,E,r,\lambda,\rho}}{%
      \FormulaRefAuto{OrderedFullBinaryTreeDef}[def]{2,3,4,5}}
  \closeproofpartwide
\end{tabproofsplitwide}

Die Funktionen \(\lambda\) und \(\rho\) wählen nicht nachträglich zwei
Kinder aus, sondern gehören zur Struktur. Sie heißen linker und rechter
Kindoperator.

\FormulaThmDeltaKR[Zugriff auf linken und rechten Kindoperator]{%
  \begin{aligned}[t]
    &\OrderedFullBinaryTreeStruct{V,E,r,\lambda,\rho}\dsep
      x\in\TreeInnerVertices{V}{E}{r}\vdash{}\\[-2pt]
    &\text{(i)}\quad\lambda(x),\rho(x)\in\TreeChildren{V}{E}{r}{x},\\[-2pt]
    &\text{(ii)}\quad\lambda(x)\neq\rho(x).
  \end{aligned}%
}{%
  \begin{aligned}[t]
  &\OrderedFullBinaryTreeStruct{V,E,r,\lambda,\rho}\dsep
    x\in\TreeInnerVertices{V}{E}{r}\vdash{}\\[-2pt]
  &\quad\lambda(x),\rho(x)\in\TreeChildren{V}{E}{r}{x}\dsep
    \lambda(x)\neq\rho(x)
  \end{aligned}%
}{OrderedFullBinaryTreeAccess}{%
  \DeltaRow{Mengen}{V\dsep E\dsep r\dsep x}
  \DeltaRow{Funktionen}{\lambda\dsep\rho}
}
\begin{tabproofwide}
  \proofstepwidestarbreak[1]{%
    \OrderedFullBinaryTreeStruct{V,E,r,\lambda,\rho}}{\rA}
  \proofstepwidestarbreak[2]{x\in\TreeInnerVertices{V}{E}{r}}{\rA}
  \proofstepwidestarbreak[1]{%
    \lambda,\rho\colon\TreeInnerVertices{V}{E}{r}\to V}{%
    \FormulaRefAuto{OrderedFullBinaryTreeCharacterization}[thm]{1}}
  \proofstepwidestarbreak[1,2]{\lambda(x)\neq\rho(x)\land
    \TreeChildren{V}{E}{r}{x}=\{\lambda(x),\rho(x)\}}{%
    \FormulaRefAuto{OrderedFullBinaryTreeCharacterization}[thm]{1,2}}
  \proofstepwidestarbreak[]{\lambda(x)\in\{\lambda(x),\rho(x)\}}{%
    \FormulaRefAuto{a\in\{a,b\}}[thm]}
  \proofstepwidestarbreak[]{\rho(x)\in\{\lambda(x),\rho(x)\}}{%
    \FormulaRefAuto{b\in\{a,b\}}[thm]}
  \proofstepwidestarbreak[1,2]{\lambda(x)\in\TreeChildren{V}{E}{r}{x}}{%
    \rIE{\FormulaRefAuto{a=b\vdash b=a}[thm]{\rAEb{4}},5}}
  \proofstepwidestarbreak[1,2]{\rho(x)\in\TreeChildren{V}{E}{r}{x}}{%
    \rIE{\FormulaRefAuto{a=b\vdash b=a}[thm]{\rAEb{4}},6}}
  \proofstepwidestarbreak[1,2]{%
    \lambda(x),\rho(x)\in\TreeChildren{V}{E}{r}{x}\dsep
    \lambda(x)\neq\rho(x)}{%
    \rAIStar{7,8,\rAEa{4}}}
\end{tabproofwide}

\FormulaThmDeltaK[Vergessen der Kinderordnung]{%
  \begin{aligned}[t]
  &\OrderedFullBinaryTreeStruct{V,E,r,\lambda,\rho}\vdash{}\\[-2pt]
  &\quad\FullBinaryTreeStruct{V,E,r}\land{}\\[-2pt]
  &\quad\BinaryTreeStruct{V,E,r}\land\RootedTreeStruct{V,E,r}
  \end{aligned}%
}{OrderedFullBinaryTreeForgetfulChain}{%
  \DeltaRow{Mengen}{V\dsep E\dsep r}
  \DeltaRow{Funktionen}{\lambda\dsep\rho}
}
\begin{tabproofwide}
  \proofstepwidestarbreak[1]{%
    \OrderedFullBinaryTreeStruct{V,E,r,\lambda,\rho}}{\rA}
  \proofstepwidestarbreak[1]{\FullBinaryTreeStruct{V,E,r}}{%
    \FormulaRefAuto{OrderedFullBinaryTreeCharacterization}[thm]{1}}
  \proofstepwidestarbreak[1]{\BinaryTreeStruct{V,E,r}}{%
    \FormulaRefAuto{FullBinaryTreeCharacterization}[thm]{2}}
  \proofstepwidestarbreak[1]{\RootedTreeStruct{V,E,r}}{%
    \FormulaRefAuto{BinaryTreeCharacterization}[thm]{3}}
  \proofstepwidestarbreak[1]{%
    \FullBinaryTreeStruct{V,E,r}\land
    \BinaryTreeStruct{V,E,r}\land
    \RootedTreeStruct{V,E,r}}{\rAIStar{2,3,4}}
\end{tabproofwide}

\chapter{Ausblick}

Die axiomatische Baumstruktur dieses Bandes trennt die Form eines Baumes von
seiner konkreten Codierung. Im folgenden Band wird daher nicht erneut eine
zweite Baumtheorie aufgebaut. Stattdessen werden die dort aus Wörtern
konstruierten Klammerungsbäume mit einem Knoten- und Kantenbestand versehen;
anschließend wird nachgewiesen, dass sie die Axiome eines geordneten vollen
Binärbaums erfüllen. Rekursion und Auswertung der Klammerungsbäume bleiben
dabei die bereits bewiesenen worttheoretischen Konstruktionen.

\end{document}
